\documentclass[11pt]{article}
% Shared preamble for the CSP dichotomy blueprint.
% House style follows the earlier `zhuk_centers` blueprint: numbered theorem
% environments shared across a single counter, a `longtheorem` upright variant
% for statements with many numbered hypotheses, and explicit `convention`
% environments for standing hypotheses.
\usepackage[a4paper,margin=27mm,headheight=14pt]{geometry}
\usepackage[T1]{fontenc}
\usepackage{lmodern}
\usepackage{amsmath,amssymb,amsthm}
\usepackage{longtable,array}
\usepackage{fancyhdr}
\usepackage[hidelinks]{hyperref}

\newtheoremstyle{mainplain}{9pt}{9pt}{\itshape}{}{\bfseries}{.}{0.55em}{}
\newtheoremstyle{mainroman}{9pt}{9pt}{\normalfont}{}{\bfseries}{.}{0.55em}{}
\theoremstyle{mainplain}
\newtheorem{theorem}{Theorem}[section]
\newtheorem{lemma}[theorem]{Lemma}
\newtheorem{proposition}[theorem]{Proposition}
\newtheorem{corollary}[theorem]{Corollary}
\theoremstyle{mainroman}
\newtheorem{longtheorem}[theorem]{Theorem}
\newtheorem{longlemma}[theorem]{Lemma}
\newtheorem{longproposition}[theorem]{Proposition}
\newtheorem{longcorollary}[theorem]{Corollary}
\theoremstyle{definition}
\newtheorem{definition}[theorem]{Definition}
\newtheorem{remark}[theorem]{Remark}
\newtheorem{convention}[theorem]{Convention}
\newtheorem{hazard}[theorem]{Formalization note}

\setlength{\parindent}{1.25em}\setlength{\parskip}{0pt}\setlength{\emergencystretch}{2em}
\newcommand{\alphlabels}{\renewcommand{\labelenumi}{\textup{(\alph{enumi})}}%
  \renewcommand{\theenumi}{\alph{enumi}}}
\newcommand{\romanlabels}{\renewcommand{\labelenumi}{\textup{(\roman{enumi})}}%
  \renewcommand{\theenumi}{\roman{enumi}}}
\providecommand{\tightlist}{\setlength{\itemsep}{0pt}\setlength{\parskip}{0pt}}
\numberwithin{equation}{section}
\setcounter{tocdepth}{2}\setcounter{secnumdepth}{3}
\pagestyle{fancy}\fancyhf{}
\fancyhead[L]{\small CSP dichotomy}\fancyhead[R]{\small\nouppercase{\leftmark}}
\fancyfoot[C]{\thepage}\renewcommand{\headrulewidth}{0.3pt}

% ---------------------------------------------------------------- operators
\DeclareMathOperator{\Sg}{Sg}
\DeclareMathOperator{\Clo}{Clo}
\DeclareMathOperator{\Pol}{Pol}
\DeclareMathOperator{\Inv}{Inv}
\DeclareMathOperator{\Var}{Var}
\DeclareMathOperator{\Con}{Con}
\DeclareMathOperator{\Sol}{Sol}
\DeclareMathOperator{\Expanded}{Exp}
\DeclareMathOperator{\CSP}{CSP}
\DeclareMathOperator{\GenSAT}{SAT}
\DeclareMathOperator*{\proj}{pr}
\newcommand{\ConOne}{\operatorname{Con}_1}

% ---------------------------------------------------------------- algebras
\newcommand{\bA}{\mathbf A}
\newcommand{\bB}{\mathbf B}
\newcommand{\bC}{\mathbf C}
\newcommand{\bD}{\mathbf D}
\newcommand{\bE}{\mathbf E}
\newcommand{\bR}{\mathbf R}
\newcommand{\bS}{\mathbf S}
\newcommand{\bZ}{\mathbf Z}
\newcommand{\Vn}{\mathcal V_{n}}

% ---------------------------------------------------------- subuniverse types
% The six types of Zhuk's theory.  Typed as upright small capitals so that a
% type never reads as a product of variables.
\newcommand{\TBA}{\textup{\textsc{ba}}}
\newcommand{\TC}{\textup{\textsc{c}}}
\newcommand{\TS}{\textup{\textsc{s}}}
\newcommand{\TD}{\textup{\textsc{d}}}
\newcommand{\TL}{\textup{\textsc{l}}}
\newcommand{\TPC}{\textup{\textsc{pc}}}
\newcommand{\MT}[1]{\mathcal M#1}
\newcommand{\TML}{\MT\TL}
\newcommand{\TMPC}{\MT\TPC}
\newcommand{\TMD}{\MT\TD}

% `C <_{T}^{A} B' and its relatives.
\newcommand{\subt}[2][]{<_{#2}^{#1}}
\newcommand{\subte}[2][]{\le_{#2}^{#1}}
\newcommand{\dsubt}[2][]{\dot<_{#2}^{#1}}
\newcommand{\dsubte}[2][]{\dot\le_{#2}^{#1}}
\renewcommand{\lll}{\mathrel{\lhd\!\lhd\!\lhd}}
\newcommand{\dlll}{\mathrel{\dot{\lhd}\!\dot{\lhd}\!\dot{\lhd}}}
\newcommand{\sdle}{\le_{\mathrm{sd}}}

% congruence cover  sigma^*
\newcommand{\cover}[1]{#1^{*}}
% bridge collapse
\newcommand{\bcol}[1]{\widetilde{#1}}
% absorption
\newcommand{\ab}{\lhd}
\newcommand{\abin}{\lhd_{\mathrm{bin}}}
\newcommand{\abc}{\lhd_{\mathrm{c}}}


\hypersetup{pdftitle={The CSP Dichotomy Theorem: a formalization blueprint},
  pdfsubject={Zhuk's simplified proof, rewritten at formalization granularity},
  pdfauthor={}}

\title{The CSP Dichotomy Theorem\\[2pt]
  \large A formalization blueprint}
\author{}\date{}

\begin{document}
\maketitle\vspace{-2.2em}
\begin{abstract}
The Constraint Satisfaction Problem over a finite constraint language is
solvable in polynomial time if the language admits a weak near-unanimity
polymorphism, and is \textup{NP}-complete otherwise. This document rewrites
Zhuk's simplified proof \cite{ZhukSimplified} --- together with the
prerequisites it imports --- as a blueprint for a formalization in Lean~4 and
Mathlib. It states precisely what a formalization can and cannot claim
(\S\ref{sec:target}); fixes the vocabulary at a granularity a proof assistant
can consume (\S\S\ref{sec:conventions}--\ref{sec:instances}); records the
theory's interface (\S\ref{sec:properties}) and its summit
(\S\ref{sec:main-statements}); states the algorithm and the hard half
(\S\S\ref{sec:algorithm}--\ref{sec:hardness}); and reports what Mathlib is
missing, what has been built, and what the remaining work is
(\S\S\ref{sec:complexity-layer}--\ref{sec:formalization}).

Throughout, \emph{formalization notes} mark the places where the informal proof
has latitude that a formal one does not: a suppressed side condition, a
quantifier whose scope is ambiguous, an object asserted to exist without an
argument, or a step whose well-foundedness is left to the reader.
\end{abstract}

\tableofcontents\clearpage

\section*{Introduction}
\addcontentsline{toc}{section}{Introduction}

The dichotomy conjecture of Feder and Vardi \cite{FederVardi} was settled
independently by Bulatov \cite{Bulatov} and Zhuk \cite{ZhukJACM} in 2017. Both
proofs are long. Zhuk's, in the form published in 2020, is seventy-eight
journal pages, and its central induction alternates between global and local
information in a way its author later described as forcing ``a very complicated
induction to prove most of the claims''.

In 2024 Zhuk published a simplified proof \cite{ZhukSimplified}. The
simplification is a redesign of the theory of strong subalgebras so that
\emph{every reduction is either strong or global}: for each domain $D_{x}$ one
can build a chain
$D_{x}\supsetneq D_{x}^{(1)}\supsetneq\dots\supsetneq\{a\}$ in which every step
is either a strong subset or the intersection with a block of an equivalence
relation with very strong properties. The types of the intermediate sets need
not be tracked --- only that such a chain exists, written $C\lll B$. This is
what makes the argument tractable, and it is what makes it a plausible
formalization target.

\medskip

\noindent\textbf{Why this route.} Among the available proofs,
\cite{ZhukSimplified} is the only complete, recent, rigorously written one whose
author has already removed the obstruction that made its predecessor
unformalizable. Bulatov's proof requires tame congruence theory and centralizer
theory, neither of which is formalized anywhere. Restricting to minimal Taylor
algebras --- which \cite{ZhukSimplified} itself suggests would simplify its
\S2 --- turns out to be a bad trade: it would simplify perhaps a quarter of one
section and none of the CSP argument, at the cost of importing the cyclic term
theorem, which the present route does not otherwise need. Brady's notes
\cite{BradyNotes} are the clearest prose in the subject and are the best source
for several of the imported prerequisites, but they stop short of the
dichotomy. The route taken here is \cite{ZhukSimplified} as the spine, with
\cite{BradyNotes} for the imports, on top of an existing formalization of Zhuk's
centre theorem.

\medskip

\noindent\textbf{What is claimed.} Not that the dichotomy has been formalized.
\S\ref{sec:scope} estimates the transitive closure of what must be proved at
$100$--$130$ printed pages, which extrapolates to tens of thousands of lines of
Lean, none of it attempted in any proof assistant. That is a large number but
not, at current throughput, a long time; what governs the schedule is the
critical path and the defects of \S\ref{sec:defects}, not the line count. What is offered is its design: a route, precise statements, a
vocabulary pinned down where the sources leave it loose, an inventory of the
missing Mathlib layers, and a Lean development whose foundational modules are
complete and whose target theorems are stated.

\medskip

\noindent\textbf{What the formalization notes are for.} A blueprint's value is
mostly in its failures of nerve --- the places where writing the statement down
carefully shows that it does not quite say what it appeared to say. The
predecessor project found four such places in twenty-four pages, and all four
were in the foundations rather than in the mathematics. The notes here record
the same kind of thing: that $\mathbf Z_{p}$ lies in $\Vn$ only when
$p\mid n-1$ (Hazard~\ref{haz:Zp}); that $\cover\sigma$ need not be a congruence
(Hazard~\ref{haz:irreducible}); that images do not commute with intersection,
and that one proof needs them to (Hazard~\ref{haz:quotient-rel}); that ``we
cannot weaken forever'' needs a measure that is not the obvious one
(Hazard~\ref{haz:weakening}); that the main induction quantifies over the
instance \emph{inside} (Hazard~\ref{haz:the-induction}). \S\ref{sec:defects} collects the
harder cases separately: ten places where the source cannot be transcribed at
all, of which one --- \cite{ZhukSimplified}'s restatement of a lemma from
\cite{ZhukStrong}, with two hypotheses dropped --- is false as printed, with a
two-line counterexample.

\clearpage
%!TeX root=csp.tex
\section{What is being proved}\label{sec:target}

\subsection{The theorem}

Fix a finite set $A$ and a finite set $\Gamma$ of relations on $A$, the
\emph{constraint language}. An instance of $\CSP(\Gamma)$ is a conjunction
$R_{1}(\bar x_{1})\wedge\dots\wedge R_{s}(\bar x_{s})$ with each
$R_{i}\in\Gamma$; the question is whether it is satisfiable.

\begin{theorem}[Dichotomy; Bulatov, Zhuk]\label{thm:dichotomy}
If some weak near-unanimity operation preserves $\Gamma$, then $\CSP(\Gamma)$ is
solvable in polynomial time. Otherwise $\CSP(\Gamma)$ is \textup{NP}-complete.
\end{theorem}

Both halves are conditional on a model of computation, and this is not a
formality that a formalization may wave through. ``Solvable in polynomial time''
quantifies over an encoding of instances as strings, a machine, and a
polynomial bound on its step count; ``\textup{NP}-complete'' additionally
quantifies over all problems in \textup{NP} and over polynomial-time many-one
reductions. Mathlib supplies Turing machines
(\texttt{Mathlib/Computability/TuringMachine/}) and a notion of computability,
but it has \emph{no} complexity class, no polynomial-time predicate, and no
notion of polynomial-time reduction.

\subsection{Where the work actually is}\label{sec:program-axis}

An earlier draft of this section concluded from that gap that the
complexity-theoretic half should be deferred because building it would dominate
the project. That inference was wrong, and it is worth recording why, because
the corrected version is what organises the rest of the document.

The complexity vocabulary itself --- languages, a cost model, $\mathrm P$,
$\mathrm{NP}$ by verifiers, $\le_{p}$, \textup{NP}-hardness,
\textup{NP}-completeness, and a generic \textup{NP}-hard problem --- has since
been built, in about $850$ lines. That is consistent with the two existing
Cook--Levin formalizations, whose class layers are $702$ lines (Isabelle) and
$992$ lines (Coq) against totals of $56\,514$ and $\approx 16\,500$. Defining
the classes is cheap in every development that has done it.

What is expensive is not complexity theory. It is any statement that requires
one to \emph{exhibit a program}. In a formalization, ``$f$ is computable in
polynomial time'' is not a property one observes of a mathematical function; it
is an existential over programs, discharged only by writing one and proving a
bound on its step count. The dividing line runs along that axis, and it cuts
across the algebra/complexity split rather than agreeing with it:

\begin{center}
\begin{tabular}{@{}lll@{}}
\hline
Target & Exhibit a program? & Cost \\
\hline
\textup{(T0)} the algebraic core & no & the mathematics \\
\textup{(T1)} $\textsc{Solve}$ is correct & no (but see Hazard~\ref{haz:t1-computable}) & moderate \\
\textup{(T2)} $\textsc{Solve}$ is polynomial & \textbf{yes --- the largest one here} & \textbf{large} \\
\textup{(H0)} $\Gamma$ pp-interprets \textsc{Nae} & no & the mathematics \\
\textup{(H1)} pp-interpretation $\Rightarrow\le_{p}$ & yes, a syntactic substitution & small \\
\hline
\end{tabular}
\end{center}

So the document still separates the machine-free statements from the rest, but
for a sharper reason than before, and with a different verdict on \textup{(T2)}.

\subsection{The layered targets}

Write $\mathbf A=(A;w)$ for a finite idempotent algebra whose basic operation
$w$ is a special weak near-unanimity operation, and $\Vn$ for the class of such
algebras with $w$ of arity $n$ (Convention~\ref{conv:standing}).

\paragraph{Tractable half.} Three layers.

\begin{description}\tightlist
\item[\textup{(T0)} \emph{The algebraic core.}] The three theorems of
  \S\ref{sec:main-statements} --- the existence of a strong subuniverse or a
  linear congruence, the safety of a reduction to a strong subuniverse, and the
  codimension-one theorem --- together with everything they rest on. These are
  statements about finite algebras, congruences, and CSP instances regarded as
  finite combinatorial objects. No computation appears in them.
\item[\textup{(T1)} \emph{Correctness of the algorithm.}] Zhuk's algorithm
  written as a total function $\textsc{Solve}$ from instances to Booleans, and
  the proposition
  \[
    \textsc{Solve}(\mathcal I)=\texttt{true}\iff
    \mathcal I\ \text{has a solution}.
  \]
  This is a statement about a recursive function, and it is provable from
  \textup{(T0)} with no cost model. It is the honest formal content of ``the
  algorithm works'' --- \emph{provided} $\textsc{Solve}$ actually computes;
  see Hazard~\ref{haz:t1-computable}.
\item[\textup{(T2)} \emph{The complexity wrapper.}] That $\textsc{Solve}$ runs
  in polynomial time in the size of the instance, for fixed $\Gamma$. This is
  the one place a cost model is needed on this side, and it is a large
  undertaking rather than a wrapper: see Hazard~\ref{haz:t2-is-not-small}.
\end{description}

\paragraph{Hard half.} Two layers.

\begin{description}\tightlist
\item[\textup{(H0)} \emph{The algebraic core.}] If no weak near-unanimity
  operation preserves $\Gamma$ then $\Gamma$ primitively positively interprets
  every finite constraint language --- in particular one whose CSP is
  \textup{NP}-hard. This is again pure algebra: the Galois connection between
  polymorphisms and primitive positive definability, the Maróti--McKenzie
  theorem that a Taylor term yields a weak near-unanimity term, and the
  Bulatov--Jeavons--Krokhin reduction.
\item[\textup{(H1)} \emph{The complexity wrapper.}] That a primitive positive
  interpretation yields a polynomial-time many-one reduction, and hence
  \textup{NP}-hardness. This does need a program, but the program performs a
  syntactic substitution of formulas whose size is linear in the instance, so
  it is genuinely small.
\end{description}

The claim this document makes is that \textup{(T0)}, \textup{(T1)} and
\textup{(H0)} are the mathematics and are what the literature actually proves.
\textup{(H1)} is a small addition on top. \textup{(T2)} is neither small nor
mathematics, but it is not \emph{unrelated} to the dichotomy either --- it is
the assertion that the algorithm one has just proved correct is also fast, and
it is the second-largest component of the project. The blueprint is written for
\textup{(T0)}, \textup{(T1)} and \textup{(H0)};
\S\ref{sec:complexity-layer} states \textup{(T2)} and \textup{(H1)} precisely
and says what building them costs.

\begin{hazard}[\textup{(T1)} is vacuous unless $\textsc{Solve}$ computes]
\label{haz:t1-computable}
``$\textsc{Solve}$ written as a total function'' hides a requirement that no
cost model would catch. If $\textsc{Solve}$ is defined using classical choice
--- for instance by deciding ``does $D_{x}$ have a nonempty proper binary
absorbing subuniverse?'' with \texttt{Classical.dec} --- then
$\textsc{Solve}(\mathcal I)=\texttt{true}\iff\mathcal I$ has a solution is a
true theorem that says nothing algorithmic, because $\textsc{Solve}$ does not
reduce. To mean what it appears to mean, \textup{(T1)} needs every predicate
the algorithm branches on to carry a genuine decision procedure.

That is a real obligation with a locatable cost. Absorption is defined by an
existential over \emph{terms} --- an infinite type --- so
``$C\subt{\TBA}B$'' is not decidable on its face. It becomes decidable through
the bound that a witness of arity at most $|A|^{|A|}$ suffices, which is a
theorem and not a definition. Hazard~\ref{haz:algorithm-formalization}(b)
records that the bound exists; what it does not record, and what matters here,
is that \textup{(T1)} is empty without it. The same applies to the searches
over congruences and over subsets.
\end{hazard}

\begin{hazard}[\textup{(T2)} is not a wrapper]\label{haz:t2-is-not-small}
Two things make \textup{(T2)} much larger than the phrase ``complexity
wrapper'' suggests.

\emph{The program is enormous.} \textup{(T2)} requires $\textsc{Solve}$ ---
\textsc{ForceConsistency}, the search for a strong subuniverse,
\textsc{SolveLinear}, and the mutual recursion between them --- to be expressed
in a cost model, with a proved polynomial bound. Every existing measurement
says this kind of work dominates: it is the $50\,000$ lines of Turing-machine
engineering in the Isabelle Cook--Levin development, against $702$ for the
classes themselves.

\emph{The statement is non-uniform.} Because $\Gamma$ is fixed, the polynomial
may depend on $\Gamma$, and the searches over subsets of $A$, over congruences
on $A$, and over term operations of bounded arity are all constant-time. But
``constant'' here means the search is compiled away into a program whose
\emph{size} depends on $|A|$. So $\textsc{Solve}$ is not one program: it is a
family $\textsc{Solve}_{\Gamma}$ with a polynomial $p_{\Gamma}$ for each
$\Gamma$, and \textup{(T2)} must be stated in that form. Attempting a bound
uniform in $\Gamma$ would be proving something else, and something false.
\end{hazard}

\begin{remark}
The separation is not a dodge. Layer \textup{(T1)} is strictly stronger than
``\,$\CSP(\Gamma)$ is decidable'', which is trivial for a finite template: it
says that one particular, highly non-obvious algorithm --- the one whose
existence is the content of the dichotomy --- is correct. Everything that makes
the dichotomy hard lives in \textup{(T0)}. Conversely, no part of
\textup{(T0)}--\textup{(T1)} becomes easier if a cost model is available.
\end{remark}

\clearpage
%!TeX root=csp.tex
\section{Standing conventions}\label{sec:conventions}

Zhuk's papers carry their hypotheses in prose: ``in this paper we assume that
every algebra is a finite idempotent algebra having a WNU term operation''.
A formalization cannot carry a hypothesis in prose, and an informal reader
cannot audit which results actually consume which hypothesis. We therefore fix
the hypotheses as a labelled convention and record, for each one, where it is
consumed.

\begin{convention}[Standing hypotheses]\label{conv:standing}
Throughout, $n\ge 3$ is a fixed integer and:
\begin{enumerate}\alphlabels\tightlist
\item\label{conv:finite} every algebra is \emph{finite} and has a
  \emph{nonempty} universe;
\item\label{conv:onebasic} every algebra has exactly one basic operation, of
  arity $n$, written $w$;
\item\label{conv:idem} $w$ is \emph{idempotent}: $w(x,\dots,x)=x$;
\item\label{conv:wnu} $w$ is a \emph{weak near-unanimity} operation:
  $w(y,x,\dots,x)=w(x,y,x,\dots,x)=\dots=w(x,\dots,x,y)$ for all $x,y$;
\item\label{conv:special} $w$ is \emph{special}:
  $w(x,\dots,x,y)=w\bigl(x,\dots,x,w(x,\dots,x,y)\bigr)$ for all $x,y$.
\end{enumerate}
The class of algebras satisfying \ref{conv:finite}--\ref{conv:special} is
written $\Vn$. Since only finite algebras are considered, $\Vn$ is not a
variety.
\end{convention}

\begin{remark}[Why one basic operation]\label{rem:one-basic}
Restricting to a single basic operation is not a loss. The dichotomy hypothesis
is that \emph{some} weak near-unanimity operation preserves $\Gamma$; one may
then discard the original signature and work in the algebra
$(A;w)$ whose only basic operation is that one, because every relation of
$\Gamma$ is a subuniverse of a power of $(A;w)$ and that is all the argument
uses. Passing from a WNU $w$ to a \emph{special} one is
Lemma~\ref{lem:special-wnu}, and costs an increase of arity from $n$ to
$n^{n!}$.

The gain is large for a formalization: a term is a finite tree with $n$-ary
branching over a single symbol, so ``term operation'' needs no signature
bookkeeping, and $\Clo(\mathbf A)$ is the clone generated by one operation.
\end{remark}

\begin{hazard}[Where each hypothesis is consumed]\label{haz:hyp-audit}
\ref{conv:finite} is used in every induction on $|A|$ and in every
minimality argument (a minimal element of a nonempty family of subsets
exists); it cannot be dropped anywhere. \ref{conv:idem} is what makes
singletons subuniverses and makes $\Sg$ of a nonempty set nonempty.
\ref{conv:wnu} is consumed only through the results it implies --- the
existence of absorption-theoretic structure --- and never directly in
\S\S\ref{sec:instances}--\ref{sec:main-statements}. \ref{conv:special} is used
where a unary polynomial must be idempotent under iteration, chiefly in the
theory of $\Sg$-closures of two-element sets. A formalization should carry
\ref{conv:finite}--\ref{conv:special} as typeclass instances on a structure and
not as ambient hypotheses, so that the audit is mechanical.
\end{hazard}

\begin{convention}[The empty set]\label{conv:empty}
A \emph{subuniverse} of $\mathbf A$ is a subset closed under $w$; the empty set
is a subuniverse under this definition, and we allow it. A
\emph{subalgebra} is a nonempty subuniverse; $\mathbf B\le\mathbf A$ always
carries $B\neq\varnothing$.

The relations $\subt{T}$, $\subte{T}$ and $\lll$ of \S\ref{sec:types} are
defined only between \emph{nonempty} sets; in particular
$\varnothing\lll A$ never holds. Where a construction can produce the empty
set --- a projection of an intersection, most often --- the dotted variants
$\dsubt{T}$, $\dsubte{T}$, $\dlll$ are used, and they mean ``the stated relation
holds, or the left-hand side is empty''.
\end{convention}

\begin{hazard}[The dot is not decoration]\label{haz:dot}
Conflating $\lll$ with $\dlll$ is the single most common way to misread
\S\ref{sec:types}. Every propagation statement of \S\ref{sec:properties} that
projects or intersects produces a dotted conclusion, and every statement that
consumes a $\lll$ hypothesis needs the undotted one. In a formalization the two
should be distinct constants, with the dotted one \emph{defined} as
\texttt{C = }$\varnothing$\texttt{ }$\vee$\texttt{ C }$\lll$\texttt{ B}, so that the case split is forced
at every use site rather than being resolved by the reader's judgement.
\end{hazard}

\begin{convention}[Relations, and indices]\label{conv:relations}
A relation is a subset of a product $A_{1}\times\dots\times A_{m}$ of universes,
indexed by $[m]=\{1,\dots,m\}$; $\mathbf R\le\mathbf A_{1}\times\dots\times
\mathbf A_{m}$ means additionally that $R$ is a subuniverse of the product
algebra. We write $\proj_{i_{1},\dots,i_{s}}(R)$ for the projection onto the
listed coordinates and $\mathbf R\sdle\mathbf A_{1}\times\dots\times\mathbf
A_{m}$ when every $\proj_{i}(R)=A_{i}$.

Products are indexed by \emph{arbitrary finite index sets}, not only by initial
segments of $\mathbb N$. Several constructions in \S\ref{sec:main-statements}
form a product over the variable set of an instance, or over a set of
subuniverses, and reindexing those to $[m]$ by hand is exactly the kind of
bookkeeping that a formalization pays for repeatedly and that buys nothing.
\end{convention}

%!TeX root=csp.tex
\section{Vocabulary}\label{sec:vocabulary}

This section fixes the universal-algebraic vocabulary. It is Zhuk's
\S2.1, rewritten so that every definition is unambiguous when read by a
machine. The substantive content begins in \S\ref{sec:types}; a reader who
knows the subject should skim to
Definition~\ref{def:irreducible} (irreducible congruences) and
Definition~\ref{def:bridge} (bridges), which are the two places where the
conventions genuinely matter.

\subsection{Algebras, subuniverses, terms}

\begin{definition}[Algebra]\label{def:algebra}
An algebra is a pair $\mathbf A=(A;w^{\mathbf A})$ with $A$ a finite nonempty
set and $w^{\mathbf A}\colon A^{n}\to A$, subject to
Convention~\ref{conv:standing}. We write $A$ for the universe of $\mathbf A$
and drop the superscript on $w$ when no confusion arises.
\end{definition}

\begin{definition}[The algebras $\mathbf Z_{p}$]\label{def:Zp}
For a prime $p$ with $p\mid n-1$, $\mathbf Z_{p}$ is the algebra on
$\{0,1,\dots,p-1\}$ with $w^{\mathbf Z_{p}}(x_{1},\dots,x_{n})
=x_{1}+\dots+x_{n}\bmod p$.
\end{definition}

\begin{hazard}[$\mathbf Z_{p}$ is only sometimes in $\Vn$]\label{haz:Zp}
Zhuk defines $\mathbf Z_{p}$ for every prime $p$ and then asserts that
``every algebra $\mathbf Z_{p}$ belongs to $\Vn$''. That assertion carries a
silent side condition. The operation $x_{1}+\dots+x_{n}\bmod p$ is idempotent
exactly when $n\equiv 1\pmod p$: idempotence says $nx\equiv x$ for all $x$,
i.e.\ $p\mid n-1$. Given that, it is also a WNU (all the identities read
$(n-1)x+y=y$) and special (both sides reduce to $y$).

So $\mathbf Z_{p}\in\Vn$ if and only if $p\mid n-1$, and every statement of the
form ``$\mathbf A/\sigma\cong\mathbf Z_{p}$ for some prime $p$''
(Lemma~\ref{lem:linear-on-top}, Definition~\ref{def:perfect-linear}) silently
produces a $p$ dividing $n-1$. A formalization must decide whether
$\mathbf Z_{p}$ is defined for all $p$ and merely fails to lie in $\Vn$ for bad
$p$, or is defined only for $p\mid n-1$. The first is cleaner; the second makes
the membership free. Either way the side condition must be visible, because
``for some prime $p$'' otherwise reads as an unconstrained existential.
\end{hazard}

\begin{definition}[Subuniverse, subalgebra]\label{def:subuniverse}
$B\subseteq A$ is a \emph{subuniverse} of $\mathbf A$ if
$w(b_{1},\dots,b_{n})\in B$ for all $b_{1},\dots,b_{n}\in B$. If moreover
$B\neq\varnothing$ then $\mathbf B=(B;w\!\restriction_{B^{n}})$ is a
\emph{subalgebra}, written $\mathbf B\le\mathbf A$.
\end{definition}

\begin{definition}[Term operations]\label{def:clone}
$\Clo(\mathbf A)$ is the least set of finitary operations on $A$ containing all
projections and closed under composition with $w$. Its $m$-ary part is written
$\Clo_{m}(\mathbf A)$. Elements of $\Clo(\mathbf A)$ are the \emph{term
operations} of $\mathbf A$.
\end{definition}

\begin{hazard}[Terms versus term operations]\label{haz:term-vs-op}
Zhuk writes $t\in\Clo(\mathbf A)$ and treats $t$ as both a syntactic term and
the operation it induces. These must be distinguished in a formalization: a
term is a finite tree, and distinct terms may induce the same operation. Every
statement below that quantifies over $\Clo(\mathbf A)$ is a statement about
\emph{operations}, so the syntactic layer may be quotiented away --- but it
cannot be skipped, because $\Sg$ is characterised through terms
(Lemma~\ref{lem:sg-terms}) and because the arity bookkeeping in the absorption
arguments is syntactic.
\end{hazard}

\begin{definition}[Generation]\label{def:sg}
For $X\subseteq A^{m}$, $\Sg_{\mathbf A}(X)$ is the least subuniverse of
$\mathbf A^{m}$ containing $X$.
\end{definition}

\begin{lemma}[Generation by terms]\label{lem:sg-terms}
Let $X=\{x_{1},\dots,x_{k}\}\subseteq A^{m}$ be nonempty. Then
\[
  \Sg_{\mathbf A}(X)=\bigl\{\,t^{\mathbf A^{m}}(x_{1},\dots,x_{k})
      \;:\;t\in\Clo_{k}(\mathbf A)\,\bigr\}.
\]
The generator list is \emph{fixed}: the same list $x_{1},\dots,x_{k}$ in the
same order serves every element of the closure, with $t$ varying.
\end{lemma}

\begin{hazard}\label{haz:sg-fixed-list}
The ``fixed list'' phrasing is what makes Lemma~\ref{lem:sg-terms} usable. The
weaker reading --- for each $y$ in the closure there are \emph{some}
generators and \emph{some} term --- is also true but useless, because every
later application applies one term simultaneously to several closure elements
and needs their generator lists to coincide. This was a defect found in review
of the predecessor blueprint and is called out here for the same reason.
\end{hazard}

\subsection{Relations}

\begin{definition}[Relational vocabulary]\label{def:rel-vocab}
Let $R\subseteq A_{1}\times\dots\times A_{m}$.
\begin{enumerate}\alphlabels\tightlist
\item $R$ is \emph{subdirect} if $\proj_{i}(R)=A_{i}$ for every $i$;
  $\mathbf R\sdle\mathbf A_{1}\times\dots\times\mathbf A_{m}$ abbreviates
  ``$R$ is a subdirect subuniverse''.
\item For $R\subseteq A^{m}$: $R$ is \emph{reflexive} if
  $(a,\dots,a)\in R$ for every $a\in A$.
\item For binary $\delta_{1}\subseteq A_{1}\times A_{2}$,
  $\delta_{2}\subseteq A_{2}\times A_{3}$:
  $\delta_{1}\circ\delta_{2}=\{(a,c):\exists b\;(a,b)\in\delta_{1}\wedge
  (b,c)\in\delta_{2}\}$, and $\delta^{-1}=\{(b,a):(a,b)\in\delta\}$.
\item For $B\subseteq A_{1}$ and $\delta\subseteq A_{1}\times A_{2}$:
  $B\circ\delta=\{c:\exists b\in B\;(b,c)\in\delta\}$.
\item A subdirect $\delta\subseteq A\times B$ is \emph{linked} if the bipartite
  graph on $A\sqcup B$ with edge set $\delta$ is connected. It is
  \emph{bijective} if $|\delta|=|A|=|B|$.
\item For binary $\sigma$ and $m\ge 2$,
  $\sigma^{[m]}=\{(a_{1},\dots,a_{m}):(a_{i},a_{j})\in\sigma\ \forall i,j\}$.
\end{enumerate}
\end{definition}

\begin{hazard}[Composition is diagrammatic]\label{haz:comp-order}
$\delta_{1}\circ\delta_{2}$ applies $\delta_{1}$ \emph{first}. This is the
order of \texttt{Rel.comp} and the opposite of \texttt{Function.comp}. Every
composite of bridges and of congruences below is written left to right, and
silently reversing it produces statements that are true-looking and wrong ---
in Corollary~\ref{cor:stable-intersection}(l) the composite
$\sigma_{k}\circ\proj_{k,\ell}(R)\circ\sigma_{\ell}$ is asymmetric in $k$ and
$\ell$ and the order is the content.
\end{hazard}

\begin{hazard}[``Linked'' is used for two different things]\label{haz:linked}
Zhuk uses \emph{linked} both for a binary relation (Definition~\ref{def:rel-vocab}(e))
and for a CSP instance (Definition~\ref{def:instance-linked}). They are related
but not the same predicate, and the coincidence of names causes real confusion
in \S\ref{sec:main-statements}, where both occur in one proof. Use two names.
\end{hazard}

\begin{definition}[Quotients of relations]\label{def:quotient-rel}
For an equivalence $\sigma$ on $A$ and $a\in A$, $a/\sigma$ is the class of
$a$; for $B\subseteq A$, $B/\sigma=\{b/\sigma:b\in B\}$; for
$R\subseteq A^{m}$, $R/\sigma=\{(b_{1}/\sigma,\dots,b_{m}/\sigma):\bar b\in R\}$.
\end{definition}

\begin{hazard}[Images do not commute with intersection]\label{haz:quotient-rel}
$B/\sigma$ and $R/\sigma$ are \emph{images}. Three consequences that the source
navigates silently:
\begin{enumerate}\alphlabels\tightlist
\item $B/\sigma$ is a subset of $A/\sigma$, not the quotient of $B$ by
  $\sigma\cap B^{2}$. For $B$ a subuniverse and $\sigma$ a congruence the two
  are canonically isomorphic as algebras, and the source moves between them
  freely. Fix $B/\sigma$ to be the image and prove the isomorphism once.
\item In general $(C_{1}\cap C_{2})/\sigma\subsetneq C_{1}/\sigma\cap
  C_{2}/\sigma$. Lemma~\ref{lem:propagation}(fm) needs equality for a family
  $C_{1},\dots,C_{t}$, and gets it from the extra hypothesis
  $(C_{i}\circ\sigma)\cap B=C_{i}$, i.e.\ that each $C_{i}$ is
  $\sigma$-saturated inside $B$. That is a content step, not notation.
\item $R/\sigma$ need not be stable under anything, and need not be a
  subuniverse of $(\mathbf A/\sigma)^{m}$ unless $\sigma$ is a congruence.
\end{enumerate}
\end{hazard}

\subsection{Rectangularity}

\begin{definition}[Parallelogram property, rectangularity]\label{def:parallelogram}
Let $R\subseteq A_{1}\times\dots\times A_{m}$.
\begin{enumerate}\alphlabels\tightlist
\item The $i$-th variable of $R$ is \emph{rectangular} if for all
  $\bar a,\bar b$,
  \[
    \bar a[i\mapsto b_{i}]\in R,\quad
    \bar b[i\mapsto a_{i}]\in R,\quad
    \bar b\in R
    \;\Longrightarrow\;\bar a\in R,
  \]
  where $\bar a[i\mapsto c]$ is $\bar a$ with its $i$-th entry replaced by $c$.
  $R$ is \emph{rectangular} if every variable is.
\item $R$ \emph{has the parallelogram property} if for every permutation of the
  coordinates giving $R'$ and every $\ell\in[m-1]$,
  \[
    (a_{1},\dots,a_{\ell},b_{\ell+1},\dots,b_{m})\in R',\quad
    (b_{1},\dots,b_{\ell},a_{\ell+1},\dots,a_{m})\in R',\quad
    \bar b\in R'
    \;\Longrightarrow\;\bar a\in R'.
  \]
\item The \emph{rectangular closure} of $R$ is the least rectangular relation
  containing $R$.
\end{enumerate}
\end{definition}

\begin{lemma}[Two assertions the source makes in passing]\label{lem:rect-basics}
\begin{enumerate}\alphlabels\tightlist
\item A relation with the parallelogram property is rectangular.
\item The rectangular closure exists.
\end{enumerate}
\end{lemma}

\begin{proof}
(a) Apply the parallelogram property to the permutation carrying coordinate
$i$ to position $1$, with $\ell=1$. (b) Rectangularity at coordinate $i$ is a
Horn implication, so the rectangular relations containing $R$ are closed under
arbitrary intersection and the family is nonempty (it contains
$A_{1}\times\dots\times A_{m}$); the intersection of all of them is the least
one.
\end{proof}

\begin{hazard}\label{haz:rect-closure}
Both parts are asserted without proof in \cite{ZhukSimplified} --- (a) as ``as
it follows from the definitions'', (b) by the definite article in ``\emph{the}
minimal rectangular relation''. Both are easy and both must be present. Note
that ``$R$ is rectangular'' unqualified --- the form used in ``rectangular
closure'' --- is never defined in the source; it means every variable is
rectangular.
\end{hazard}

\begin{hazard}[Quantifying over permutations]\label{haz:permutations}
The parallelogram property is stated ``any permutation of its variables gives a
relation $R'$ satisfying \dots''. Formally this is a quantification over
$\mathfrak S_{m}$ combined with a quantification over the split point $\ell$;
equivalently, and more usably, over all ways of splitting $[m]$ into two
nonempty blocks. The second formulation is the one to formalize: it avoids
transporting $R$ along a permutation, which in a dependently typed setting
means an equivalence of index types and a coercion at every use.
\end{hazard}

\subsection{Congruences}

\begin{definition}[Stability]\label{def:stable}
Let $R\subseteq A_{1}\times\dots\times A_{m}$ and let $\sigma$ be a binary
relation on $A_{i}$. The $i$-th variable of $R$ is \emph{stable under $\sigma$}
if $\bar a\in R$ and $(a_{i},b)\in\sigma$ imply $\bar a[i\mapsto b]\in R$.
$R$ is \emph{stable under $\sigma$} if every variable is (each with respect to
$\sigma$ on the corresponding factor, when the factors agree).
\end{definition}

\begin{definition}[Congruence]\label{def:congruence}
A \emph{congruence} on $\mathbf A$ is an equivalence relation on $A$ that is a
subuniverse of $\mathbf A\times\mathbf A$. It is \emph{nontrivial} if it is
neither the equality relation $0_{\mathbf A}$ nor $A^{2}$. The quotient
$\mathbf A/\sigma$ is the algebra on $A/\sigma$ with
$w(\bar a/\sigma)=w(\bar a)/\sigma$.
\end{definition}

\begin{definition}[Irreducible congruence and its cover]\label{def:irreducible}
A congruence $\sigma$ on $\mathbf A$ is \emph{irreducible} if it cannot be
written as an intersection of binary subuniverses of $\mathbf A\times\mathbf A$
each stable under $\sigma$ and each properly containing $\sigma$. Equivalently:
there are no $\mathbf S_{1},\dots,\mathbf S_{k}\le\mathbf A/\sigma\times
\mathbf A/\sigma$ with $0_{\mathbf A/\sigma}=S_{1}\cap\dots\cap S_{k}$ and
$S_{i}\neq 0_{\mathbf A/\sigma}$ for every $i$.

For irreducible $\sigma$, $\cover{\sigma}$ denotes the least
$\delta\le\mathbf A\times\mathbf A$ with $\delta\supsetneq\sigma$ and $\delta$
stable under $\sigma$.
\end{definition}

\begin{hazard}[Three traps in Definition~\ref{def:irreducible}]\label{haz:irreducible}
\begin{enumerate}\alphlabels\tightlist
\item Irreducibility is \emph{not} meet-irreducibility in the congruence
  lattice. The $\mathbf S_{i}$ range over binary \emph{subuniverses} stable
  under $\sigma$, which need not be equivalence relations. The two notions
  differ, and every use below is the stated one.
\item $\cover{\sigma}$ need not be a congruence. It is a subuniverse of
  $\mathbf A^{2}$ stable under $\sigma$, nothing more; that $\cover{\sigma}$
  \emph{is} a congruence is an extra hypothesis in the definition of a linear
  congruence (Definition~\ref{def:linear-congruence}), which would be
  redundant otherwise.
\item Existence and uniqueness of $\cover{\sigma}$ need an argument, given in
  Proposition~\ref{prop:cover}. Uniqueness is exactly what irreducibility buys;
  without it there is no ``the'' cover. A formalization must produce
  $\cover{\sigma}$ as data attached to a proof of irreducibility, not as a
  standalone function.
\item The family $S_{1},\dots,S_{k}$ may be \emph{empty}
  (item \ref{leg:emptyfamily} of
  \S\ref{sec:legislation}). This is not a free choice: the two formulations
  above agree only under it.
\end{enumerate}
\end{hazard}

\begin{proposition}[The cover exists, and is a tolerance]\label{prop:cover}
Let $\sigma$ be an irreducible congruence on $\mathbf A$. Then
\begin{enumerate}\alphlabels\tightlist
\item $\sigma\neq A^{2}$;
\item the family $\mathcal F$ of $\delta\le\mathbf A\times\mathbf A$ with
  $\delta\supsetneq\sigma$ and $\delta$ stable under $\sigma$ is nonempty and
  closed under intersection, so it has a least element $\cover{\sigma}$;
\item $\cover{\sigma}$ is reflexive and symmetric.
\end{enumerate}
\end{proposition}

\begin{proof}
(a) The empty family is allowed, and its intersection is $A^{2}$; so if
$\sigma=A^{2}$ then $\sigma$ is an intersection of a family of relations each
properly containing it, vacuously, and $\sigma$ is reducible.

(b) $\mathcal F\ni A^{2}$ by (a). If $\delta_{1},\delta_{2}\in\mathcal F$ then
$\delta_{1}\cap\delta_{2}$ is again a subuniverse of $\mathbf A^{2}$ stable
under $\sigma$ and contains $\sigma$; it contains $\sigma$ \emph{properly},
since otherwise $\sigma=\delta_{1}\cap\delta_{2}$ exhibits $\sigma$ as an
intersection of two members of $\mathcal F$, contradicting irreducibility. So
$\mathcal F$ is closed under binary, hence finite, intersection, and
$\cover{\sigma}=\bigcap\mathcal F\in\mathcal F$ is its least element.

(c) $\cover{\sigma}\supseteq\sigma\supseteq 0_{\mathbf A}$ is reflexive. For
symmetry, $(\cover{\sigma})^{-1}$ is a subuniverse of $\mathbf A^{2}$, is
stable under $\sigma$ because $\sigma$ is symmetric, and properly contains
$\sigma^{-1}=\sigma$; so $(\cover{\sigma})^{-1}\in\mathcal F$ and therefore
$\cover{\sigma}\subseteq(\cover{\sigma})^{-1}$, whence equality.
\end{proof}

\begin{remark}[Why not a congruence]\label{rem:cover-not-cong}
Nothing in Proposition~\ref{prop:cover} gives transitivity, and the source
supplies the internal evidence that none is available: that $\cover{\sigma}$ is
a congruence is item (b) of Definition~\ref{def:linear-congruence} and
conclusion (a) of Lemma~\ref{lem:nontrivial-bridge}, both of which would be
redundant otherwise. Symmetry of $\cover{\sigma}$, by contrast, is used
silently in \cite{ZhukSimplified} --- in the proof of
Lemma~\ref{lem:pc-bridges-trivial}, where the argument is run again ``switching
$x_{1}$ and $x_{2}$''.
\end{remark}

\begin{lemma}[Passing to the quotient]\label{lem:irred-quotient}
Let $\sigma$ be a congruence on $\mathbf A$. Then $\delta\mapsto\delta/\sigma$
is an inclusion-preserving bijection, commuting with intersections, from the
subuniverses of $\mathbf A^{2}$ that contain $\sigma$ and are stable under
$\sigma$ onto the subuniverses of $(\mathbf A/\sigma)^{2}$, sending $\sigma$ to
$0_{\mathbf A/\sigma}$. Consequently $\sigma$ is irreducible, linear or PC
exactly when $0_{\mathbf A/\sigma}$ is, and then
$\cover{\sigma}/\sigma=\cover{(0_{\mathbf A/\sigma})}$.
\end{lemma}

\begin{proof}
A subuniverse $\delta$ of $\mathbf A^{2}$ containing $\sigma$ and stable under
$\sigma$ is a union of products of $\sigma$-blocks, so it is the preimage of
$\delta/\sigma$; conversely the preimage of a subuniverse of
$(\mathbf A/\sigma)^{2}$ is a subuniverse containing $\sigma$ and stable under
$\sigma$. The two constructions are mutually inverse and preserve inclusion and
intersection, and $\sigma$ is the preimage of $0_{\mathbf A/\sigma}$. The
statement of Definition~\ref{def:irreducible} is a statement about that family
alone, so it transports; and both the least element in
Proposition~\ref{prop:cover}(b) and the data of
Definition~\ref{def:linear-congruence} transport with it.
\end{proof}

\begin{hazard}[This is the licence for ``assume $\sigma=0$'']\label{haz:factor-out}
Three proofs in \S\ref{sec:properties} open with ``we may factor by $\sigma$
and assume $\sigma$ is the equality relation''.
Lemma~\ref{lem:irred-quotient} is what licenses that, and it is one of the
sentences \cite{ZhukSimplified} states without proof
(\S\ref{sec:legislation}, item \ref{leg:notice}).
\end{hazard}

\subsection{Bridges}

\begin{definition}[Bridge]\label{def:bridge}
Let $\sigma_{1},\sigma_{2}$ be congruences on $\mathbf D_{1},\mathbf D_{2}$. A
relation $\delta\le\mathbf D_{1}^{2}\times\mathbf D_{2}^{2}$ is a \emph{bridge
from $\sigma_{1}$ to $\sigma_{2}$} if
\begin{enumerate}\alphlabels\tightlist
\item the first two variables of $\delta$ are stable under $\sigma_{1}$;
\item the last two variables of $\delta$ are stable under $\sigma_{2}$;
\item $\proj_{1,2}(\delta)\supsetneq\sigma_{1}$ and
      $\proj_{3,4}(\delta)\supsetneq\sigma_{2}$;
\item $(a_{1},a_{2},a_{3},a_{4})\in\delta$ implies
      $\bigl((a_{1},a_{2})\in\sigma_{1}\iff(a_{3},a_{4})\in\sigma_{2}\bigr)$.
\end{enumerate}
For a bridge $\delta$ put $\bcol{\delta}=\{(x,y):(x,x,y,y)\in\delta\}$.
A bridge $\delta\subseteq D^{4}$ is \emph{reflexive} if $(a,a,a,a)\in\delta$
for all $a\in D$.
\end{definition}

\begin{remark}[The motivating example]\label{rem:bridge-example}
On $\mathbb Z_{4}$, the relation
$\delta=\{(a_{1},a_{2},a_{3},a_{4}):a_{1}-a_{2}=2a_{3}-2a_{4}\}$ is a bridge
from $0_{\mathbb Z_{4}}$ to congruence mod~$2$. Condition (d) is the content:
$a_{1}=a_{2}$ exactly when $a_{3}\equiv a_{4}\pmod 2$. Bridges are the
formal shadow of a centralizer relation; see \cite{RossSlides}.
\end{remark}

\begin{definition}[Composition of bridges]\label{def:bridge-comp}
For bridges $\delta_{1}$ from $\sigma_{0}$ to $\sigma_{1}$ and $\delta_{2}$
from $\sigma_{1}$ to $\sigma_{2}$, set
\[
  (\delta_{1}\ast\delta_{2})(x_{1},x_{2},z_{1},z_{2})
  =\exists y_{1}\exists y_{2}\;
   \delta_{1}(x_{1},x_{2},y_{1},y_{2})\wedge\delta_{2}(y_{1},y_{2},z_{1},z_{2}).
\]
\end{definition}

\begin{lemma}[Bridge composition; {\cite[Lem.~6.3]{ZhukJACM}}]\label{lem:bridge-comp}
If $\sigma_{0},\sigma_{1},\sigma_{2}$ are irreducible then
$\delta_{1}\ast\delta_{2}$ is a bridge from $\sigma_{0}$ to $\sigma_{2}$, and
$\bcol{\delta_{1}\ast\delta_{2}}=\bcol{\delta_{1}}\circ\bcol{\delta_{2}}$.
\end{lemma}

\begin{hazard}\label{haz:bridge-comp}
Lemma~\ref{lem:bridge-comp} is imported, not proved, in
\cite{ZhukSimplified}; it is \cite[Lemma 6.3]{ZhukJACM}. Both conclusions are
used: the first in every chain of bridges, the second --- the identity
$\bcol{\cdot}$ commutes with composition --- in
Lemma~\ref{lem:perfect-from-linked}, where a bridge with $\bcol{\delta}$ linked
is powered up until $\bcol{\delta}=A^{2}$. The irreducibility hypothesis on
\emph{all three} congruences is needed and is easy to drop by accident.
\end{hazard}

\begin{definition}[Symmetric and pair-reflexive bridges]\label{def:pair-reflexive}
Let $\delta$ be a bridge from $\sigma$ to $\sigma$ on $\mathbf D$. Call $\delta$
\emph{symmetric} if
$\delta(x_{1},x_{2},x_{3},x_{4})\iff\delta(x_{3},x_{4},x_{1},x_{2})$, and
\emph{pair-reflexive} if
\[
  (x_{1},x_{2})\in\proj_{1,2}(\delta)\;\Longrightarrow\;
  (x_{1},x_{2},x_{1},x_{2})\in\delta .
\]
Write $\delta^{-1}(x_{1},x_{2},x_{3},x_{4})=\delta(x_{3},x_{4},x_{1},x_{2})$.
\end{definition}

\begin{lemma}[Symmetrisation]\label{lem:symmetrise}
Let $\delta$ be a reflexive bridge from $\sigma$ to $\sigma$ on $\mathbf D$ and
put $\delta^{\circ}=\delta\ast\delta^{-1}$, that is
\[
  \delta^{\circ}(x_{1},x_{2},x_{3},x_{4})=\exists x_{5}\exists x_{6}\;
  \delta(x_{1},x_{2},x_{5},x_{6})\wedge\delta(x_{3},x_{4},x_{5},x_{6}).
\]
Then $\delta^{\circ}$ is a reflexive, symmetric, pair-reflexive bridge from
$\sigma$ to $\sigma$ with
$\proj_{1,2}(\delta^{\circ})=\proj_{3,4}(\delta^{\circ})=\proj_{1,2}(\delta)$
and $\bcol{\delta^{\circ}}=\bcol{\delta}\circ\bcol{\delta}^{-1}
\supseteq\bcol{\delta}$. Moreover any pair-reflexive bridge $\varepsilon$
satisfies $\varepsilon\subseteq\varepsilon\ast\varepsilon$.
\end{lemma}

\begin{proof}
Symmetry is visible in the formula. For pair-reflexivity, if
$(x_{1},x_{2})\in\proj_{1,2}(\delta^{\circ})$ then
$\delta(x_{1},x_{2},x_{5},x_{6})$ for some $x_{5},x_{6}$, and taking the same
witnesses twice gives $(x_{1},x_{2},x_{1},x_{2})\in\delta^{\circ}$; the same
computation shows $\proj_{1,2}(\delta^{\circ})=\proj_{1,2}(\delta)$, the
inclusion $\subseteq$ being immediate, and symmetry then gives
$\proj_{3,4}(\delta^{\circ})=\proj_{1,2}(\delta^{\circ})$. Clause (c) of
Definition~\ref{def:bridge} follows, and clause (d) because
$(x_{1},x_{2})\in\sigma\iff(x_{5},x_{6})\in\sigma\iff(x_{3},x_{4})\in\sigma$.
Stability is inherited from $\delta$, and reflexivity from
$(a,a,a,a)\in\delta$. The identity
$\bcol{\delta^{\circ}}=\bcol{\delta}\circ\bcol{\delta^{-1}}
=\bcol{\delta}\circ\bcol{\delta}^{-1}$ is Lemma~\ref{lem:bridge-comp};
it contains $\bcol{\delta}$ because $\bcol{\delta}$ is reflexive, so
$(x,y)\in\bcol{\delta}$ gives $(x,y)\in\bcol{\delta}\circ\bcol{\delta}^{-1}$ by
taking $y$ as the intermediate point. Finally, if $\varepsilon$ is
pair-reflexive and $\varepsilon(x_{1},x_{2},x_{3},x_{4})$, then taking
$(x_{1},x_{2})$ itself as the intermediate pair --- legitimate because
$\varepsilon(x_{1},x_{2},x_{1},x_{2})$ --- gives
$(\varepsilon\ast\varepsilon)(x_{1},x_{2},x_{3},x_{4})$.
\end{proof}

\begin{hazard}[Pair-reflexivity is the hypothesis three lemmas need]\label{haz:pair-reflexive}
\cite{ZhukSimplified} states pair-reflexivity only once, as a hypothesis of
Lemma~\ref{lem:pc-inductive-step}, but three further statements require it and
do not say so; without it one of them is false. See
Hazard~\ref{haz:bridge-repairs}. Lemma~\ref{lem:symmetrise} is what makes this
costless: every bridge these statements are ever applied to is a
$\delta^{\circ}$.
\end{hazard}

\begin{definition}[Perfect linear congruence]\label{def:perfect-linear}
A congruence $\sigma$ on $\mathbf A$ is a \emph{perfect linear congruence} if
it is irreducible and there are a prime $p$ and
$\zeta\le\mathbf A\times\mathbf A\times\mathbf Z_{p}$ with
$\proj_{1,2}(\zeta)=\cover{\sigma}$ and
\[
  (a_{1},a_{2},b)\in\zeta\;\Longrightarrow\;
  \bigl((a_{1},a_{2})\in\sigma\iff b=0\bigr).
\]
\end{definition}

Perfect linear congruences are the payoff of the whole bridge machinery: they
let the passage from $\sigma$ to $\cover{\sigma}$ be tracked by a single
element of $\mathbf Z_{p}$, which is what turns a CSP instance into a system of
linear equations in \S\ref{sec:codim-one}.

\clearpage
%!TeX root=csp.tex
\section{The six types of strong subuniverse}\label{sec:types}

This section is Zhuk's \S2.2. It defines six binary relations
$C\subt[A]{T}B$ between subsets of a fixed algebra $\mathbf A$, one for each
type $T$, and the transitive-closure relation $C\lll^{A}B$. Everything in
\S\S\ref{sec:properties}--\ref{sec:main-statements} is expressed in this
language, so the definitions repay care.

The informal reading is: $C\subt[A]{T}B$ says that $C$ is a \emph{strong}
subset of $B$ --- strong enough that the CSP algorithm may reduce a variable's
domain from $B$ to $C$ without losing all solutions --- and $T$ records
\emph{why} it is strong. Three of the reasons are absorption-theoretic
($\TBA$, $\TC$, $\TS$) and three are congruence-theoretic
($\TD$, $\TL$, $\TPC$).

\subsection{Absorption}

\begin{definition}[Absorbing subuniverse]\label{def:absorbing}
A subuniverse $B$ of $\mathbf A$ is \emph{absorbing} if there is
$t\in\Clo_{k}(\mathbf A)$, for some $k$, such that for every $i\in[k]$
\[
  t(B,\dots,B,\underbrace{A}_{i},B,\dots,B)\subseteq B .
\]
We then say \emph{$B$ absorbs $\mathbf A$ with the term $t$}. If $t$ may be
chosen binary, $B$ is a \emph{binary absorbing} subuniverse ($\TBA$); if
ternary, a \emph{ternary absorbing} subuniverse.
\end{definition}

\begin{hazard}[The tuple form]\label{haz:absorption-tuples}
Definition~\ref{def:absorbing} must be formalized over \emph{tuples
constrained coordinatewise}, not as ``for all $b_{1},\dots,b_{k}\in B$ and
$a\in A$, $t(b_{1},\dots,b_{i-1},a,b_{i+1},\dots,b_{k})\in B$''. The two agree,
but the second phrasing invites the variant in which $b_{i}$ is quantified and
then discarded, which is vacuously true when $B=\varnothing$ and at $k=1$
where the intended content is not vacuous. This was an actual defect in the
predecessor blueprint. State it as: for every $\bar a\in A^{k}$ such that
$a_{j}\in B$ for all $j\neq i$, one has $t(\bar a)\in B$.
\end{hazard}

\begin{definition}[Central subuniverse]\label{def:central}
A subuniverse $C$ of $\mathbf A$ is \emph{central} ($\TC$) if it is absorbing
and for every $a\in A\setminus C$,
\[
  (a,a)\notin\Sg_{\mathbf A}\bigl((\{a\}\times C)\cup(C\times\{a\})\bigr).
\]
\end{definition}

\begin{lemma}[{\cite[Cor.~6.11.1]{ZhukStrong}}]\label{lem:central-ternary}
A central subuniverse is a ternary absorbing subuniverse.
\end{lemma}

\begin{remark}[Already formalized]\label{rem:center-done}
Lemma~\ref{lem:central-ternary} is the one deep import of this theory that is
already available in Lean: it is the main theorem of the predecessor project
(\texttt{zhuk\_main} in \texttt{ZhukLean}), proved there from the definitions
via the Barto--Kazda relational description of absorption and the Zhuk--Kozik
doubling trick. See \S\ref{sec:reuse}.
\end{remark}

\begin{definition}[BA and centre free; $\TS$]\label{def:ba-center-free}
$\mathbf A$ is \emph{BA and centre free} if it has no proper nonempty binary
absorbing subuniverse and no proper nonempty central subuniverse.

For $\varnothing\neq C\lneq B\le A$ write $C\subt[A]{\TS}B$ if there is a
subuniverse $D\subseteq C$ of $\mathbf B$ that is simultaneously binary
absorbing and central in $\mathbf B$. $\mathbf A$ is \emph{$\TS$-free} if there
is no $C\subt{\TS}\mathbf A$; equivalently, no nonempty proper $D\le A$ is
both binary absorbing and central.
\end{definition}

\begin{hazard}[$\TS$ is not the conjunction of $\TBA$ and $\TC$]\label{haz:S-type}
$C\subt{\TS}B$ does not say that $C$ is binary absorbing and central in
$\mathbf B$; it says that $C$ \emph{contains} some $D$ that is. The type
$\TS$ exists precisely so that it is closed under enlarging $C$, which
$\TBA$ and $\TC$ are not. Note also the asymmetry with $\TS$-freeness, which
\emph{is} stated as a conjunction. Getting this backwards makes
Lemma~\ref{lem:propagation}(ft) unprovable.
\end{hazard}

\subsection{Linear and PC congruences}

\begin{definition}[Linear congruence]\label{def:linear-congruence}
A congruence $\sigma$ on $\mathbf A$ is \emph{linear} if
\begin{enumerate}\alphlabels\tightlist
\item $\sigma$ is irreducible;
\item $\cover{\sigma}$ is a congruence;
\item there are a prime $p$ and $S\le(\cover{\sigma})^{[4]}$ such that for
  every block $B$ of $\cover{\sigma}$ there is $m\ge 0$ with
  \[
    \bigl(B/\sigma;\;S\cap(B/\sigma)^{4}\bigr)\;\cong\;
    \bigl(\mathbb Z_{p}^{m};\;x_{1}-x_{2}=x_{3}-x_{4}\bigr).
  \]
\end{enumerate}
An irreducible congruence that is not linear is a \emph{PC congruence}.
\end{definition}

The relation $S$ of (c) is a bridge from $\sigma$ to $\sigma$ with
$\bcol{S}=\proj_{1,2}(S)=\proj_{3,4}(S)=\cover{\sigma}$. This is the useful
reformulation:

\begin{lemma}[{\cite[Lem.~\ref*{lem:linear-equiv}]{ZhukSimplified}}]\label{lem:linear-equiv}
For an irreducible congruence $\sigma$ the following are equivalent:
\begin{enumerate}\alphlabels\tightlist
\item $\sigma$ is linear;
\item there is a bridge $\delta$ from $\sigma$ to $\sigma$ with
  $\bcol{\delta}\supsetneq\sigma$.
\end{enumerate}
\end{lemma}

\begin{hazard}[Definition (c) versus criterion (b)]\label{haz:linear-def}
Definition~\ref{def:linear-congruence}(c) is a per-block isomorphism statement
with an existential $m$ \emph{depending on the block}, and it quantifies over
blocks of $\cover{\sigma}$, not of $\sigma$. Lemma~\ref{lem:linear-equiv}
replaces all of that by one bridge. A formalization should take (b) as the
\emph{definition} and prove (c) as a structure theorem, reversing the paper's
order: (b) is a $\Sigma_1$ statement over finite data and is what every
consumer actually uses, while (c) drags in a quotient, a choice of prime, and a
family of isomorphisms.

The dichotomy ``linear or PC'' is then \emph{by definition} exhaustive on
irreducible congruences, which is how the paper uses it. Note that PC-ness is a
negative condition; every lemma about PC congruences is proved by contradiction
against Lemma~\ref{lem:linear-equiv}(b).
\end{hazard}

\begin{lemma}[No bridge crosses the types]\label{lem:no-cross-bridge}
If $\sigma_{1}$ is linear, $\sigma_{2}$ is irreducible, and there is a bridge
from $\sigma_{1}$ to $\sigma_{2}$, then $\sigma_{2}$ is linear.
\end{lemma}

\begin{lemma}[Bridges from PC congruences are trivial]\label{lem:pc-bridges-trivial}
Let $\sigma$ be a PC congruence on $A$ and $\delta$ a reflexive bridge from
$\sigma$ to $\sigma$ with $\proj_{1,2}(\delta)=\proj_{3,4}(\delta)=
\cover{\sigma}$. Then
\[
  \delta(x_{1},x_{2},x_{3},x_{4})=\sigma(x_{1},x_{3})\wedge\sigma(x_{2},x_{4})
  \quad\text{or}\quad
  \delta(x_{1},x_{2},x_{3},x_{4})=\sigma(x_{1},x_{4})\wedge\sigma(x_{2},x_{3}).
\]
\end{lemma}

\begin{lemma}[Bridges into PC congruences]\label{lem:bridge-to-pc}
Let $\delta$ be a bridge from a PC congruence $\sigma_{1}$ on $\mathbf A_{1}$
to an irreducible congruence $\sigma_{2}$ on $\mathbf A_{2}$, with
$\proj_{1,2}(\delta)=\cover{\sigma_{1}}$ and
$\proj_{3,4}(\delta)=\cover{\sigma_{2}}$. Then
\begin{enumerate}\alphlabels\tightlist
\item $\sigma_{2}$ is a PC congruence;
\item $\mathbf A_{1}/\sigma_{1}\cong\mathbf A_{2}/\sigma_{2}$;
\item $\{(a/\sigma_{1},b/\sigma_{2}):(a,b)\in\bcol{\delta}\}$ is bijective;
\item $\delta(x_{1},x_{2},x_{3},x_{4})=
      \bcol{\delta}(x_{1},x_{3})\wedge\bcol{\delta}(x_{2},x_{4})$ or
      $\delta(x_{1},x_{2},x_{3},x_{4})=
      \bcol{\delta}(x_{1},x_{4})\wedge\bcol{\delta}(x_{2},x_{3})$.
\end{enumerate}
\end{lemma}

Two results connect the new vocabulary to the linear and PC subuniverses of the
original proof.

\begin{lemma}\label{lem:linear-on-top}
If $\sigma$ is a linear congruence on $\mathbf A\in\Vn$ with
$\cover{\sigma}=A^{2}$, then $\mathbf A/\sigma\cong\mathbf Z_{p}$ for a prime
$p$.
\end{lemma}

\begin{lemma}\label{lem:pc-on-top}
If $\sigma$ is a PC congruence on $\mathbf A$ with $\cover{\sigma}=A^{2}$, then
$\mathbf A/\sigma$ is polynomially complete.
\end{lemma}

\begin{hazard}[Polynomial completeness]\label{haz:pc-algebra}
An algebra is \emph{polynomially complete} if the clone generated by its basic
operations together with all constants is the clone of \emph{all} operations on
its universe. Lemma~\ref{lem:pc-on-top} is where the name ``PC'' comes from,
and it is the only place in the development where polynomial completeness is
used as more than a label. Formalizing it needs the Istinger--Kaiser
characterization \cite{IstingerKaiser}; formalizing everything else needs only
the \emph{name}. A staged development should therefore treat ``PC congruence''
as ``irreducible and not linear'' throughout, and prove
Lemma~\ref{lem:pc-on-top} last or not at all --- it is used only to connect
with \cite{ZhukJACM}, not by any argument here.
\end{hazard}

\subsection{The types}

\begin{definition}[The six types]\label{def:types}
Let $\varnothing\neq C\lneq B\le A$. We write:
\begin{itemize}\tightlist
\item $C\subt[A]{\TBA}B$ if $C$ is a binary absorbing subuniverse of
  $\mathbf B$;
\item $C\subt[A]{\TC}B$ if $C$ is a central subuniverse of $\mathbf B$;
\item $C\subt[A]{\TD}B$ if there is an irreducible congruence $\sigma$ on
  $\mathbf A$ with
  \begin{enumerate}\romanlabels\tightlist
  \item $B^{2}\subseteq\cover{\sigma}$,
  \item $C=B\cap E$ for some block $E$ of $\sigma$,
  \item $\mathbf B/\sigma$ is BA and centre free;
  \end{enumerate}
\item $C\subt[A]{\TL}B$ if $C\subt[A]{\TD}B$ with the witnessing $\sigma$
  linear;
\item $C\subt[A]{\TPC}B$ if $C\subt[A]{\TD}B$ with the witnessing $\sigma$ a PC
  congruence;
\item $C\subt[A]{\TS}B$ as in Definition~\ref{def:ba-center-free}.
\end{itemize}
When the witnessing congruence matters we write $C\subt[A]{T(\sigma)}B$; for
$T\in\{\TBA,\TC,\TS\}$, where no congruence is involved, $\sigma$ is taken to
be $A^{2}$.
\end{definition}

\begin{hazard}[The superscript $A$ is not decoration]\label{haz:superscript}
For $T\in\{\TD,\TL,\TPC\}$ the witnessing congruence lives on $\mathbf A$, not
on $\mathbf B$, and condition (i) relates the two: $B^{2}\subseteq
\cover{\sigma}$ says $B$ sits inside a single block of $\cover{\sigma}$. This
is why the relation carries the ambient algebra as a superscript. For
$T\in\{\TBA,\TC,\TS\}$ the superscript is genuinely irrelevant and Zhuk drops
it. A formalization should not overload: define $\subt[A]{T}$ with $A$ always
present, and prove the independence for the three absorption types as a lemma.

Note also that $\mathbf B/\sigma$ in (iii) means $B/\sigma$ with the induced
algebra structure, which requires $B$ to be a subuniverse and $\sigma$
restricted to $B$; since $B^{2}\subseteq\cover{\sigma}$ and not necessarily
$B^{2}\subseteq\sigma$, the quotient $B/\sigma$ is in general \emph{not} a
single point, and its being BA and centre free is a real condition.
\end{hazard}

\begin{definition}[Multi-types]\label{def:multitypes}
For $T\in\{\TL,\TPC,\TD\}$ write $C\subt[A]{\MT T}B$ if $C\neq\varnothing$ and
$C=C_{1}\cap\dots\cap C_{t}$ with $C_{i}\subt[A]{T}B$ for every $i\in[t]$.
\end{definition}

\begin{definition}[The strong-reduction relation]\label{def:lll}
Write $C\lll^{A}B$ if there are $B_{0},\dots,B_{m}\subseteq B$ and
$T_{1},\dots,T_{m}\in\{\TBA,\TC,\TS,\TD\}$ with
\[
  C=B_{m}\subt[A]{T_{m}}B_{m-1}\subt[A]{T_{m-1}}\dots\subt[A]{T_{1}}B_{0}=B .
\]
Here $m=0$ is allowed, so $\lll^{A}$ is reflexive. A congruence \emph{comes
from} $C\lll^{A}B$ if it is one of the witnessing congruences in such a chain.
We write $B\lll A$ for $B\lll^{A}A$, and $C\subte[A]{T(\sigma)}B$ for
``$C=B$ or $C\subt[A]{T(\sigma)}B$''.
\end{definition}

\begin{hazard}[$\lll$ is existential over chains]\label{haz:lll-chain}
$C\lll^{A}B$ asserts the \emph{existence} of a chain; the intermediate types
are forgotten. This is deliberate --- it is the paper's central simplification
over \cite{ZhukJACM}, where the types in the middle had to be tracked --- and
it is why $\lll$ is usable as a hypothesis. In Lean it is the reflexive
transitive closure of $\bigcup_{T\in\{\TBA,\TC,\TS,\TD\}}\subt[A]{T}$, i.e.\
\texttt{Relation.ReflTransGen}, and the induction principle that comes with
that is exactly the one the proofs use. Note that $\TL$ and $\TPC$ are
\emph{absent} from the list because $\TD$ subsumes them.
\end{hazard}

\begin{convention}[Dotted variants]\label{conv:dotted}
$C\dsubt[A]{T}B$, $C\dsubte[A]{T}B$ and $C\dlll^{A}B$ mean that the undotted
relation holds \emph{or} $C=\varnothing$. See Hazard~\ref{haz:dot}.
\end{convention}

\clearpage
%!TeX root=csp-proof.tex
\section{Instances}\label{sec:instances}

This section is entirely combinatorial: no algebra is used beyond the fact that
each domain carries a $w$ and each constraint relation is a subuniverse of a
product of domains. It is also where the most representation decisions are
forced, because instances are syntactic objects manipulated by operations ---
weakening, covering, renaming --- that are usually described in words.

\subsection{Instances and solutions}

\begin{definition}[Instance]\label{def:instance}
Fix a finite set $V$ of variables and, for each $x\in V$, an algebra
$\mathbf D_{x}\in\Vn$. A \emph{constraint} is a pair
$C=(\bar x,R)$ where $\bar x=(x_{1},\dots,x_{m})$ is a tuple of variables and
$R\le\mathbf D_{x_{1}}\times\dots\times\mathbf D_{x_{m}}$. An
\emph{instance} $\mathcal I$ is a finite list of constraints; we write
$C\in\mathcal I$, and $\Var(C)$, $\Var(\mathcal I)$ for the variables
occurring. A \emph{subinstance} is a sublist.

A \emph{solution} is a map $s$ with $s(x)\in D_{x}$ for all $x$ such that
$(s(x_{1}),\dots,s(x_{m}))\in R$ for every constraint. The solution set of
$\mathcal I$ is \emph{subdirect} if for every $x$ and every $a\in D_{x}$ there
is a solution with $s(x)=a$.
\end{definition}

\begin{hazard}[Lists, not sets; tuples, not sets of variables]\label{haz:instance-list}
Two representation choices are forced by the way the arguments below go, and
both differ from the naive one.

\emph{Constraints form a list, not a set.} Proofs delete ``a constraint
$C\in\mathcal I$'' and reason about $\mathcal I\setminus\{C\}$, and the same
relation may occur twice on different variable tuples --- and, in the expanded
coverings of \S\ref{sec:coverings}, the same relation on the \emph{same}
tuple can legitimately occur twice. Multiset or list semantics is required.

\emph{A constraint's scope is a tuple, not a set.} Repeated variables in a
scope are allowed and occur --- $\zeta(x,x',z)$ in
\S\ref{sec:main-statements} arises by splitting one variable into two, and
$R(x,x,\dots,x)$ appears explicitly in Definition~\ref{def:expanded}. So
$\Var(C)$ is the image of the scope, and $|{\Var(C)}|$ may be smaller than the
arity.

Once repeated variables are allowed, every semantic operation below that names
a \emph{variable} rather than a \emph{position} --- $\proj_{z}(C)$ in
$1$-consistency, $\proj_{z_{i},z_{i+1}}(C)$ along a path, $\ConOne(C,x)$,
adjacency at a common variable, weakening by a set of variable names, and the
projections in Definition~\ref{def:instance-linked}(ii) --- must be read as
pulled back along the equal-variable diagonal first. Otherwise a constraint
$R(x,x)$ has a full first-coordinate projection even when $R$ contains no
diagonal tuple and the constraint admits no value for $x$. Fixing this notation
uniformly is work this document has not yet done; see \S\ref{sec:status}.
\end{hazard}

\begin{definition}[Relations defined by an instance]\label{def:instance-relation}
For $x_{1},\dots,x_{m}\in\Var(\mathcal I)$, $\mathcal I(x_{1},\dots,x_{m})$ is
the set of $(a_{1},\dots,a_{m})$ such that $\mathcal I$ has a solution with
$s(x_{i})=a_{i}$ for all $i$. It is a subuniverse of
$\mathbf D_{x_{1}}\times\dots\times\mathbf D_{x_{m}}$, being defined by a
primitive positive formula over the constraint relations.
\end{definition}

\subsection{Reductions}

\begin{definition}[Reduction]\label{def:reduction}
A \emph{reduction} for $\mathcal I$ is a map $D^{(\top)}$ assigning to each
$x\in\Var(\mathcal I)$ a subuniverse $D^{(\top)}_{x}\subseteq D_{x}$, possibly
empty. It is \emph{nonempty} if every $D^{(\top)}_{x}\neq\varnothing$. The
instance $\mathcal I^{(\top)}$ is $\mathcal I$ with each domain replaced by
$D^{(\top)}_{x}$ and each constraint relation intersected with the
corresponding box. For reductions we write $D^{(\bot)}\lll D^{(\top)}$ and
$D^{(\bot)}\subte{T}D^{(\top)}$ when the corresponding relation holds at every
variable.
\end{definition}

\begin{hazard}[Empty reductions]\label{haz:empty-reduction}
Reductions are allowed to be empty at some or all variables --- the maximal
$1$-consistent reduction below $D^{(B,\top)}$ in the proof of
Theorem~\ref{thm:main-inductive} is constructed exactly so that it may come out
empty, and the case split on whether it does is the branch point of the whole
argument. But $D^{(\bot)}\lll D^{(\top)}$ requires nonemptiness at every
variable (Convention~\ref{conv:empty}). Both notions are needed and they must
be different objects.
\end{hazard}

\subsection{Consistency}

\begin{definition}[Paths]\label{def:path}
A \emph{path} in $\mathcal I$ is an alternating sequence
$z_{1}-C_{1}-z_{2}-\dots-C_{l-1}-z_{l}$ with $z_{i},z_{i+1}\in\Var(C_{i})$.
It \emph{connects $b$ and $c$} if there are $a_{i}\in D_{z_{i}}$ with
$a_{1}=b$, $a_{l}=c$, and $(a_{i},a_{i+1})\in\proj_{z_{i},z_{i+1}}(C_{i})$ for
every $i$.
$\mathcal I$ is a \emph{tree-instance} if it has no path with $l\ge 3$,
$z_{1}=z_{l}$ and $C_{1},\dots,C_{l-1}$ pairwise distinct.
\end{definition}

\begin{definition}[Consistency]\label{def:consistency}
$\mathcal I$ is \emph{$1$-consistent} if $\proj_{z}(C)=D_{z}$ for every
constraint $C$ and every $z\in\Var(C)$. A reduction $D^{(\top)}$ is
$1$-consistent for $\mathcal I$ if $\mathcal I^{(\top)}$ is $1$-consistent.
$\mathcal I$ is \emph{cycle-consistent} if it is $1$-consistent and every path
from $z$ to $z$ connects $a$ to $a$, for every variable $z$ and every
$a\in D_{z}$.
\end{definition}

\begin{definition}[Linked, fragmented, irreducible]\label{def:instance-linked}
$\mathcal I$ is \emph{linked} if for every $z\in\Var(\mathcal I)$ and every
$a,b\in D_{z}$ some path from $z$ to $z$ connects $a$ and $b$.
$\mathcal I$ is \emph{fragmented} if $\Var(\mathcal I)$ splits into two disjoint
nonempty sets $X_{1},X_{2}$ with $\Var(C)\subseteq X_{1}$ or
$\Var(C)\subseteq X_{2}$ for every $C$.
$\mathcal I$ is \emph{irreducible} if there is no instance $\mathcal I'$ with
\begin{enumerate}\romanlabels\tightlist
\item $\Var(\mathcal I')\subseteq\Var(\mathcal I)$;
\item every constraint of $\mathcal I'$ is a projection of a constraint of
  $\mathcal I$ onto some of its variables;
\item $\mathcal I'$ is not fragmented;
\item $\mathcal I'$ is not linked;
\item the solution set of $\mathcal I'$ is not subdirect.
\end{enumerate}
\end{definition}

\begin{hazard}[Irreducibility is a $\Pi_{1}$ condition over a large family]
\label{haz:irreducible-instance}
Conditions (i)--(v) quantify over \emph{all} instances built from projections
of constraints of $\mathcal I$ --- a finite but exponentially large family
(any multiset of projections of any constraints onto any subsets of their
variables). It is finite, so the predicate is decidable, but it is not
something to unfold. Every use in \S\ref{sec:main-statements} goes through
Lemma~\ref{lem:crucial-irreducible} and \emph{never} through the definition.
\end{hazard}

\subsection{Weakening and cruciality}

\begin{definition}[Weakening]\label{def:weakening}
A constraint $R_{1}(y_{1},\dots,y_{t})$ is \emph{weaker or equivalent to}
$R_{2}(z_{1},\dots,z_{s})$ if $\{y_{i}\}\subseteq\{z_{j}\}$ and
$R_{2}(\bar z)$ implies $R_{1}(\bar y)$; it is \emph{weaker} if in addition
$R_{1}(\bar y)$ does not imply $R_{2}(\bar z)$. \emph{Weakening} a constraint
$C$ in $\mathcal I$ means replacing $C$ by all constraints weaker than it. An
instance $\mathcal I'$ is a \emph{weakening} of $\mathcal I$ if
$\Var(\mathcal I')\subseteq\Var(\mathcal I)$ and every constraint of
$\mathcal I'$ is weaker or equivalent to a constraint of $\mathcal I$.
\end{definition}

\begin{hazard}[``Replace by all weaker constraints'', and its termination]
\label{haz:weakening}
The operation is finite but enormous: there are $2^{|D|^{t}}$ relations weaker
than a given one on a given scope, and one must also range over all sub-scopes.
No algorithm performs it --- it appears only inside proofs, as a device for
reaching a crucial instance (Remark~\ref{rem:get-crucial}). Weakening is
therefore defined above as a \emph{relation between instances}, not as a
function producing one.

The termination of Remark~\ref{rem:get-crucial} is a genuine obligation and is
\emph{not} discharged here. Two things make it delicate, and both must be
faced before ``we cannot weaken forever'' --- which is used three times in
\S\ref{sec:main-statements} --- can be relied on.

\emph{The obvious measure points the wrong way.} A properly weaker constraint
admits \emph{more} tuples, so the total number of tuples in all constraint
relations \emph{increases}. A measure with the right orientation would be a
rank in the finite implication order on constraints, or the number of
assignments forbidden by the instance.

\emph{Replacing a constraint by \emph{all} its strict weakenings need not make
progress.} The conjunction of all strict weakenings of a relation can be the
relation itself --- intersecting all one-tuple extensions of $R$ recovers $R$.
So the construction must choose one proper weakening, or a finite antichain
with a proved rank decrease, and the choice must be part of the statement.

See \S\ref{sec:status}.
\end{hazard}

\begin{definition}[Crucial]\label{def:crucial}
Let $D^{(\top)}$ be a reduction for $\mathcal I$. A variable $y_{i}$ of a
constraint $R(y_{1},\dots,y_{t})$ is \emph{dummy} if $R$ does not depend on its
$i$-th coordinate. A constraint $C\in\mathcal I$ is \emph{crucial in
$D^{(\top)}$} if it has no dummy variables, $\mathcal I^{(\top)}$ has no
solution, and weakening $C$ in $\mathcal I$ yields an instance with a solution
in $D^{(\top)}$. The instance $\mathcal I$ is \emph{crucial in $D^{(\top)}$} if
it has at least one constraint and all of its constraints are crucial in
$D^{(\top)}$.
\end{definition}

\begin{remark}[Reaching a crucial instance]\label{rem:get-crucial}
If $\mathcal I^{(\top)}$ has no solution, then iteratively weakening
constraints that are not crucial, dropping dummy variables as they appear,
terminates in an instance crucial in $D^{(\top)}$. Every relation so introduced
is still a subuniverse of the corresponding product of domains. The termination
is Caveat~\ref{haz:weakening} and is not proved here.
\end{remark}

\subsection{Coverings}\label{sec:coverings}

\begin{definition}[Expanded covering]\label{def:expanded}
$\mathcal I'$ is an \emph{expanded covering} of $\mathcal I$, written
$\mathcal I'\in\Expanded(\mathcal I)$, if there is
$S\colon\Var(\mathcal I')\to\Var(\mathcal I)$ with
\begin{enumerate}\romanlabels\tightlist
\item $S(x)=x$ for $x\in\Var(\mathcal I)\cap\Var(\mathcal I')$;
\item $D_{x}=D_{S(x)}$ for every $x\in\Var(\mathcal I')$;
\item for every constraint $R(x_{1},\dots,x_{m})$ of $\mathcal I'$, either
  $S(x_{1}),\dots,S(x_{m})$ are pairwise distinct and
  $R(S(x_{1}),\dots,S(x_{m}))$ is weaker or equivalent to a constraint of
  $\mathcal I$, or $S(x_{1})=\dots=S(x_{m})$ and
  $\{(a,\dots,a):a\in D_{x_{1}}\}\subseteq R$.
\end{enumerate}
It is a \emph{covering} if moreover $R(S(x_{1}),\dots,S(x_{m}))\in\mathcal I$
for every constraint; a \emph{tree-covering} if it is a covering and a
tree-instance. $S(x)$ is the \emph{parent} of $x$, and $x$ a \emph{child}.
\end{definition}

\begin{proposition}[Basic properties]\label{prop:covering-basics}
Let $\mathcal I'$ be an expanded covering of $\mathcal I$.
\begin{enumerate}\alphlabels\tightlist
\item Replacing every $x$ by $S(x)$ and deleting the constraints
  $R(x,\dots,x)$ yields a weakening of $\mathcal I$.
\item A weakening is an expanded covering with $S=\mathrm{id}$.
\item Every solution of $\mathcal I$ lifts to a solution of $\mathcal I'$.
\item If $\mathcal I$ is $1$-consistent and $\mathcal I'$ is a tree-covering,
  the solution set of $\mathcal I'$ is subdirect.
\item The union of two expanded coverings is an expanded covering.
\item An expanded covering of an expanded covering is an expanded covering.
\item An expanded covering of a cycle-consistent irreducible instance is
  cycle-consistent and irreducible.
\item Every reduction of $\mathcal I$ extends to $\mathcal I'$, and a
  $1$-consistent one extends to a $1$-consistent one.
\end{enumerate}
\end{proposition}

\begin{hazard}[Variable renaming carries obligations]\label{haz:renaming}
The constructions of \S\ref{sec:main-statements} say ``rename the variables so
that the only common variable of $\Upsilon_{x}$ and $\mathcal J$ is $x$'', and
form conjunctions $\bigwedge_{\mathcal J\in\Omega}\mathcal J$ of instances
whose variable sets must first be made disjoint. Renaming is not free: each of
$1$-consistency, cycle-consistency, cruciality and being a covering has to be
shown invariant under it. Those four facts are used throughout
\S\ref{sec:main-statements} and are nowhere stated.
\end{hazard}

\subsection{Induced congruences and connectedness}

\begin{definition}[$\ConOne$]\label{def:conone}
For $R$ of arity $m$ and $i\in[m]$, $\ConOne(R,i)$ is the binary relation
\[
  \{(y,y')\;:\;\exists\,\bar x\ \ R(\bar x[i\mapsto y])\wedge R(\bar x[i\mapsto y'])\}.
\]
For a constraint $C=R(x_{1},\dots,x_{m})$ we write $\ConOne(C,x_{i})$ for
$\ConOne(R,i)$. For an instance, $\Con(\mathcal I,x)=\{\ConOne(C,x):C\in
\mathcal I\}$ and $\Con(\mathcal I)=\bigcup_{x}\Con(\mathcal I,x)$.
\end{definition}

\begin{lemma}\label{lem:conone-basics}
The $i$-th variable of $R$ is rectangular if and only if $R$ is stable under
$\ConOne(R,i)$. If $R$ is subdirect and its $i$-th variable is rectangular,
then $\ConOne(R,i)$ is a congruence.
\end{lemma}

\begin{hazard}\label{haz:conone-not-congruence}
$\ConOne(R,i)$ is always reflexive and symmetric but is \emph{not} transitive
in general, hence not always a congruence; Lemma~\ref{lem:conone-basics} is
what makes it one, and its hypotheses (subdirect \emph{and} rectangular at $i$)
are both needed. The notation invites the reader to assume congruence-hood, and
every use below must check the hypotheses. Note also that
$\ConOne(C,x_{i})$ depends on the \emph{position} $i$, not on the variable, so
it is ill-defined if $x$ occurs twice in the scope
(Caveat~\ref{haz:instance-list}).
\end{hazard}

\begin{definition}[Types of relations and instances]\label{def:rel-type}
$R$ is \emph{of PC type} (resp.\ \emph{linear type}) if $R$ is rectangular and
every $\ConOne(R,i)$ is a PC (resp.\ linear) congruence. An instance is of PC
or linear type if all its constraints are.
\end{definition}

\begin{definition}[Adjacency, connectedness]\label{def:connected}
Two congruences $\sigma_{1},\sigma_{2}$ on $\mathbf D_{x}$ are \emph{adjacent}
if there is a reflexive bridge from $\sigma_{1}$ to $\sigma_{2}$. Two
rectangular constraints $C_{1},C_{2}$ are \emph{adjacent in a common variable
$x$} if $\ConOne(C_{1},x)$ and $\ConOne(C_{2},x)$ are adjacent. An instance
$\mathcal I$ is \emph{connected} if all its constraints are rectangular, all
congruences in $\Con(\mathcal I)$ are irreducible, and the graph on the
constraints with edges the adjacent pairs is connected.
\end{definition}

Note that every proper congruence is adjacent to itself, via
$\delta(x_{1},x_{2},x_{3},x_{4})=\sigma(x_{1},x_{3})\wedge\sigma(x_{2},x_{4})$.

\clearpage
%!TeX root=csp.tex
\section{Properties of strong subuniverses}\label{sec:properties}

This section is Zhuk's \S2.3. It is the \emph{interface} between the algebraic
theory and the CSP argument: \S\ref{sec:main-statements} uses the results
below and nothing else from \S\S\ref{sec:types}. Proofs are deferred to \S5 of \cite{ZhukSimplified}, which is not
reproduced here; \S\ref{sec:formalization} budgets it.

Its shape is worth naming. Four of the five groups say that the relation
$\lll$ and the types $T$ \emph{propagate}: forward and backward along
surjective homomorphisms, down to quotients, across products of a congruence,
and through subdirect relations. The fifth --- Theorem~\ref{thm:stable-intersection}
--- is the one genuinely hard statement, and it is what converts a failure of
consistency into a bridge. Everything the CSP proof does with bridges traces
back to it.

\begin{convention}\label{conv:type-quantifier}
Where a type is not specified, $T$ ranges over
$\{\TBA,\TC,\TS,\TPC,\TL,\TD\}$. Where a multi-type $\MT T$ appears,
$T$ ranges over $\{\TPC,\TL,\TD\}$.
\end{convention}

\subsection{Ubiquity}

\begin{lemma}[Ubiquity]\label{lem:ubiquity}
If $B\lll A$ and $|B|>1$, then $C\subt[A]{T}B$ for some $C$ and some
$T\in\{\TBA,\TC,\TL,\TPC\}$.
\end{lemma}

Lemma~\ref{lem:ubiquity} is the formal content of Informal
Claim~\ref{ic:existence}: every domain of size at least $2$ that has been
reached by a chain of strong reductions admits another strong reduction. It
is what guarantees the outer loop of the algorithm makes progress, and it is
the reason the algorithm terminates with singleton domains or a linear
instance.

\begin{hazard}\label{haz:ubiquity-hyp}
The hypothesis is $B\lll A$, not $B\le A$. A subalgebra reached by no chain
of strong reductions need admit nothing. In the algorithm this hypothesis is
maintained as an invariant of the reduction loop; in a formalization it should
be carried in the state, not re-derived.
\end{hazard}

\subsection{Propagation}

\begin{longlemma}[Propagation along a homomorphism]\label{lem:propagation}
Let $f\colon\mathbf A\to\mathbf A'$ be a surjective homomorphism. Then
\begin{enumerate}\tightlist
\item[\textup{(f)}] $C\lll^{A}B\Rightarrow f(C)\lll f(B)$;
\item[\textup{(b)}] $C'\lll^{A'}B'\Rightarrow f^{-1}(C')\lll f^{-1}(B')$;
\item[\textup{(ft)}] $C\subt[A]{T(\sigma)}B\lll A\Rightarrow
  \bigl(f(C)=f(B)$ or $f(C)\subt{\TS}f(B)$ or
  $f(C)\subt[A']{T}f(B)\bigr)$;
\item[\textup{(bt)}] $C'\subt[A']{T(\sigma)}B'\Rightarrow
  f^{-1}(C')\subt[A]{T(f^{-1}(\sigma))}f^{-1}(B')$;
\item[\textup{(fs)}] $T\in\{\TBA,\TC,\TS\}$ and $C\subt{T}B$
  $\Rightarrow f(C)\subte{T}f(B)$;
\item[\textup{(fm)}] $C\subte[A]{\MT T}B\lll A$ and $f(B)$ $\TS$-free
  $\Rightarrow f(C)\subte[A']{\MT T}f(B)$;
\item[\textup{(bm)}] $C'\subte[A']{\MT T}B'\lll A'\Rightarrow
  f^{-1}(C')\subte[A]{\MT T}f^{-1}(B')$.
\end{enumerate}
\end{longlemma}

\begin{hazard}[The three-way disjunction in (ft)]\label{haz:ft}
Clause (ft) is the characteristic shape of this theory and the reason it is
usable. Pushing a strong subset forward along a homomorphism can
\begin{itemize}\tightlist
\item collapse it ($f(C)=f(B)$: the reduction becomes vacuous),
\item degrade it ($f(C)\subt{\TS}f(B)$: still strong, but the type is lost),
\item or preserve it ($f(C)\subt[A']{T}f(B)$: same type).
\end{itemize}
The type $\TS$ exists to be the value of the middle case. Every downstream
proof that pushes a reduction through a quotient does a three-way case split
here, and the middle case is the one that is easy to forget; in \cite{ZhukJACM}
it is the source of the ``go forward and backward from global to local''
induction that this paper eliminates.

Note the asymmetry: the backward clauses (b), (bt), (bm) are clean
implications with no disjunction. Preimages behave; images do not.
\end{hazard}

\begin{corollary}[Propagation to a quotient]\label{cor:prop-quotient}
Let $\delta$ be a congruence on $\mathbf A$. Then
\textup{(f)} $C\lll^{A}B\Rightarrow C/\delta\lll^{A/\delta}B/\delta$;
\textup{(t)} $C\subt[A]{T(\sigma)}B\lll A\Rightarrow
 \bigl(C/\delta=B/\delta$ or $C/\delta\subt{\TS}B/\delta$ or
 $C/\delta\subt[A/\delta]{T}B/\delta\bigr)$;
\textup{(s)} for $T\in\{\TBA,\TC,\TS\}$, $C\subt{T}B\Rightarrow
 C/\delta\subte{T}B/\delta$;
\textup{(m)} $C\subte[A]{\MT T}B\lll A$ and $B/\delta$ $\TS$-free
 $\Rightarrow C/\delta\subte[A/\delta]{\MT T}B/\delta$.
\end{corollary}

\begin{corollary}[Reflection from a quotient]\label{cor:prop-from-factor}
Let $\delta$ be a congruence on $\mathbf A$ and $B,C\le\mathbf A$. Then
$C/\delta\lll^{A/\delta}B/\delta\iff C\circ\delta\lll^{A}B\circ\delta$, and
$C/\delta\subt[A/\delta]{T}B/\delta\iff C\circ\delta\subt[A]{T}B\circ\delta$.
\end{corollary}

\begin{corollary}[Multiplying by a congruence]\label{cor:prop-times-cong}
Let $\delta$ be a congruence on $\mathbf A$. Then
\textup{(f)} $C\lll^{A}B\Rightarrow C\circ\delta\lll^{A}B\circ\delta$;
\textup{(t)} $C\subt[A]{T(\sigma)}B\lll^{A}A\Rightarrow
 \bigl(C\circ\delta=B\circ\delta$ or $C\circ\delta\subt[A]{\TS}B\circ\delta$
 or $C\circ\delta\subt[A]{T}B\circ\delta\bigr)$;
\textup{(e)} if $\delta\subseteq\sigma$ and $C\subt[A]{T(\sigma)}B\lll^{A}A$
 then $C\circ\delta\subt[A]{T}B\circ\delta$.
\end{corollary}

\begin{longcorollary}[Propagation to relations]\label{cor:prop-relations}
Let $R\sdle\mathbf A_{1}\times\dots\times\mathbf A_{m}$ and $B_{i}\lll A_{i}$.
Then
\begin{enumerate}\tightlist
\item[\textup{(r)}] $R\cap(B_{1}\times\dots\times B_{m})\dlll R$;
\item[\textup{(r1)}] $\proj_{1}\bigl(R\cap(B_{1}\times\dots\times B_{m})\bigr)
  \dlll A_{1}$;
\item[\textup{(b)}] if $C_{i}\lll^{A_{i}}B_{i}$ for all $i$ then
  $R\cap(C_{1}\times\dots\times C_{m})\dlll^{R}R\cap(B_{1}\times\dots\times B_{m})$;
\item[\textup{(b1)}] under the same hypothesis,
  $\proj_{1}\bigl(R\cap\prod C_{i}\bigr)\dlll^{A_{1}}
   \proj_{1}\bigl(R\cap\prod B_{i}\bigr)$;
\item[\textup{(m)}] if $C_{i}\subte[A_{i}]{\MT T}B_{i}$ for all $i$ then
  $R\cap\prod C_{i}\dsubte[R]{\MT T}R\cap\prod B_{i}$;
\item[\textup{(m1)}] if in addition $\proj_{1}(R\cap\prod B_{i})$ is
  $\TS$-free, then $\proj_{1}(R\cap\prod C_{i})
   \dsubte[A_{1}]{\MT T}\proj_{1}(R\cap\prod B_{i})$.
\end{enumerate}
\end{longcorollary}

For the two absorption types the conclusion is stronger --- no chain, no dot on
the type, and no subdirectness hypothesis:

\begin{lemma}[{\cite[Cor.~6.1.2, 6.9.2]{ZhukStrong}}]\label{lem:ba-center-implies}
Let $R\le\mathbf A_{1}\times\dots\times\mathbf A_{m}$ with
$\proj_{1}(R)=A_{1}$, and let $C_{i}\subte{T}A_{i}$ for every $i$, where
$T\in\{\TBA,\TC\}$ --- and, when $T=\TBA$, with a \emph{single} binary term $t$
witnessing all $m$ absorptions. Then
$\proj_{1}\bigl(R\cap(C_{1}\times\dots\times C_{m})\bigr)\dsubte{T}A_{1}$,
again with the same $t$ when $T=\TBA$.
\end{lemma}

\begin{hazard}[The source's restatement is false; this is the repair]
\label{haz:ba-center-implies}
Lemma~\ref{lem:ba-center-implies} is the workhorse of case (2) of
Theorem~\ref{thm:main-inductive}, and \cite[Lemma~19]{ZhukSimplified} restates
it from \cite{ZhukStrong} \emph{dropping two hypotheses}: the subdirectness
$\proj_{1}(R)=A_{1}$, and --- in the $\TBA$ case --- that the $C_{i}$ absorb
with a common term. Both are present in the cited source
(\cite[Cor.~6.1.2]{ZhukStrong} reads ``Suppose $R\le A_{1}\times\dots\times
A_{n}$, $\proj_{1}(R)=A_{1}$, $B_{i}\le_{\TBA(t)}A_{i}$ for every $i$''), and
without the first the statement is false.

\emph{Counterexample to the printed form.} Take $R=\{(0,0)\}\le\mathbf Z_{p}
\times\mathbf Z_{p}$ --- a subuniverse, by idempotence --- and $C_{i}=A_{i}$,
which is allowed because $\subte{T}$ admits equality. Then
$R\cap(C_{1}\times C_{2})=R$ and $\proj_{1}(R)=\{0\}$, which is neither empty
nor all of $\mathbf Z_{p}$, so the conclusion asserts $\{0\}\subt{\TBA}
\mathbf Z_{p}$ --- contradicting Lemma~\ref{lem:no-abs-in-linear}, the paper's
own Lemma~29. The failure is exactly the missing subdirectness: $\proj_{1}(R)
=\{0\}\neq\mathbf Z_{p}$.

The repair is to restore both hypotheses, as above. A formalization must also
check that they hold at each of the call sites --- in particular in the proof of
Theorem~\ref{thm:main-inductive}(2), where the lemma is applied to the relation
computing $E_{2}$ from $E_{1}$. The common-term condition is the reason
\cite{ZhukStrong} writes $\le_{\TBA(t)}$ with the term in the notation, and
carrying $t$ in the type is the cheapest way to keep it.

This is the single highest-value import to formalize first: small,
self-contained, used everywhere, and adjacent to what the predecessor project
already proves.
\end{hazard}

\subsection{Intersections}

\begin{lemma}[Intersecting with a strong subset]\label{lem:intersect-all}
If $B\lll A$ and $D\lll A$ then \textup{(i)} $B\cap D\dlll A$, and
\textup{(t)} $C\subt[A]{T(\sigma)}B\Rightarrow C\cap D\dsubte[A]{T(\sigma)}B\cap D$.
\end{lemma}

The next theorem is the heart of \S\ref{sec:properties}. Read it as follows:
if a family of strong reductions is \emph{jointly} inconsistent but
\emph{minimally} so --- dropping any one of them restores consistency --- then
the types must agree, and in the linear case one gets bridges between the
witnessing congruences. Bridges are produced nowhere else.

\begin{longtheorem}[Stable intersection]\label{thm:stable-intersection}
Suppose
\begin{enumerate}\romanlabels\tightlist
\item $C_{i}\subt[A]{T_{i}(\sigma_{i})}B_{i}\lll A$ with
  $T_{i}\in\{\TBA,\TC,\TS,\TL,\TPC\}$, for $i\in[m]$, $m\ge 2$;
\item $\bigcap_{i\in[m]}C_{i}=\varnothing$;
\item $B_{j}\cap\bigcap_{i\neq j}C_{i}\neq\varnothing$ for every $j\in[m]$.
\end{enumerate}
Then one of:
\begin{enumerate}\tightlist
\item[\textup{(ba)}] $T_{1}=\dots=T_{m}=\TBA$;
\item[\textup{(l)}] $T_{1}=\dots=T_{m}=\TL$, and for all $k,\ell\in[m]$ there
  is a bridge $\delta$ from $\sigma_{k}$ to $\sigma_{\ell}$ with
  $\bcol{\delta}=\sigma_{k}\circ\sigma_{\ell}$;
\item[\textup{(c)}] $m=2$ and $T_{1}=T_{2}=\TC$;
\item[\textup{(pc)}] $m=2$, $T_{1}=T_{2}=\TPC$ and $\sigma_{1}=\sigma_{2}$.
\end{enumerate}
\end{longtheorem}

\begin{longcorollary}[Stable intersection, relational form]\label{cor:stable-intersection}
Suppose $R\sdle\mathbf A_{1}\times\dots\times\mathbf A_{m}$,
$C_{i}\subt[A_{i}]{T_{i}(\sigma_{i})}B_{i}\lll A_{i}$ with
$T_{i}\in\{\TBA,\TC,\TS,\TL,\TPC\}$ and $m\ge 2$,
$R\cap(C_{1}\times\dots\times C_{m})=\varnothing$, and
$R\cap(C_{1}\times\dots\times B_{j}\times\dots\times C_{m})\neq\varnothing$ for
every $j$. Then one of:
\textup{(ba)} all $T_{i}=\TBA$;
\textup{(l)} all $T_{i}=\TL$ and for all $k,\ell$ a bridge $\delta$ from
$\sigma_{k}$ to $\sigma_{\ell}$ with
$\bcol{\delta}=\sigma_{k}\circ\proj_{k,\ell}(R)\circ\sigma_{\ell}$;
\textup{(c)} $m=2$ and $T_{1}=T_{2}=\TC$;
\textup{(pc)} $m=2$, $T_{1}=T_{2}=\TPC$,
$\mathbf A_{1}/\sigma_{1}\cong\mathbf A_{2}/\sigma_{2}$, and
$\{(a/\sigma_{1},b/\sigma_{2}):(a,b)\in R\}$ is bijective.
\end{longcorollary}

\begin{remark}\label{rem:duplicate-coordinate}
Several applications need more than one restriction on a single coordinate of
$R$. The statement is kept simple; to apply it, duplicate the coordinate (form
$R'=\{(\bar a,a_{k}):\bar a\in R\}$) and restrict the copies separately.
\end{remark}

\begin{hazard}[Minimality is a hypothesis, not a construction]\label{haz:minimality}
Hypotheses (ii)--(iii) of Theorem~\ref{thm:stable-intersection} say the family
is \emph{critically} inconsistent. In applications one starts with an
inconsistent family and shrinks it to a critical subfamily. That shrinking step
is left implicit in the paper (``let $\mathcal M$ be a minimal set of children
of $z$ we need to restrict''). It is easy --- a minimal subset of a finite set
with a monotone property --- but it must be stated as a lemma, because it is
performed at least four times and each time the property being minimised is
slightly different. State once: \emph{if $P$ is a monotone predicate on subsets
of a finite set $X$ and $P(X)$ holds, there is $M\subseteq X$ minimal with
$P(M)$}, and instantiate.
\end{hazard}

\subsection{Multi-types}

\begin{lemma}\label{lem:multitype-chain}
If $C\subte[A]{\MT T}B$ then $C\subt[A]{T}\dots\subt[A]{T}B$ and
$C\lll^{A}B$.
\end{lemma}

That is: an intersection of type-$T$ subsets, though not itself obtained by a
single reduction of type $T$, is reachable by a \emph{chain} of type-$T$
reductions. This is what makes $\MT T$ compatible with $\lll$, and it is used
whenever the algorithm reduces several variables at once.

\begin{lemma}[Preserving linkedness]\label{lem:preserve-linked}
Let $R\sdle\mathbf A_{1}\times\mathbf A_{2}$,
$C_{i}\subte[A_{i}]{\TMD}B_{i}\lll A_{i}$ for $i=1,2$, let $S$ be the
rectangular closure of $R$, and suppose $R\cap(B_{1}\times B_{2})\neq
\varnothing$ and $S\cap(C_{1}\times C_{2})\neq\varnothing$. Then
$R\cap(C_{1}\times C_{2})\neq\varnothing$.
\end{lemma}

\begin{lemma}[Maximal multi-type extension]\label{lem:maximal-mult}
Let $C_{1}\subt[A]{\MT T}B_{1}\lll A$, $B_{2}\lll A$,
$C_{1}\cap B_{2}=\varnothing$, $B_{1}\cap B_{2}\neq\varnothing$, and let
$\sigma$ be a maximal congruence on $\mathbf A$ with
$(C_{1}\circ\sigma)\cap B_{2}=\varnothing$. Then
$\sigma=\omega_{1}\cap\dots\cap\omega_{s}$ where each $\omega_{i}$ is a
congruence of type $T$ on $\mathbf A$ with $\cover{\omega_{i}}\supseteq
B_{1}^{2}$.
\end{lemma}

\subsection{Auxiliary imports}

The following are used but not proved in \cite{ZhukSimplified}.

\begin{lemma}[{\cite[Lem.~4.7]{MarotiMcKenzie}}]\label{lem:special-wnu}
If $w$ is an idempotent WNU operation of arity $n$ on a finite set $A$, then
$\Clo(w)$ contains a special idempotent WNU operation of arity $n^{n!}$.
\end{lemma}

\begin{lemma}[{\cite[Cor.~8.17.1]{ZhukJACM}}]\label{lem:building-perfect}
If $\sigma$ is an irreducible congruence on $\mathbf A\in\Vn$ and $\delta$ is a
bridge from $\sigma$ to $\sigma$ with $\bcol{\delta}=A^{2}$, then $\sigma$ is a
perfect linear congruence.
\end{lemma}

\begin{lemma}\label{lem:perfect-from-linked}
If $\sigma$ is an irreducible congruence on $\mathbf A\in\Vn$ and $\delta$ is a
bridge from $\sigma$ to $\sigma$ with $\bcol{\delta}$ linked, then $\sigma$ is
a perfect linear congruence.
\end{lemma}

\begin{proof}
Let $\delta^{-1}(y_{1},y_{2},x_{1},x_{2})=\delta(x_{1},x_{2},y_{1},y_{2})$,
which is again a bridge from $\sigma$ to $\sigma$. Since $\bcol{\delta}$ is
linked, $(\bcol{\delta}\circ\bcol{\delta^{-1}})^{N}=A^{2}$ for $N$ large.
Composing $2N$ bridges $\delta,\delta^{-1},\dots$ by
Lemma~\ref{lem:bridge-comp} gives a bridge $\delta'$ from $\sigma$ to $\sigma$
with $\bcol{\delta'}=A^{2}$; apply Lemma~\ref{lem:building-perfect}.
\end{proof}

\begin{hazard}[``for sufficiently large $N$'']\label{haz:linked-power}
The step $(\bcol{\delta}\circ\bcol{\delta^{-1}})^{N}=A^{2}$ deserves a lemma of
its own: if a reflexive symmetric relation $\rho$ on a finite set $A$ has
connected graph, then $\rho^{N}=A^{2}$ for $N\ge|A|$. Reflexivity is what makes
the powers increase monotonically; $\bcol{\delta}$ is reflexive because
$(a,a,a,a)\in\delta$ is not assumed --- so the reflexivity must come from
$\delta$ being a bridge from $\sigma$ to $\sigma$ with $\sigma$ reflexive, via
$\proj_{1,2}(\delta)\supsetneq\sigma$. Check this: it is the sort of step where
the informal proof is right for a reason it does not give.
\end{hazard}

\begin{lemma}\label{lem:no-abs-in-linear}
If $B\le\mathbf Z_{p}\times\dots\times\mathbf Z_{p}$ then there is no
$C\subt{T}B$ with $T\in\{\TBA,\TC\}$.
\end{lemma}

\begin{proof}
It suffices to check that for every term $\tau$, if the last variable of
$\tau^{\mathbf Z_{p}}(a,\dots,a,x)$ is not dummy then it takes all values;
hence no absorbing subuniverse exists.
\end{proof}


%!TeX root=csp-proof.tex
\section{Proofs of the strong-subalgebra theory}\label{sec:strongproofs}

This section proves the statements of \S\ref{sec:properties}. Throughout,
Convention~\ref{conv:standing} is in force --- every algebra is finite,
idempotent, and has a WNU term operation --- and \S\ref{sec:legislation} fixes
the readings.

\subsection{What is imported}\label{sec:imports}

The development is self-contained except for the following, which are proved in
the literature and used here as black boxes. Two are substantial; the rest are
routine. \S\ref{sec:imports-ledger} lists them again with their exact
hypotheses.

\begin{longtheorem}[Absorption Theorem]\label{thm:brady-absorption}
Let $R\sdle\mathbf A\times\mathbf B$ be a subdirect relation between finite
idempotent Taylor algebras. If $R$ is linked then one of the following holds:
$R=A\times B$; or $\mathbf A$ has a proper nonempty binary absorbing or
centrally absorbing subalgebra; or $\mathbf B$ has a proper nonempty subalgebra
that is both binary absorbing and centrally absorbing.
\end{longtheorem}

\begin{corollary}[{a corollary of Theorem~\ref{thm:brady-absorption}}]
\label{lem:linked-implies-abs}
Let $R\lneq_{sd}\mathbf A\times\mathbf B$ be a \emph{proper} subdirect
relation between finite idempotent Taylor algebras. If $R$ is linked then
$\mathbf A$ or $\mathbf B$ has a proper nonempty binary absorbing or central
subuniverse.
\end{corollary}

\begin{proof}
Properness excludes the first alternative of
Theorem~\ref{thm:brady-absorption}. The second is the conclusion on
$\mathbf A$; the third is stronger than the conclusion on $\mathbf B$.
\end{proof}

\begin{hazard}[Both words in ``proper nonempty'' are load-bearing]
\label{haz:linked-abs-proper}
Without them Corollary~\ref{lem:linked-implies-abs} is vacuous, since every
algebra is a binary absorbing subuniverse of itself and $\varnothing$ is
vacuously one (\S\ref{sec:legislation}, item \ref{leg:empty}); every use below
is the proper nonempty one. Theorem~\ref{thm:brady-absorption} is the single
heaviest import of this section, and it is stated here in the form its source
gives rather than in the weaker symmetric form actually used, so that the
strengthening is visible.
\end{hazard}

\begin{lemma}[{\cite[Lem.~2.9]{DecidingAbsorption}}, {\cite[Lem.~6.1, Thm.~6.9]{ZhukStrong}}]\label{lem:pp-propagation}
Let $R\le A^{n}$ be defined by a pp-formula $\Phi$ in which a relation $S$
occurs, and let $R'$ be defined by the formula obtained from $\Phi$ by
replacing \emph{every} occurrence of $S$ by some $S'\subt{T}S$, where
$T\in\{\TBA,\TC\}$. Then $R'\dsubte{T}R$.
\end{lemma}

\begin{hazard}[The conclusion is dotted]\label{haz:pp-dotted}
Replacing a relation by a proper strong subrelation can make the pp-defined
result \emph{empty}, so the undotted $R'\subte{T}R$ --- which under
Convention~\ref{conv:empty} entails $R'\neq\varnothing$ --- is not what the
lemma gives.

Every use below has been checked against this. Lemma~\ref{lem:ba-center-intersect}
states its conclusion dotted and passes it on. Lemmas~\ref{lem:absorb-linked-sub},
\ref{lem:block-good-bridge} and \ref{lem:pc-inductive-step} exhibit a point of
the smaller relation and so upgrade to the undotted form.
Lemma~\ref{lem:bacon-left} needs both nonemptiness \emph{and} properness, and
gets each from a different clause of Lemma~\ref{lem:barto-kazda}
(Caveat~\ref{haz:bacon-left}). The remaining use is inside
Lemma~\ref{lem:center-pushed-in}, which is an outline
(\S\ref{sec:status}).
\end{hazard}

\begin{lemma}[{\cite[Prop.~2.14]{DecidingAbsorption}}, {\cite[Lem.~3.2]{ZhukStrong}}]\label{lem:barto-kazda}
For $B\le\mathbf A$ and $n\ge2$: $B$ absorbs $\mathbf A$ with an operation of
arity $n$ if and only if there is no $S\le A^{n}$ with $S\cap B^{n}=\varnothing$
and $S\cap(B^{i-1}\times A\times B^{n-i})\neq\varnothing$ for every $i$.
\end{lemma}

\begin{lemma}[{\cite[Lem.~6.8]{ZhukStrong}}]\label{lem:factor-type}
If $B\subte{T}\mathbf A$ with $T\in\{\TBA,\TC,\TS\}$ and $\sigma$ is a
congruence on $\mathbf A$, then $B/\sigma\subte{T}\mathbf A/\sigma$.
\end{lemma}

\begin{hazard}[Only the $\TC$ case is imported]\label{haz:factor-type}
For $T=\TBA$ the statement is one line: if $t(B,A)\subseteq B$ and
$t(A,B)\subseteq B$ then the same $t$ works on images. The case $T=\TS$ follows
from the other two by Definition~\ref{def:ba-center-free}. Only $\TC$ is
genuinely imported.
\end{hazard}

\begin{lemma}[{\cite[Lem.~6.24]{ZhukStrong}}]\label{lem:power-to-base}
If $R$ is a nontrivial strong subuniverse of
$\mathbf A_{1}\times\dots\times\mathbf A_{n}$ of a type $T\neq\TPC$, then some
$\mathbf A_{i}$ has a nontrivial subuniverse of type $T$. In particular
$B\subt{T}\mathbf A^{n}$ with $T\in\{\TBA,\TC,\TS\}$ yields some
$C\subt{T}\mathbf A$.
\end{lemma}

\begin{hazard}[One run of the proof serves all three types]\label{haz:power-to-base}
The cited lemma is stated for every type $T\neq\TPC$, and its proof is
\emph{type-uniform}: it takes $\proj_{1}(R)$ if that is proper and otherwise a
fibre $R'=\{(a_{2},\dots,a_{n}):(a_{1},\dots,a_{n})\in R\}$, and neither
construction mentions the type --- the type enters only through the facts that
projections and fibres preserve $\TBA$ and preserve $\TC$. So a single run
produces one witness carrying \emph{every} strong type that $R$ has; in
particular a subuniverse that is simultaneously binary absorbing and central
yields one that is too, which is the case $T=\TS$.
\end{hazard}

\begin{longtheorem}[{\cite[Thm.~6.15]{ZhukStrong}}]\label{thm:central-relation-source}
Let $R\sdle\mathbf A\times\mathbf B$ and
$C=\{c\in A:\{c\}\times B\subseteq R\}$. Then $C$ is a central subuniverse of
$\mathbf A$, or $\mathbf B$ has a nontrivial binary absorbing subuniverse, or
$\mathbf B$ has a nontrivial \emph{projective} subuniverse.
\end{longtheorem}

\begin{hazard}[The heaviest import after Theorem~\ref{thm:brady-absorption}]
\label{haz:615-todo}
Theorem~\ref{thm:central-relation-source} is taken on faith here. Its proof in
\cite{ZhukStrong} is two pages and uses only the Barto--Kazda criterion
(Lemma~\ref{lem:barto-kazda}) and the theory of projective subuniverses; none
of that theory is needed anywhere else below, which is why absorbing it has
been deferred. The same applies to \cite[Lem.~3.4]{ZhukStrong}, used once, in
the proof of Lemma~\ref{lem:central-relation}.
\end{hazard}

\begin{lemma}[{\cite{HobbyMcKenzie}}]\label{lem:hobby-mckenzie}
An algebra is Abelian if and only if some congruence of its square has the
diagonal as a block; and a finite algebra with a WNU term operation is Abelian
if and only if it is affine.
\end{lemma}

\begin{lemma}\label{lem:bridge-abelian}
Let $|A|>1$. Then $\mathbf A$ is Abelian if and only if there is a bridge
$\delta$ from $0_{\mathbf A}$ to $0_{\mathbf A}$ with
$\bcol\delta=\proj_{1,2}(\delta)=\proj_{3,4}(\delta)=A^{2}$.
\end{lemma}

\begin{proof}
If $\mathbf A$ is Abelian, a congruence of $\mathbf A^{2}$ having the diagonal
as a block is such a bridge; clause (c) needs $|A|>1$, since for $|A|=1$ no
relation properly contains $0_{\mathbf A}=A^{2}$.

Conversely, regard $\delta$ as a binary relation $\hat\delta$ on $A^{2}$,
$\hat\delta\bigl((x_{1},x_{2}),(x_{3},x_{4})\bigr)\iff\delta(x_{1},x_{2},
x_{3},x_{4})$, and put $\rho=\hat\delta\circ\hat\delta^{-1}$, a reflexive
symmetric compatible relation on $\mathbf A^{2}$ --- reflexive because
$\proj_{1,2}(\delta)=A^{2}$ makes $\hat\delta$ left-total. Its powers increase,
so $\theta=\rho^{N}$ for $N\ge|A|^{2}$ is reflexive, symmetric and transitive,
hence a congruence of $\mathbf A^{2}$. Clause (d) of
Definition~\ref{def:bridge} is preserved by each composition, so every pair of
$\theta$ relates a diagonal element only to diagonal elements: the diagonal is
a $\theta$-block. Now apply Lemma~\ref{lem:hobby-mckenzie}.
\end{proof}

\subsection{Absorption preliminaries}\label{sec:abs-prelim}

\begin{lemma}\label{lem:ba-center-intersect}
If $B\subte{T}\mathbf A$ with $T\in\{\TBA,\TC\}$ and $C\le\mathbf A$, then
$B\cap C\dsubte{T}C$.
\end{lemma}

\begin{proof}
$C$ is defined by the pp-formula $A(x)\wedge C(x)$; replacing the occurrence of
the full unary relation $A$ by $B$ yields $B\cap C$, so
Lemma~\ref{lem:pp-propagation} applies.
\end{proof}

\begin{lemma}[No absorbing diagonal]\label{lem:no-absorbing-diagonal}
Let $\mathbf C$ be a finite idempotent algebra with $|C|\ge 2$. Then
$0_{\mathbf C}$ is not an absorbing subuniverse of $\mathbf C^{2}$, of any
arity; in particular $0_{\mathbf C}\not\subte{\TBA}\mathbf C^{2}$ and
$0_{\mathbf C}\not\subte{\TC}\mathbf C^{2}$.
\end{lemma}

\begin{proof}
Suppose $t$ is $k$-ary and absorbs at every position. Fix $i$, take
$a_{j}\in C$ for $j\neq i$ and $b,c\in C$, and apply $t$ coordinatewise to the
tuples $(a_{j},a_{j})\in 0_{\mathbf C}$ for $j\neq i$ and $(b,c)\in C^{2}$ at
position $i$. The result
$\bigl(t(\bar a;b),\,t(\bar a;c)\bigr)$ lies in $0_{\mathbf C}$, so
$t(\bar a;b)=t(\bar a;c)$. As $b,c$ were arbitrary, $t$ does not depend on its
$i$-th argument; and this holds for every $i$, so $t$ is constant. Idempotence
then forces $|C|=1$. A central subuniverse is absorbing by
Definition~\ref{def:central}, so the central case is included.
\end{proof}

\begin{lemma}[Absorption preserves linkedness, subrelation form]\label{lem:absorb-linked-sub}
Let $W\sdle\mathbf P\times\mathbf A$ be linked and let
$\varnothing\neq W'\subte{T}W$ with $T\in\{\TBA,\TC\}$. Put
$P'=\proj_{1}(W')$ and $A'=\proj_{2}(W')$. Then $W'\sdle\mathbf P'\times
\mathbf A'$ is linked.
\end{lemma}

\begin{proof}
Let $\Phi_{k}(x,y)$ be the pp-formula in one binary relation symbol obtained by
iterating $x\sim y\iff\exists z\,(W(x,z)\wedge W(y,z))$ $k$ times. Because $W$
is subdirect, $W\circ W^{-1}$ is reflexive and symmetric on $P$, so for
$k\ge|P|$ the formula $\Phi_{k}$ interpreted in $W$ defines the transitive
closure of $W\circ W^{-1}$, which is $P^{2}$ since $W$ is linked. Interpreted
in $W'$ and for the same $k\ge|P|\ge|P'|$ it defines
$\lambda:=\LeftLinked(W')$. By Lemma~\ref{lem:pp-propagation},
replacing every occurrence of $W$ by $W'$ gives $\lambda\dsubte{T}P^{2}$, and
$\lambda\neq\varnothing$ because $W'\neq\varnothing$ makes $P'$ nonempty and
$\lambda$ reflexive on it, so $\lambda\subte{T}P^{2}$. Since
$\lambda\subseteq (P')^{2}\le P^{2}$, Lemma~\ref{lem:ba-center-intersect} gives
$\lambda\subte{T}(P')^{2}$.

Now $\lambda$ is a congruence on $\mathbf P'$: reflexive because $W'$ is
subdirect over $P'$, symmetric by construction, transitive because the powers
stabilise. So $\lambda\times\lambda$ is a congruence on the algebra $(P')^{2}$,
and the image of $\lambda$ under $(P')^{2}\twoheadrightarrow(P'/\lambda)^{2}$ is
exactly $0_{\mathbf P'/\lambda}$; by Lemma~\ref{lem:factor-type},
$0_{\mathbf P'/\lambda}\subte{T}(\mathbf P'/\lambda)^{2}$. Lemma
\ref{lem:no-absorbing-diagonal} forces $|P'/\lambda|=1$, i.e.\
$\lambda=(P')^{2}$, which says $W'$ is linked.
\end{proof}

\begin{hazard}[The box form is not enough]\label{haz:box-vs-sub}
The familiar form of this principle, \cite[Prop.~2.15(i)]{BartoKozik},
restricts by a \emph{box}: if $W\sdle\mathbf A_{1}\times\mathbf A_{2}$ is
linked, $B_{i}$ absorbs $\mathbf A_{i}$, and $W\cap(B_{1}\times B_{2})$ is
subdirect in $B_{1}\times B_{2}$, then $W\cap(B_{1}\times B_{2})$ is linked.
That form cannot reach the restrictions used in
Lemma~\ref{lem:pc-inductive-step}, and the two really do differ: for the
meet-semilattice on $\{0,1,2\}$ with $B=\{0,1\}$, the subalgebra
$K=\{(0,0),(1,2)\}$ of $A^{2}$ meets $B^{2}$ and has
$\proj_{1}(K)\cap B=\{0,1\}$ while $\proj_{1}(K\cap B^{2})=\{0\}$. The
subrelation form also discharges the box form's subdirectness hypothesis for
free.
\end{hazard}

\begin{lemma}[A WNU forbids essentially unary quotients]\label{lem:no-unary-quotient}
Let $\mathbf U$ be a finite algebra with a WNU term operation and $|U|\ge2$.
Then some term operation of $\mathbf U$ depends on more than one variable.
Consequently, if $\mathbf A$ has a WNU term operation then no
$\mathbf U\in\mathrm{HS}(\mathbf A)$ with $|U|\ge2$ is essentially unary.
\end{lemma}

\begin{proof}
Suppose every term operation of $\mathbf U$ has at most one non-dummy variable,
and let $w$ be a WNU term operation of arity $k\ge2$; it is idempotent. If $w$
has no non-dummy coordinate it is constant, and idempotence gives $|U|=1$.
Otherwise $w(\bar x)=g(x_{i})$ for a unary $g$, and $g(x)=w(x,\dots,x)=x$, so
$w$ is the $i$-th projection. The WNU identities equate the $k$ terms in which
a single $y$ sits at each position against a background of $x$'s; two of them
are $w$ with $y$ at position $i$, which evaluates to $y$, and $w$ with $y$ at
some position $\neq i$, which evaluates to $x$. Hence $x=y$ for all $x,y$ and
$|U|=1$. The consequence follows because idempotence and the WNU identities are
identities, hence hold in every member of $\mathrm{HS}(\mathbf A)$.
\end{proof}

\begin{longlemma}[Central relations]\label{lem:central-relation}
Let $\mathbf A,\mathbf B$ satisfy Convention~\ref{conv:standing},
let $R\sdle\mathbf A\times\mathbf B$, and put
$C=\{c\in A:\{c\}\times B\subseteq R\}$. Then
\begin{enumerate}\alphlabels\tightlist
\item $C=\varnothing$, or
\item $C$ is a central subuniverse of $\mathbf A$, or
\item $\mathbf B$ has a proper nonempty binary absorbing subuniverse.
\end{enumerate}
\end{longlemma}

\begin{proof}
Apply Theorem~\ref{thm:central-relation-source}. Its first alternative gives
(a) or (b) according as $C$ is empty or not, its second gives (c). In its third,
$\mathbf B$ has a nontrivial projective subuniverse $P$; by
\cite[Lem.~3.4]{ZhukStrong}, either $P$ is a binary absorbing subuniverse of
$\mathbf B$, which is (c) since $P$ is nontrivial, or there is an essentially
unary $\mathbf U\in\mathrm{HS}(\mathbf B)$ with $|U|\ge2$, which
Lemma~\ref{lem:no-unary-quotient} forbids.
\end{proof}

\begin{hazard}[Alternative (a) is not optional]\label{haz:central-relation}
Under Convention~\ref{conv:empty} the two-alternative form (b)--(c) is false:
$\mathbf A=\mathbf B=\mathbf Z_{p}$ with $R=\{(x,x+1)\}$ is subdirect, has
$C=\varnothing$, and $\mathbf Z_{p}$ has no nontrivial binary absorbing
subuniverse at all. It is true in \cite{ZhukStrong} only because $\varnothing$
counts as central there. Every call site below discharges (a) by exhibiting a
point of $C$ first, and two further obligations go with each contrapositive
use: the centre produced must be shown \emph{proper}, and a strong subuniverse
of a \emph{power} must be moved down by Lemma~\ref{lem:power-to-base}.
\end{hazard}

\begin{lemma}\label{lem:bacon-left}
Let $R\sdle\mathbf A\times\mathbf B$ with $\mathbf A$ BA and centre free,
$\LeftLinked(R)=A^{2}$, and
$C=\{c\in B:A\times\{c\}\subseteq R\}$. Then $C\neq\varnothing$ and
$C\subte{\TBA,\TC}\mathbf B$.
\end{lemma}

\begin{proof}
If $|A|=1$ then subdirectness gives $A\times\{b\}\subseteq R$ for every
$b\in B$, so $C=B$ and the conclusion holds with equality. Assume $|A|\ge2$.

\emph{Step 1: $C\neq\varnothing$.} For $k\ge1$ put
$W_{k}=\{(a_{1},\dots,a_{k}):\exists b\,\forall i\,R(a_{i},b)\}$, a
subuniverse of $\mathbf A^{k}$, with $W_{1}=\proj_{1}(R)=A$. If
$W_{|A|}=A^{|A|}$ then a witness $b$ for the tuple listing all of $A$ has
$A\times\{b\}\subseteq R$, so $b\in C$ and we are done. Otherwise choose
$k\ge2$ minimal with $W_{k}\neq A^{k}$, so $W_{k-1}=A^{k-1}$.

View $W_{k}$ as a binary relation in $\mathbf A\times\mathbf A^{k-1}$. It is
\emph{subdirect}: for $a\in A$, subdirectness of $R$ gives $b$ with $R(a,b)$
and hence $(a,a,\dots,a)\in W_{k}$, so the left projection is $A$; and for
$\bar a\in A^{k-1}$, $W_{k-1}=A^{k-1}$ gives $b$ with $R(a_{i},b)$ for all
$i$, while $\proj_{2}(R)=B$ gives some $a_{0}$ with $R(a_{0},b)$, so
$(a_{0},\bar a)\in W_{k}$ and the right projection is $A^{k-1}$.

It is \emph{linked}. If $a,c\in A$ share an $R$-neighbour $b$ --- that is,
$R(a,b)$ and $R(c,b)$ --- then $b$ also witnesses
$(a,a,\dots,a),(c,a,\dots,a)\in W_{k}$, so $(a,c)\in\LeftLinked(W_{k})$ by
Definition~\ref{def:linked}. Since $\LeftLinked(R)$ is the least equivalence
closed under that step and $\LeftLinked(W_{k})$ is an equivalence containing
each such pair, $\LeftLinked(R)\subseteq\LeftLinked(W_{k})$; and
$\LeftLinked(R)=A^{2}$ by hypothesis, so $\LeftLinked(W_{k})=A^{2}$.

So $W_{k}$ is a proper linked subdirect relation, and
Corollary~\ref{lem:linked-implies-abs} gives a proper nonempty binary absorbing
or central subuniverse of $\mathbf A$ or of $\mathbf A^{k-1}$;
Lemma~\ref{lem:power-to-base} moves the second case down to $\mathbf A$. Either
way this contradicts the hypothesis on $\mathbf A$. Hence
$W_{|A|}=A^{|A|}$ after all, and $C\neq\varnothing$.

\emph{Step 2: $C$ is central in $\mathbf B$.} Apply
Lemma~\ref{lem:central-relation} to $R^{-1}\sdle\mathbf B\times\mathbf A$,
whose associated set is exactly $C$. Alternative (a) is excluded by Step 1 and
alternative (c) --- a proper nonempty binary absorbing subuniverse of
$\mathbf A$ --- by hypothesis, so (b) holds.

\emph{Step 3: $C$ is binary absorbing in $\mathbf B$.} Suppose not. By
Lemma~\ref{lem:barto-kazda} at $n=2$ there is $S\le B\times B$ with
\[
  S\cap(C\times C)=\varnothing,\qquad
  S\cap(C\times B)\neq\varnothing,\qquad
  S\cap(B\times C)\neq\varnothing .
\]
Write $U_{X}=\{b\in B:\exists c\,(S(b,c)\wedge X(c))\}$ for a unary $X$ on
$B$, and consider the pp-formula
\[
  \Phi_{X}(a_{1},\dots,a_{|A|})=\exists b\;\Bigl(U_{X}(b)\wedge
    \bigwedge_{i}R(a_{i},b)\Bigr),
\]
in which the unary relation $X$ occurs once. Put $W=\Phi_{C}$ and
$W^{\circ}=\Phi_{B}$. Since $C\subt{\TC}\mathbf B$ by Step 2 ---
\emph{proper}, because $C=B$ would put $S\cap(C\times C)=S\neq\varnothing$ ---
Lemma~\ref{lem:pp-propagation} gives $W\dsubte{\TC}W^{\circ}$. Three
observations turn that into what is wanted, and each spends one of the three
displayed conditions.

\emph{$W^{\circ}=A^{|A|}$.} $S\cap(C\times B)\neq\varnothing$ supplies
$b_{0}\in C$ with $b_{0}\in U_{B}$. Since $b_{0}\in C$ we have
$A\times\{b_{0}\}\subseteq R$, so $b_{0}$ witnesses $\Phi_{B}(\bar a)$ for
\emph{every} $\bar a\in A^{|A|}$.

\emph{$W\neq\varnothing$.} $S\cap(B\times C)\neq\varnothing$ supplies
$(b_{1},c_{1})\in S$ with $c_{1}\in C$, so $b_{1}\in U_{C}$; and
$\proj_{2}(R)=B$ supplies $a$ with $R(a,b_{1})$, so
$(a,\dots,a)\in W$.

\emph{$W\neq A^{|A|}$.} If it were, the tuple enumerating $A$ would give
$b\in U_{C}$ with $A\times\{b\}\subseteq R$, i.e.\ $b\in C$; and $b\in U_{C}$
means $S(b,c)$ for some $c\in C$, so $(b,c)\in S\cap(C\times C)=\varnothing$.

So $W\subt{\TC}A^{|A|}$ is proper and nonempty, and
Lemma~\ref{lem:power-to-base} produces a proper nonempty central subuniverse of
$\mathbf A$ --- again a contradiction.
\end{proof}

\begin{hazard}[All three Barto--Kazda conditions are spent]\label{haz:bacon-left}
Step 3 is the only place in this section where the three side conditions of
Lemma~\ref{lem:barto-kazda} are used separately, one apiece: $S$ meeting
$C\times B$ is what makes the \emph{ambient} of the pp-propagation step the
full power $A^{|A|}$ rather than some subrelation of it; $S$ meeting
$B\times C$ is what makes $W$ nonempty; and $S\cap(C\times C)=\varnothing$ is
what makes it proper. Dropping any one leaves a conclusion that
Lemma~\ref{lem:power-to-base} cannot consume, since $\subt{\TC}$ requires a
proper nonempty subuniverse of a product.
\end{hazard}

\begin{lemma}\label{lem:strong-nonempty}
If $C\lll^{A}B\lll A$ and $D\subt{\TBA,\TC}B$ then $C\cap D\neq\varnothing$.
\end{lemma}

\begin{proof}
Suppose not, and choose $C''$ minimal with
$C\lll^{A}C'\subt[A]{T(\sigma)}C''\lll^{A}B$ and $C''\cap D\neq\varnothing$.
By Lemma~\ref{lem:ba-center-intersect} $C''\cap D\subt{\TBA,\TC}C''$, so by
Lemma~\ref{lem:possible-intersections} $T=\TD$; and then
Lemma~\ref{lem:factor-type} gives
$(C''\cap D)/\sigma\subt{\TBA}C''/\sigma$, contradicting the requirement in
Definition~\ref{def:types} that $C''/\sigma$ be BA and centre free.
\end{proof}

\begin{lemma}[No critical central family of size $\ge3$]\label{lem:no-central-critical}
Let $n\ge3$, let $B_{1},\dots,B_{n}\le\mathbf A$, and let
$C_{i}\subt{\TC}B_{i}$ for $i\in[n]$. Then
$\bigcap_{i}C_{i}=\varnothing$ and $B_{j}\cap\bigcap_{i\neq j}C_{i}\neq
\varnothing$ for every $j$ cannot both hold.
\end{lemma}

\begin{proof}
Suppose they do. Put $D=\bigcap_{i}B_{i}$ and $E_{i}=C_{i}\cap D$. Then $D$ is
a nonempty subuniverse, containing $B_{1}\cap C_{2}\cap\dots\cap C_{n}$; each
$E_{i}\subte{\TC}D$ by Lemma~\ref{lem:ba-center-intersect}, and properly, since
$E_{j}=D$ would give $\bigcap_{i}E_{i}=\bigcap_{i\neq j}E_{i}\neq\varnothing$.
Moreover $\bigcap_{i}E_{i}=\varnothing$ and, because $C_{i}\subseteq B_{i}$,
\[
  \textstyle\bigcap_{i\neq j}E_{i}=D\cap\bigcap_{i\neq j}C_{i}
  =B_{j}\cap\bigcap_{i\neq j}C_{i}\neq\varnothing .
\]
Now the diagonal $R=\{(a,\dots,a):a\in D\}\le D^{n}$ is a subuniverse by
idempotence, and the two displays say exactly that $R$ is
$(E_{1},\dots,E_{n})$-essential. By \cite[Lem.~6.11]{ZhukStrong} no such
relation exists for $n\ge3$ when the $E_{i}$ are proper nonempty central
subuniverses.
\end{proof}

\begin{corollary}[Relational form]\label{cor:no-central-critical}
Let $R\sdle\mathbf A_{1}\times\dots\times\mathbf A_{n}$ with $n\ge3$,
$B_{i}\le\mathbf A_{i}$, and $C_{i}\subt{\TC}B_{i}$ for $i\in[n]$. If
$R\cap(C_{1}\times\dots\times B_{j}\times\dots\times C_{n})\neq\varnothing$
for every $j$, then $R\cap(C_{1}\times\dots\times C_{n})\neq\varnothing$.
\end{corollary}

\begin{proof}
Write $f_{i}\colon\mathbf R\to\mathbf A_{i}$ for the projections and
$P_{i}=f_{i}^{-1}(B_{i})$, $Q_{i}=f_{i}^{-1}(C_{i})$, subuniverses of
$\mathbf R$. The formula $R(x_{1},\dots,x_{n})\wedge X(x_{i})$ is positive
primitive in $X$ and defines $P_{i}$ at $X=B_{i}$ and $Q_{i}$ at $X=C_{i}$, so
Lemma~\ref{lem:pp-propagation} gives $Q_{i}\dsubte{\TC}P_{i}$; and $Q_{i}$ is
nonempty, containing $R\cap(C_{1}\times\dots\times B_{j}\times\dots\times
C_{n})$ for any $j\neq i$, which exists since $n\ge3$. So
$Q_{i}\subte{\TC}P_{i}$.

If $Q_{i}=P_{i}$ for some $i$, the conclusion is immediate:
$R\cap(C_{1}\times\dots\times C_{n})
=Q_{i}\cap\bigcap_{j\neq i}Q_{j}
=P_{i}\cap\bigcap_{j\neq i}Q_{j}
=R\cap(C_{1}\times\dots\times B_{i}\times\dots\times C_{n})\neq\varnothing$.
Otherwise every $Q_{i}\subt{\TC}P_{i}$ is proper, and
Lemma~\ref{lem:no-central-critical} applies inside $\mathbf R$ with
$B_{i}:=P_{i}$ and $C_{i}:=Q_{i}$: its two forbidden conditions are exactly
$R\cap(C_{1}\times\dots\times C_{n})=\varnothing$ and the displayed
nonemptiness, read through the $f_{i}$. So the first fails.
\end{proof}

\begin{hazard}[Why this is a separate statement]\label{haz:central-critical-split}
Lemma~\ref{lem:no-central-critical} is the type-$\TC$ half of case (c) of
Theorem~\ref{thm:stable-intersection}, and it is stated on its own because
Lemma~\ref{lem:block-good-bridge} needs \emph{only} this half. Citing the whole
theorem there would make the bridge theory rest on the bridge classification,
which rests back on it:
\[
  \text{\ref{cor:stable-intersection}}\to
  \text{\ref{thm:stable-intersection}}\to
  \text{\ref{lem:bridge-to-pc}}\to\text{\ref{lem:no-cross-bridge}}\to
  \text{\ref{lem:linear-equiv}}\to\text{\ref{lem:nontrivial-bridge}}\to
  \text{\ref{lem:block-good-bridge}}\to\text{\ref{cor:stable-intersection}} .
\]
Every edge of that is a real citation, and the loop is broken only by observing
that the appeal at the last step uses no bridge at all: it is $n=3$ with all
types $\TC$, and the argument for that case runs through
\cite[Lem.~6.11]{ZhukStrong} and nothing else in this section. Splitting the
statement is what makes the observation checkable rather than a matter of
trust.

A second observation, found by the layer check rather than by reading: the
$\lll$ hypotheses these two statements used to carry are consumed by
\emph{nothing} in either proof --- nonemptiness of $\bigcap_{i}B_{i}$ comes
from the displayed conditions, and the essential-relations import needs only
proper nonempty central subuniverses. They are now stated for plain
subalgebras, and the corollary pulls back through
Lemma~\ref{lem:pp-propagation} instead of the $\lll$-pullback of
Lemma~\ref{lem:reverse-hom}, with the equality case of the pullback giving the
conclusion outright. That is what places the entire bridge classification
below the strong-subuniverse theory rather than beside it.
\end{hazard}

\begin{lemma}[{\cite[Lem.~6.25]{ZhukStrong}}]\label{lem:possible-intersections}
If $B\subt{T_{1}}\mathbf A$, $C\subt{T_{2}}\mathbf A$, $B\cap C=\varnothing$
and $T_{1},T_{2}\in\{\TBA,\TC,\TS\}$, then $T_{1}=T_{2}\in\{\TBA,\TC\}$.
\end{lemma}

\begin{lemma}[A BA-and-central subuniverse meets every strong one]\label{lem:ba-central-meets}
If $D\subt{\TBA,\TC}\mathbf A$ and $E\subt{T}\mathbf A$ with
$T\in\{\TBA,\TC,\TS\}$, then $D\cap E\neq\varnothing$.
\end{lemma}

\begin{proof}
Suppose $D\cap E=\varnothing$. Lemma~\ref{lem:possible-intersections} applied
to the pair $(D,E)$ with $T_{1}=\TBA$ gives $T=\TBA$; applied to the same pair
with $T_{1}=\TC$ it gives $T=\TC$. The two conclusions are incompatible.
\end{proof}

\begin{hazard}[Where the conjunction reading of $\subt{\TBA,\TC}$ pays]\label{haz:ba-central-meets}
Lemma~\ref{lem:ba-central-meets} is false with $\subt{\TBA,\TC}$ read as a
disjunction: two disjoint binary absorbing subuniverses are exactly what
Lemma~\ref{lem:possible-intersections} permits. What makes the argument run is
that $D$ carries \emph{both} types at once, so whichever of $\TBA$ and $\TC$
the subuniverse $E$ has, the other is available for $D$ and the two types
differ. This is the reading fixed in Caveat~\ref{haz:S-type}; it is the
mechanism already used inside the proof of Lemma~\ref{lem:strong-nonempty},
and the two appeals in Lemma~\ref{lem:intersection-good}(s) are what make it
worth naming.
\end{hazard}

\subsection{The intersection property}\label{sec:intersection-property}

The goal is Lemma~\ref{lem:intersection-good}: if $B\lll A$, $C\lll A$ meet and
$\delta$ is a dividing congruence for $B$, then $(B\cap C)/\delta$ is a single
point or all of $B/\delta$. Everything in \S\ref{sec:strongproofs} that
produces a bridge goes through it.

\begin{lemma}[The shift manoeuvre]\label{lem:multiset-shift}
Let $k\ge2$ and let $R\le A^{k}$ be totally symmetric and closed under
$(a_{1},\dots,a_{k})\mapsto(a_{1},a_{1},a_{2},\dots,a_{k-1})$. Identify a tuple
of $R$ with its entry multiset, a $k$-element multiset over $A$. If
$\bar a\in R$ has entry multiset $M$, then for any $u,v\in M$ ---
\emph{occupying distinct positions of $\bar a$, though possibly with $u=v$ as
elements of $A$} --- the multiset $M+\{u\}-\{v\}$ is again the entry multiset of
a tuple of $R$. Consequently every $k$-element multiset over the set $S$ of
distinct entries of $\bar a$ is reached, provided it is reached by such steps;
in particular, for any $u\in S$, the constant multiset $k\cdot\{u\}$ is
reached, in at most $k-1$ steps.
\end{lemma}

\begin{proof}
Total symmetry lets us place $u$ first and $v$ last; the shift then duplicates
$u$ and deletes $v$, giving $M+\{u\}-\{v\}$. For the last clause, apply this
$k-1$ times with the same $u$, deleting a different position each time.
\end{proof}

\begin{hazard}[The shift deletes the last entry]\label{haz:shift}
A single application of the shift to $(a,b_{2},\dots,b_{k-1},c)$ destroys $c$,
so $(a,\dots,a,c)\in R$ does not follow from one step; it needs $k-2$ steps
interleaved with permutations, which is what Lemma~\ref{lem:multiset-shift}
packages. The lemma is used twice below and the one-step reading fails both
times.
\end{hazard}

\begin{lemma}\label{lem:tot-sym-full}
Let $k\ge2$, let $\mathbf A$ be BA and centre free, and let $R\le A^{k}$ be
totally symmetric, closed under the shift, with $\proj_{1,2}(R)=A^{2}$. Then
$R=A^{k}$.
\end{lemma}

\begin{proof}
Induct on $k$; for $k=2$ the hypothesis is the conclusion. Let $k>2$.

\emph{Step 1: $(a,\dots,a,c)\in R$ for all $a,c$.} By total symmetry
$\proj_{1,k}(R)=A^{2}$, so there is a tuple of $R$ with first entry $a$ and
last entry $c$. Apply Lemma~\ref{lem:multiset-shift} $k-2$ times with $u=a$,
deleting on each step one position other than the last: the entry multiset
becomes $(k-1)\cdot\{a\}+\{c\}$, which by total symmetry gives
$(a,\dots,a,c)\in R$. Deleting $c$ as well gives $(a,\dots,a)\in R$.

\emph{Step 2: the left projection is full.} $R'=\proj_{1,\dots,k-1}(R)$
inherits total symmetry, has $\proj_{1,2}(R')=A^{2}$, and is shift-closed: from
$\bar a\in R'$ extend to $R$, shift there, and project. By the inductive
hypothesis $R'=A^{k-1}$.

\emph{Step 3.} Only now may $R$ be viewed as a subdirect relation
$R\sdle A^{k-1}\times A$: Step 2 gives the left projection and Step 1 gives the
right. It is linked, since by Step 1 every $c$ is adjacent to the single left
vertex $(a,\dots,a)$. If $R\neq A^{k}$ then
Corollary~\ref{lem:linked-implies-abs} gives a proper nonempty binary
absorbing or central subuniverse on $\mathbf A^{k-1}$ or on $\mathbf A$, and
Lemma~\ref{lem:power-to-base} moves the first case to $\mathbf A$ ---
contradicting the hypothesis either way.
\end{proof}

\begin{lemma}\label{lem:tot-sym-irred}
Let $\sigma$ be a dividing congruence for $B\le A$ and let
$R\le(\mathbf A/\sigma)^{k}$ be reflexive, totally symmetric and shift-closed.
Then $(B/\sigma)^{k}\subseteq R$, or $R$ is the set of constant tuples.
\end{lemma}

\begin{proof}
If $\proj_{1,2}(R)$ is the equality relation then, by total symmetry, all
coordinates of every tuple of $R$ agree, and reflexivity supplies all constant
tuples.

Otherwise $\proj_{1,2}(R)$ is a subuniverse of $(\mathbf A/\sigma)^{2}$
containing $0_{\mathbf A/\sigma}$ properly and stable under it, so by
minimality of the cover (Proposition~\ref{prop:cover}, transported by
Lemma~\ref{lem:irred-quotient}) it contains $\cover{\sigma}/\sigma$, hence
$(B/\sigma)^{2}$ --- the last step by Definition~\ref{def:types}(i). Now apply
Lemma~\ref{lem:tot-sym-full} to $R\cap(B/\sigma)^{k}$ over the algebra
$B/\sigma$, which is BA and centre free by Definition~\ref{def:types}(iii). Its
hypothesis $\proj_{1,2}\bigl(R\cap(B/\sigma)^{k}\bigr)=(B/\sigma)^{2}$ is not
immediate from $\proj_{1,2}(R)\supseteq(B/\sigma)^{2}$, since a witnessing
tuple may leave $B/\sigma$; Lemma~\ref{lem:multiset-shift} repairs it, pulling
$(x_{1},x_{2},y_{3},\dots,y_{k})$ back to $(x_{1},x_{2},x_{1},\dots,x_{1})$.
\end{proof}

\begin{lemma}[Full fibre]\label{lem:full-fibre}
Let $F$ be a finite set with $|F|\le N$ and let $W\subseteq F^{N}\times E$
satisfy $\proj_{1..N}(W)=F^{N}$ and be closed under coordinate substitution: if
$(x_{1},\dots,x_{N},e)\in W$ and $g:[N]\to[N]$ then
$(x_{g(1)},\dots,x_{g(N)},e)\in W$. Then $F^{N}\times\{d\}\subseteq W$ for some
$d\in E$.
\end{lemma}

\begin{proof}
Write $F=\{F_{1},\dots,F_{t}\}$ with $t\le N$ and put
$\bar x=(F_{1},\dots,F_{t},F_{t},\dots,F_{t})\in F^{N}$. By surjectivity pick
$d$ with $(\bar x,d)\in W$. Every $\bar y\in F^{N}$ is $\bar x\circ g$ for some
$g:[N]\to[N]$, since every entry of $\bar y$ occurs among the entries of
$\bar x$; closure gives $(\bar y,d)\in W$.
\end{proof}

\begin{hazard}[Why the arity is $|A|$]\label{haz:full-fibre}
Lemma~\ref{lem:full-fibre} is used three times below, always in the shape
``since $\proj_{1,\dots,|A|}(R')=(B/\delta)^{|A|}$, there is $d$ with
$(B/\delta)^{|A|}\times\{d\}\subseteq R'$''. Surjectivity of a projection does
not by itself produce a full fibre; what makes the step valid is that each such
$R'$ is cut out by a conjunction of unary and binary conditions on its first
$|A|$ entries, hence closed under substituting entries for one another, and
that $|B/\delta|\le|A|$. \textbf{The exponent $|A|$ is load-bearing}: it must be
at least the number of blocks, so that one tuple can enumerate them.
Normalising the arity away breaks all three proofs.
\end{hazard}

\begin{longlemma}\label{lem:self-intersection-pc}
Let $B_{k}\subt[A]{T_{k}(\sigma_{k})}\dots\subt[A]{T_{1}(\sigma_{1})}B_{0}=A$
with $T_{i}\in\{\TBA,\TC,\TS,\TD\}$, let $\delta$ be a congruence on
$\mathbf A$, and let $m\in[k]$ with $T_{m}=\TD$. Then
$\bigl((B_{k}\circ\delta)\cap B_{m-1}\bigr)/\sigma_{m}$ is either
$B_{m-1}/\sigma_{m}$ or $B_{m}/\sigma_{m}$; and in the second case
$\sigma_{m}\supseteq\delta\cap\sigma_{1}\cap\dots\cap\sigma_{m-1}$.
\end{longlemma}

\begin{proof}
The induction is on $k$, and $\delta$, $m$ and the chain are all quantified
\emph{inside} the inductive statement: subsubcase 1B3 below applies the
hypothesis with a different congruence in the role of $\delta$, so $\delta$ must
not be held fixed.

Two conventions are in force throughout. For every $j$ with $T_{j}=\TD$ we have
$B_{j}=B_{j-1}\cap E_{j}$ for a block $E_{j}$ of $\sigma_{j}$, so
$B_{j}^{2}\subseteq\sigma_{j}$, $|B_{j}/\sigma_{j}|=1$, and
$\mathbf B_{j-1}/\sigma_{j}$ is BA and centre free with
$|B_{j-1}/\sigma_{j}|\ge2$ (Definition~\ref{def:types}(iii), and properness of
$B_{j}$ in $B_{j-1}$); in particular $|A|\ge2$. For every other $j$,
$\sigma_{j}=A^{2}$.

\medskip\noindent\textbf{Case 1: $k=m$,} so $T_{k}=\TD$. Put
$\tau=\delta\cap\sigma_{1}\cap\dots\cap\sigma_{k-1}$ and, for
$0\le n\le k$,
\[
  Q_{n}=\bigl\{(a_{1},\dots,a_{|A|})\in B_{n}^{|A|}:
    (a_{i},a_{j})\in\tau\ \text{for all }i,j\bigr\},
  \qquad S_{n}=Q_{n}/\sigma_{k},
\]
the quotient again an image. Each $Q_{n}$ is cut out of $\mathbf A^{|A|}$ by
unary and binary conditions, so each $S_{n}$ is a subuniverse of
$(\mathbf A/\sigma_{k})^{|A|}$ closed under substitution of coordinates and
hence totally symmetric and shift-closed;
$S_{0}$, whose $Q_{0}$ ranges over $B_{0}=A$, is in addition reflexive; and
$S_{n}\supseteq S_{n'}$ for $n\le n'$. Since $\sigma_{k}$ is a dividing
congruence for $B_{k-1}\le A$, Lemma~\ref{lem:tot-sym-irred} applies to
$S_{0}$.

\emph{Subcase 1A: $S_{0}=\{(a/\sigma_{k},\dots,a/\sigma_{k}):a\in A\}$.} For
$(a,b)\in\tau$ the tuple $(a,b,b,\dots,b)$ lies in $Q_{0}$, so its image is
constant and $a/\sigma_{k}=b/\sigma_{k}$: that is
$\sigma_{k}\supseteq\tau$, which is the additional clause. It also settles the
main one. Every $\sigma_{\ell}$ with $\ell<k$ contains $B_{k-1}^{2}$ --- for
$T_{\ell}=\TD$ because $B_{k-1}\subseteq B_{\ell}$ and
$B_{\ell}^{2}\subseteq\sigma_{\ell}$, otherwise because
$\sigma_{\ell}=A^{2}$ --- so
\[
  \delta\cap B_{k-1}^{2}\;\subseteq\;\tau\cap B_{k-1}^{2}\;\subseteq\;\sigma_{k}.
\]
Hence if $a\in B_{k-1}$ and $(a,b)\in\delta$ with $b\in B_{k}\subseteq B_{k-1}$
then $(a,b)\in\sigma_{k}$, so $a\in E_{k}$ and therefore
$a\in B_{k-1}\cap E_{k}=B_{k}$. Thus
$(B_{k}\circ\delta)\cap B_{k-1}=B_{k}$ and the projection is
$B_{k}/\sigma_{k}=B_{m}/\sigma_{m}$.

\emph{Subcase 1B: $(B_{k-1}/\sigma_{k})^{|A|}\subseteq S_{0}$.} Suppose first
that $(B_{k-1}/\sigma_{k})^{|A|}\subseteq S_{k-1}$. Since $|B_{k-1}/\sigma_{k}|
\le|A|$, one tuple $x\in(B_{k-1}/\sigma_{k})^{|A|}$ enumerates every class of
$B_{k-1}/\sigma_{k}$; a witness for $x$ in $Q_{k-1}$ is a family
$a_{1},\dots,a_{|A|}\in B_{k-1}$, pairwise $\tau$-related and so in particular
pairwise $\delta$-related, whose $\sigma_{k}$-classes exhaust
$B_{k-1}/\sigma_{k}$. Let $E$ be the $\delta$-class containing all of them, so
that
\[
  (E\cap B_{k-1})/\sigma_{k}=B_{k-1}/\sigma_{k}.
\]
One of those classes is the class of $B_{k}$, so some $a_{i}$ lies in
$B_{k-1}\cap E_{k}=B_{k}$; hence $E$ meets $B_{k}$, so
$E\subseteq B_{k}\circ\delta$, and
$((B_{k}\circ\delta)\cap B_{k-1})/\sigma_{k}\supseteq(E\cap B_{k-1})/\sigma_{k}
=B_{k-1}/\sigma_{k}$. That is the first alternative, and this branch is done.

So assume $(B_{k-1}/\sigma_{k})^{|A|}\not\subseteq S_{k-1}$ and choose $n$
minimal with $(B_{k-1}/\sigma_{k})^{|A|}\not\subseteq S_{n}$; then
$1\le n\le k-1$ and $(B_{k-1}/\sigma_{k})^{|A|}\subseteq S_{n-1}$. For every
$b\in B_{k-1}\subseteq B_{n}$ the constant tuple $(b,\dots,b)$ lies in
$Q_{n}$, so
\begin{equation}\label{eq:sip-nonempty}
  \varnothing\neq S_{n}\cap(B_{k-1}/\sigma_{k})^{|A|}
  \subsetneq(B_{k-1}/\sigma_{k})^{|A|} .
\end{equation}
We show each of the three types of $T_{n}$ impossible, so that this branch does
not occur.

\emph{Subsubcase 1B1: $T_{n}\in\{\TBA,\TC\}$.} The formula
$\bigwedge_{i,j}\tau(x_{i},x_{j})\wedge\bigwedge_{i}X(x_{i})$ defines
$Q_{n-1}$ at $X=B_{n-1}$ and $Q_{n}$ at $X=B_{n}$, so
Lemma~\ref{lem:pp-propagation} and \eqref{eq:sip-nonempty} give
$Q_{n}\subte{T_{n}}Q_{n-1}$; Lemma~\ref{lem:factor-type} pushes this through
$\sigma_{k}$ and Lemma~\ref{lem:ba-center-intersect} restricts it along
$(B_{k-1}/\sigma_{k})^{|A|}\le\mathbf S_{n-1}$, giving by
\eqref{eq:sip-nonempty} a proper nonempty
$\subt{T_{n}}(B_{k-1}/\sigma_{k})^{|A|}$. Lemma~\ref{lem:power-to-base} moves
it to $B_{k-1}/\sigma_{k}$, which is BA and centre free --- a contradiction.

\emph{Subsubcase 1B2: $T_{n}=\TS$.} Choose a nonempty $G\subseteq B_{n}$ with
$G\subt{\TBA,\TC}B_{n-1}$. Lemma~\ref{lem:strong-nonempty}, applied along
$B_{k}\lll^{A}B_{n-1}\lll A$, gives $G\cap B_{k}\neq\varnothing$, and
$B_{k}\subseteq B_{k-1}$, so replacing $B_{n}$ by $G$ leaves the intersection
with $(B_{k-1}/\sigma_{k})^{|A|}$ nonempty; it stays proper because
$G\subseteq B_{n}$. Now 1B1 applies verbatim with $G$ for $B_{n}$ and
contradicts BA-freeness of $B_{k-1}/\sigma_{k}$.

\emph{Subsubcase 1B3: $T_{n}=\TD$.} Let $R'$ be the image in
$(\mathbf A/\sigma_{k})^{|A|}\times\mathbf A/\sigma_{n}$ of
\[
  \bigl\{(a_{1},\dots,a_{|A|},b):b\in B_{n-1},\ \forall i\;
    a_{i}\in B_{n-1}\cap(B_{k-1}\circ\sigma_{k}),\ \forall i,j\;
    (a_{i},a_{j})\in\tau,\ \forall i\;(a_{i},b)\in\sigma_{n}\bigr\} .
\]
Because $n\le k-1$, the congruence $\sigma_{n}$ is one of those intersected in
$\tau$, so $\tau\subseteq\sigma_{n}$. Four facts.

\emph{(i) $\proj_{1,\dots,|A|}(R')=(B_{k-1}/\sigma_{k})^{|A|}$.} A tuple
$x$ on the right lies in $S_{n-1}$ by minimality of $n$, so it has witnesses
$a_{i}\in B_{n-1}$, pairwise $\tau$-related, with $a_{i}/\sigma_{k}=x_{i}$ and
hence $a_{i}\in B_{k-1}\circ\sigma_{k}$; and $\tau\subseteq\sigma_{n}$ makes
$b=a_{1}\in B_{n-1}$ a legitimate last entry.

\emph{(ii) $\proj_{|A|+1}(R')=((B_{k-1}\circ\sigma_{k})\cap B_{n-1})/\sigma_{n}$.}
A witnessing tuple has $(a_{1},b)\in\sigma_{n}$ with
$a_{1}\in B_{n-1}\cap(B_{k-1}\circ\sigma_{k})$; conversely such a $c$ is
witnessed by $(c,\dots,c,c)$.

\emph{(iii) $(B_{k-1}/\sigma_{k})^{|A|}\times B_{n}/\sigma_{n}\not\subseteq R'$.}
Otherwise, fixing $b_{0}\in B_{k-1}\subseteq B_{n}$, every
$x\in(B_{k-1}/\sigma_{k})^{|A|}$ has a witness with $(a_{i},b_{0})\in
\sigma_{n}$; then $b_{0}\in B_{n}\subseteq E_{n}$ puts each $a_{i}$ in
$B_{n-1}\cap E_{n}=B_{n}$, so $x\in S_{n}$, contradicting the choice of $n$.

\emph{(iv) $\proj_{|A|+1}(R')=B_{n-1}/\sigma_{n}$.} The inductive hypothesis at
$k-1$ --- applied to the chain $B_{k-1}\subt[A]{T_{k-1}(\sigma_{k-1})}\dots
\subt[A]{T_{1}(\sigma_{1})}B_{0}$, the congruence $\sigma_{k}$ in the role of
$\delta$, and the index $n$ --- leaves $B_{n-1}/\sigma_{n}$ or
$B_{n}/\sigma_{n}$ for the set computed in (ii). The latter is a single point,
which with (i) would give the containment excluded by (iii).

Now $R'$ is subdirect by (i) and (iv), it is closed under substitution in its
first $|A|$ coordinates, and $|B_{k-1}/\sigma_{k}|\le|A|$, so
Lemma~\ref{lem:full-fibre} supplies $d$ with
$(B_{k-1}/\sigma_{k})^{|A|}\times\{d\}\subseteq R'$. The associated set
$C^{\star}\subseteq B_{n-1}/\sigma_{n}$ is therefore nonempty, and proper by
(iii); so Lemma~\ref{lem:central-relation} applied to $R'^{-1}$ yields either a
proper nonempty central subuniverse of $B_{n-1}/\sigma_{n}$ or, through
Lemma~\ref{lem:power-to-base}, a proper nonempty binary absorbing subuniverse
of $B_{k-1}/\sigma_{k}$ --- both excluded. So subcase 1B ends at its first
branch.

\medskip\noindent\textbf{Case 2: $k>m$.} Then $B_{k}\subseteq B_{k-1}\subseteq
B_{m}\subseteq B_{m-1}$. The inductive hypothesis at $k-1$, with the same
$\delta$ and $m$, makes $((B_{k-1}\circ\delta)\cap B_{m-1})/\sigma_{m}$ either
$B_{m}/\sigma_{m}$ or $B_{m-1}/\sigma_{m}$.

If it is $B_{m}/\sigma_{m}$, that set is a single point, and
$\varnothing\neq B_{k}\subseteq(B_{k}\circ\delta)\cap B_{m-1}\subseteq
(B_{k-1}\circ\delta)\cap B_{m-1}$ forces
$((B_{k}\circ\delta)\cap B_{m-1})/\sigma_{m}=B_{m}/\sigma_{m}$ as well; the
additional clause is carried over unchanged from the inductive hypothesis,
which had the same $\delta$ and the same $m$. So assume
\begin{equation}\label{eq:sip-case2}
  ((B_{k-1}\circ\delta)\cap B_{m-1})/\sigma_{m}=B_{m-1}/\sigma_{m},
\end{equation}
and split on $T_{k}$. In each subcase the answer is $B_{m-1}/\sigma_{m}$, so
the additional clause is not triggered.

\emph{Subcase 2A: $T_{k}\in\{\TBA,\TC\}$.} The formula
$\exists y\,(X(y)\wedge\delta(x,y))\wedge B_{m-1}(x)$ defines
$(B_{k-1}\circ\delta)\cap B_{m-1}$ at $X=B_{k-1}$ and
$(B_{k}\circ\delta)\cap B_{m-1}$ at $X=B_{k}$, and the latter is nonempty as it
contains $B_{k}$; so Lemmas~\ref{lem:pp-propagation} and
\ref{lem:factor-type} give
$((B_{k}\circ\delta)\cap B_{m-1})/\sigma_{m}\subte{T_{k}}
((B_{k-1}\circ\delta)\cap B_{m-1})/\sigma_{m}=B_{m-1}/\sigma_{m}$. A
\emph{proper} such containment is a proper nonempty binary absorbing or central
subuniverse of $B_{m-1}/\sigma_{m}$, which is excluded, so it is an equality.

\emph{Subcase 2B: $T_{k}=\TS$.} Choose a nonempty $G\subseteq B_{k}$ with
$G\subt{\TBA,\TC}B_{k-1}$ and run 2A with $G$ for $B_{k}$: it gives
$((G\circ\delta)\cap B_{m-1})/\sigma_{m}=B_{m-1}/\sigma_{m}$, and
$G\subseteq B_{k}$ makes the left side a subset of
$((B_{k}\circ\delta)\cap B_{m-1})/\sigma_{m}$.

\emph{Subcase 2C: $T_{k}=\TD$.} Put
\[
  R=\bigl\{(a/\sigma_{k},b/\sigma_{m}):a\in B_{k-1},\ b\in B_{m-1},\
    (a,b)\in\delta\bigr\}\le\mathbf B_{k-1}/\sigma_{k}\times
    \mathbf B_{m-1}/\sigma_{m} .
\]
It is subdirect: $B_{k-1}\subseteq B_{m-1}$ lets $b=a$ serve for the left
projection, and the right projection is $B_{m-1}/\sigma_{m}$ by
\eqref{eq:sip-case2}. Moreover $B_{k-1}/\sigma_{k}\times B_{m}/\sigma_{m}
\subseteq R$: here $k>m$ gives $B_{k-1}\subseteq B_{m}$, so for $a\in B_{k-1}$
the pair $(a,a)$ witnesses $(a/\sigma_{k},a/\sigma_{m})\in R$ with
$a/\sigma_{m}$ the single class of $B_{m}/\sigma_{m}$. \textbf{This is the only
use of $k>m$ in the subcase, and without it $R$ need not have a centre at
all.}

So the set $C^{\star}=\{c\in B_{m-1}/\sigma_{m}:\{c\}\times B_{k-1}/\sigma_{k}
\subseteq R^{-1}\}$ is nonempty. Lemma~\ref{lem:central-relation} applied to
$R^{-1}\sdle\mathbf B_{m-1}/\sigma_{m}\times\mathbf B_{k-1}/\sigma_{k}$ leaves
only its alternative (b), alternative (c) being excluded by BA-freeness of
$B_{k-1}/\sigma_{k}$; and since $B_{m-1}/\sigma_{m}$ is centre free the central
subuniverse $C^{\star}$ cannot be proper. Hence
$C^{\star}=B_{m-1}/\sigma_{m}$ and $R$ is the full product.

Finally, fullness transfers from $B_{k-1}$ to $B_{k}$. Let $a\in B_{k}$ and
$b\in B_{m-1}$. Fullness gives $a'\in B_{k-1}$ and $b'\in B_{m-1}$ with
$(a',b')\in\delta$, $a'/\sigma_{k}=a/\sigma_{k}$ and
$b'/\sigma_{m}=b/\sigma_{m}$. From $a\in B_{k}\subseteq E_{k}$ and
$a'\,\sigma_{k}\,a$ we get $a'\in E_{k}$, so $a'\in B_{k-1}\cap E_{k}=B_{k}$;
hence $b'\in(B_{k}\circ\delta)\cap B_{m-1}$ and $b'/\sigma_{m}=b/\sigma_{m}$.
As $b$ was arbitrary, $((B_{k}\circ\delta)\cap B_{m-1})/\sigma_{m}=
B_{m-1}/\sigma_{m}$.
\end{proof}

\begin{hazard}[Three steps this proof supplies]\label{haz:self-intersection}
Subcase 1B's first branch is where the projection is actually computed, and it
needs two things beyond the displayed identity
$(E\cap B_{k-1})/\sigma_{k}=B_{k-1}/\sigma_{k}$: that a \emph{single}
$\delta$-class $E$ works, which comes from the witnesses for one tuple
enumerating all of $B_{k-1}/\sigma_{k}$ and so from $|B_{k-1}/\sigma_{k}|\le
|A|$ again; and that $E$ \emph{meets} $B_{k}$, without which $E$ has nothing to
do with $B_{k}\circ\delta$. The second follows from the first, since the
classes exhausted include the class of $B_{k}$.

Subcase 2C ends with a transfer step: $R$ full is a statement about
$B_{k-1}$, and the conclusion is about $B_{k}$. It goes through only because
$B_{k}$ is cut out of $B_{k-1}$ by a $\sigma_{k}$-block, so a $\sigma_{k}$-move
inside $B_{k-1}$ cannot leave $B_{k}$ --- the same observation that drives
subcase 1A.

The proof does not need $\delta$ to be irreducible, or comparable with any
$\sigma_{i}$; it is an arbitrary congruence, and the additional clause is the
only place any relation between $\delta$ and $\sigma_{m}$ is asserted.
\end{hazard}

\begin{longlemma}[The intersection property]\label{lem:intersection-good}
Let $B\lll A$, $C\lll A$ with $B\cap C\neq\varnothing$. Then
\begin{enumerate}\alphlabels\tightlist
\item[\textup{(d)}] if $\delta$ is a dividing congruence for $B\le A$ then
  $|(B\cap C)/\delta|=1$ or $(B\cap C)/\delta=B/\delta$; and if
  $|(B\cap C)/\delta|=1$ then $\delta\supseteq\delta_{1}\cap\dots\cap
  \delta_{s}$, where $\delta_{1},\dots,\delta_{s}$ are the congruences of the
  \emph{chosen} chains witnessing $B\lll A$ and $C\lll A$;
\item[\textup{(s)}] if $G\subt{\TBA,\TC}B$ then $G\cap C\neq\varnothing$.
\end{enumerate}
\end{longlemma}

\begin{proof}
Fix chains
\[
  B=B_{k}\subt[A]{T_{k}(\sigma_{k})}\dots\subt[A]{T_{1}(\sigma_{1})}B_{0}=A,
  \qquad
  C=C_{\ell}\subt[A]{\mathcal T_{\ell}(\omega_{\ell})}\dots
  \subt[A]{\mathcal T_{1}(\omega_{1})}C_{0}=A
\]
witnessing $B\lll A$ and $C\lll A$, with all $T_{i},\mathcal T_{j}\in
\{\TBA,\TC,\TS,\TD\}$ and, by Definition~\ref{def:types}, $\sigma_{i}=A^{2}$
whenever $T_{i}\neq\TD$ and $\omega_{j}=A^{2}$ whenever
$\mathcal T_{j}\neq\TD$. Put $\sigma=\bigcap_{i=1}^{k}\sigma_{i}$ and
$\omega=\bigcap_{j=1}^{\ell}\omega_{j}$, so that
$\sigma\cap\omega=\delta_{1}\cap\dots\cap\delta_{s}$ is exactly the
intersection named in the ``moreover'' clause. The induction is on $k+\ell$
over \emph{pairs of chains} --- every appeal below names the shorter pair it is
made to --- and proves (s) and (d) together.

Three facts are used throughout and are recorded once.

\emph{(F1) $B^{2}\subseteq\sigma$ and $C^{2}\subseteq\omega$.} If $T_{i}=\TD$
then $B_{i}=B_{i-1}\cap E_{i}$ for a block $E_{i}$ of $\sigma_{i}$, so
$B_{i}^{2}\subseteq\sigma_{i}$ and a fortiori $B^{2}\subseteq\sigma_{i}$;
otherwise $\sigma_{i}=A^{2}$. Likewise for $C$. In particular
$(a,b)\in\sigma\cap\omega$ for all $a,b\in B\cap C$.

\emph{(F2) $\sigma\subseteq\sigma_{i}$ and $\omega\subseteq\omega_{j}$} for
every $i$ and $j$, by construction.

\emph{(F3) $|A|\ge2$.} That $\delta$ is a dividing congruence for $B\le A$
means, by Definition~\ref{def:types}, that $B'\subt[A]{\TD(\delta)}B$ for some
$B'$; such a $B'$ is a nonempty proper subset $B\cap E$ of $B$ cut out by a
$\delta$-block, so $|B/\delta|\ge2$, and hence $|A|\ge|B|\ge2$. The same
appeal supplies the two properties of $\delta$ used below: $\mathbf B/\delta$
is BA and centre free, and $B^{2}\subseteq\cover{\delta}$.

\medskip\noindent\textbf{Part (s).} Let $G\subt{\TBA,\TC}B$. If $\ell=0$ then
$C=A\supseteq G\neq\varnothing$. So let $\ell\ge1$. The chain for $C$
truncated at $C_{\ell-1}$ witnesses $C_{\ell-1}\lll A$ with $\ell-1$ steps, and
$B\cap C_{\ell-1}\supseteq B\cap C\neq\varnothing$; so the inductive hypothesis
(s), at $k+(\ell-1)$, gives
\begin{equation}\label{eq:ig-s-start}
  G\cap C_{\ell-1}\neq\varnothing .
\end{equation}
Write $A'=B\cap C_{\ell-1}$, a nonempty subuniverse of $\mathbf A$. Since
$G\subseteq B$ we have $G\cap X=G\cap A'\cap X$ for every $X\subseteq
C_{\ell-1}$, and $B\cap C_{\ell}\supseteq B\cap C\neq\varnothing$. By
Lemma~\ref{lem:ba-center-intersect} applied twice, once for $\TBA$ and once for
$\TC$,
\begin{equation}\label{eq:ig-s-cut}
  G\cap C_{\ell-1}=G\cap A'\dsubte{\TBA,\TC}A',
\end{equation}
and by \eqref{eq:ig-s-start} the dot may be dropped. Split on
$\mathcal T_{\ell}$.

\emph{Case $\mathcal T_{\ell}=\TD$,} with $C_{\ell}=C_{\ell-1}\cap E$ for a
block $E$ of $\omega_{\ell}$. Apply the inductive hypothesis (d) to the pair
$(C_{\ell-1},B)$ --- that is, with $C_{\ell-1}$ in the role of the first
subuniverse and $B$ in the role of the second --- and to the dividing
congruence $\omega_{\ell}$ for $C_{\ell-1}\le A$, at $(\ell-1)+k$. Either
$|A'/\omega_{\ell}|=1$ or $A'/\omega_{\ell}=C_{\ell-1}/\omega_{\ell}$.

In the first case the sole $\omega_{\ell}$-class meeting $A'$ is $E$, because
$\varnothing\neq B\cap C\subseteq A'\cap E$; hence $A'\subseteq E$ and
$A'=A'\cap E=B\cap C_{\ell}$, so
$G\cap C_{\ell}=G\cap A'\cap C_{\ell}=G\cap A'=G\cap C_{\ell-1}\neq\varnothing$.

In the second case Lemma~\ref{lem:factor-type} applied to
\eqref{eq:ig-s-cut} gives
$(G\cap C_{\ell-1})/\omega_{\ell}\subte{\TBA,\TC}A'/\omega_{\ell}=
C_{\ell-1}/\omega_{\ell}$. If that containment is proper it exhibits a proper
nonempty binary absorbing subuniverse of $C_{\ell-1}/\omega_{\ell}$,
contradicting the requirement of Definition~\ref{def:types}(iii) for the step
$C_{\ell}\subt[A]{\TD(\omega_{\ell})}C_{\ell-1}$ that
$\mathbf C_{\ell-1}/\omega_{\ell}$ be BA and centre free. If it is an equality
then every $\omega_{\ell}$-class meeting $C_{\ell-1}$ already meets
$G\cap C_{\ell-1}$; the class $E$ does meet $C_{\ell-1}$, since
$C_{\ell}\neq\varnothing$, so there is
$g\in G\cap C_{\ell-1}\cap E=G\cap C_{\ell}$.

\emph{Case $\mathcal T_{\ell}\in\{\TBA,\TC\}$.} By
Lemma~\ref{lem:ba-center-intersect}, $B\cap C_{\ell}=A'\cap C_{\ell}
\dsubte{\mathcal T_{\ell}}A'$, and it is nonempty. If
$G\cap C_{\ell-1}=A'$ then $A'\subseteq G$, and $\varnothing\neq B\cap
C_{\ell}\subseteq A'$ gives $G\cap C_{\ell}\neq\varnothing$. If
$B\cap C_{\ell}=A'$ then $G\cap C_{\ell}=G\cap A'\cap C_{\ell}=G\cap A'\neq
\varnothing$. Otherwise both containments are proper, so
$G\cap C_{\ell-1}\subt{\TBA,\TC}A'$ and $B\cap C_{\ell}
\subt{\mathcal T_{\ell}}A'$, and Lemma~\ref{lem:ba-central-meets} inside
$\mathbf A'$ makes their intersection nonempty; that intersection is contained
in $G\cap C_{\ell}$.

\emph{Case $\mathcal T_{\ell}=\TS$.} By Definition~\ref{def:ba-center-free}
there is a nonempty $G'\subseteq C_{\ell}$ with $G'\subt{\TBA,\TC}C_{\ell-1}$
--- proper, since $G'\subseteq C_{\ell}\lneq C_{\ell-1}$. The inductive
hypothesis (s) applied to the pair $(C_{\ell-1},B)$ with the subuniverse $G'$,
at $(\ell-1)+k$, gives $B\cap G'\neq\varnothing$; so
$B\cap G'=A'\cap G'\subte{\TBA,\TC}A'$ by
Lemma~\ref{lem:ba-center-intersect}. If $G\cap C_{\ell-1}=A'$ then
$\varnothing\neq B\cap G'\subseteq A'\subseteq G$ and $B\cap G'\subseteq
C_{\ell}$, so $G\cap C_{\ell}\neq\varnothing$. If $B\cap G'=A'$ then
$A'\subseteq G'\subseteq C_{\ell}$, so $G\cap C_{\ell}\supseteq G\cap A'\neq
\varnothing$. Otherwise both are proper and
Lemma~\ref{lem:ba-central-meets} gives
$\varnothing\neq(G\cap C_{\ell-1})\cap(B\cap G')\subseteq G\cap C_{\ell}$.

\medskip\noindent\textbf{Part (d).} For $0\le m\le k$ and $0\le n\le\ell$ put
\[
  P_{m,n}=\bigl\{(a_{1},\dots,a_{|A|})\in(B_{m}\cap C_{n})^{|A|}:
    (a_{i},a_{j})\in\sigma\cap\omega\ \text{for all }i,j\bigr\},
  \qquad S_{m,n}=P_{m,n}/\delta ,
\]
the quotient being coordinatewise and, as always, an \emph{image}
(\S\ref{sec:legislation}, item~\ref{leg:quotient}): a tuple
$x\in(A/\delta)^{|A|}$ lies in $S_{m,n}$ exactly when \emph{one} tuple of
representatives of its entries lies in $P_{m,n}$. Each $P_{m,n}$ is a
subuniverse of $\mathbf A^{|A|}$, being cut out by unary and binary
conditions, so each $S_{m,n}$ is a subuniverse of $(\mathbf A/\delta)^{|A|}$;
each is closed under substitution $x\mapsto x\circ g$ along any
$g:[|A|]\to[|A|]$, because every defining condition is unary in $a_{i}$ or
binary in $(a_{i},a_{j})$, and in particular each is totally symmetric and
shift-closed; and $S_{m,n}\supseteq S_{m',n'}$ whenever $m\le m'$ and
$n\le n'$. Write $S=S_{0,0}$; since $B_{0}\cap C_{0}=A$ and
$(a,a)\in\sigma\cap\omega$ for every $a$, this one --- and only this one --- is
also reflexive.

By (F3) the arity $|A|$ is at least $2$, so Lemma~\ref{lem:tot-sym-irred}
applies to $S$ and the dividing congruence $\delta$ for $B\le A$, leaving two
cases.

\emph{Case 1: $S=\{(a/\delta,\dots,a/\delta):a\in A\}$.} Let
$(a,b)\in\sigma\cap\omega$. Since $|A|\ge2$ the tuple $(a,b,b,\dots,b)$ lies
in $P_{0,0}$, its entries being pairwise $\sigma\cap\omega$-related; so its
image lies in $S$ and is therefore constant, giving $a/\delta=b/\delta$. Hence
$\sigma\cap\omega\subseteq\delta$, which is the ``moreover'' clause. By (F1)
any $a,b\in B\cap C$ satisfy $(a,b)\in\sigma\cap\omega\subseteq\delta$, and
$B\cap C\neq\varnothing$, so $|(B\cap C)/\delta|=1$.

\emph{Case 2: $(B/\delta)^{|A|}\subseteq S$.} We show that the only surviving
subcase is 2C, in which $(B\cap C)/\delta=B/\delta$; by (F3) that alternative
has $|(B\cap C)/\delta|=|B/\delta|\ge2$, so the ``moreover'' clause is not
triggered in Case 2 and nothing further is owed.

\emph{Subcase 2A: $(B/\delta)^{|A|}\not\subseteq S_{k,0}$.} Choose $m$ minimal
with $(B/\delta)^{|A|}\not\subseteq S_{m,0}$; since $S_{0,0}=S$ we have
$m\ge1$ and $(B/\delta)^{|A|}\subseteq S_{m-1,0}$, so $(B/\delta)^{|A|}$ is a
subalgebra of $\mathbf S_{m-1,0}$. Note that for every $b\in B$ the constant
tuple $(b,\dots,b)$ lies in $P_{m,0}$, because $B\subseteq B_{m}$ and
$(b,b)\in\sigma\cap\omega$; so
\begin{equation}\label{eq:ig-2A-nonempty}
  \varnothing\neq S_{m,0}\cap(B/\delta)^{|A|}\subsetneq(B/\delta)^{|A|},
\end{equation}
the properness being the choice of $m$. Split on $T_{m}$.

\emph{Subsubcase 2A1: $T_{m}\in\{\TBA,\TC\}$.} The formula
\[
  \Phi_{X}(x_{1},\dots,x_{|A|})=
  \bigwedge_{i,j}(\sigma\cap\omega)(x_{i},x_{j})\wedge\bigwedge_{i}X(x_{i})
\]
is positive primitive in the unary relation $X$, and defines $P_{m-1,0}$ at
$X=B_{m-1}$ and $P_{m,0}$ at $X=B_{m}$. Since $B_{m}\subt{T_{m}}B_{m-1}$,
Lemma~\ref{lem:pp-propagation} gives $P_{m,0}\dsubte{T_{m}}P_{m-1,0}$, and
$P_{m,0}\neq\varnothing$ by the constant tuples, so
$P_{m,0}\subte{T_{m}}P_{m-1,0}$. Lemma~\ref{lem:factor-type}, applied to that
containment inside $\mathbf P_{m-1,0}$ with the congruence induced by
$\delta^{|A|}$, gives $S_{m,0}\subte{T_{m}}S_{m-1,0}$, and then
Lemma~\ref{lem:ba-center-intersect} restricts along the subalgebra
$(B/\delta)^{|A|}\le\mathbf S_{m-1,0}$:
\[
  S_{m,0}\cap(B/\delta)^{|A|}\dsubte{T_{m}}(B/\delta)^{|A|} .
\]
By \eqref{eq:ig-2A-nonempty} the left side is nonempty and proper, so this is
$\subt{T_{m}}$, and Lemma~\ref{lem:power-to-base} produces a proper nonempty
subuniverse of $\mathbf B/\delta$ of type $T_{m}\in\{\TBA,\TC\}$. That
contradicts (F3).

\emph{Subsubcase 2A2: $T_{m}=\TS$.} By Definition~\ref{def:ba-center-free}
choose a nonempty $G\subseteq B_{m}$ with $G\subt{\TBA,\TC}B_{m-1}$. Since
$B\lll^{A}B_{m-1}\lll A$ along the two halves of the chain,
Lemma~\ref{lem:strong-nonempty} gives $G\cap B\neq\varnothing$. Run 2A1 with
$G$ in place of $B_{m}$: writing $P^{G}=\Phi_{G}$ and $S^{G}=P^{G}/\delta$,
the same three lemmas give $S^{G}\cap(B/\delta)^{|A|}\dsubte{\TBA}
(B/\delta)^{|A|}$; it is nonempty, since $b\in G\cap B$ contributes its
constant tuple, and proper, since $G\subseteq B_{m}$ makes it a subset of the
proper set in \eqref{eq:ig-2A-nonempty}. Lemma~\ref{lem:power-to-base} again
contradicts (F3). (Only the $\TBA$ half of $G\subt{\TBA,\TC}B_{m-1}$ is spent
here.)

\emph{Subsubcase 2A3: $T_{m}=\TD$,} with $B_{m}=B_{m-1}\cap E$ for a block $E$
of $\sigma_{m}$. Let $B\circ\delta$ denote the $\delta$-saturation of $B$,
a subuniverse of $\mathbf A$, and note $a/\delta\in B/\delta\iff a\in
B\circ\delta$. Put
\[
  P'=\bigl\{(a_{1},\dots,a_{|A|},b): b\in B_{m-1},\ \forall i\;
    a_{i}\in B_{m-1}\cap(B\circ\delta),\ \forall i,j\;(a_{i},a_{j})\in
    \sigma\cap\omega,\ \forall i\;(a_{i},b)\in\sigma_{m}\bigr\}
\]
and let $R'\le(\mathbf A/\delta)^{|A|}\times\mathbf A/\sigma_{m}$ be its image
under $(a_{1},\dots,a_{|A|},b)\mapsto(a_{1}/\delta,\dots,a_{|A|}/\delta,
b/\sigma_{m})$, so that $R'\subseteq(B/\delta)^{|A|}\times B_{m-1}/\sigma_{m}$.
Four facts about $R'$.

\emph{(i) $\proj_{1,\dots,|A|}(R')=(B/\delta)^{|A|}$.} Let
$x\in(B/\delta)^{|A|}$. By minimality of $m$, $x\in S_{m-1,0}$, so there are
$a_{i}\in B_{m-1}$, pairwise $\sigma\cap\omega$-related, with
$a_{i}/\delta=x_{i}$; each $a_{i}$ then lies in $B\circ\delta$. By (F2),
$(a_{i},a_{j})\in\sigma\subseteq\sigma_{m}$, so $b=a_{1}\in B_{m-1}$ satisfies
$(a_{i},b)\in\sigma_{m}$ for every $i$, and $(a_{1},\dots,a_{|A|},b)\in P'$.

\emph{(ii) $\proj_{|A|+1}(R')=((B\circ\delta)\cap B_{m-1})/\sigma_{m}$.} For
$\subseteq$: a witnessing tuple has $a_{1}\in B_{m-1}\cap(B\circ\delta)$ and
$(a_{1},b)\in\sigma_{m}$, so $b/\sigma_{m}=a_{1}/\sigma_{m}$. For
$\supseteq$: given $c$ in that set, the tuple $(c,\dots,c,c)$ lies in $P'$.

\emph{(iii) $(B/\delta)^{|A|}\times B_{m}/\sigma_{m}\not\subseteq R'$.} Suppose
otherwise and fix $b_{0}\in B\subseteq B_{m}$. For any
$x\in(B/\delta)^{|A|}$ the pair $(x,b_{0}/\sigma_{m})$ lies in $R'$, so it has
a witness $(a_{1},\dots,a_{|A|},b)\in P'$ with $b/\sigma_{m}=b_{0}/\sigma_{m}$;
then $(a_{i},b_{0})\in\sigma_{m}$, and $b_{0}\in B_{m}\subseteq E$ forces
$a_{i}\in E$, whence $a_{i}\in B_{m-1}\cap E=B_{m}$. So the witness lies in
$P_{m,0}$ and $x\in S_{m,0}$ --- for every $x$, contradicting the choice of
$m$.

\emph{(iv) $\proj_{|A|+1}(R')=B_{m-1}/\sigma_{m}$.} Apply
Lemma~\ref{lem:self-intersection-pc} to the chain for $B$, the congruence
$\delta$ and the index $m$: by (ii) the projection is either
$B_{m-1}/\sigma_{m}$ or $B_{m}/\sigma_{m}$. In the second case
$B_{m}\subseteq E$ makes $B_{m}/\sigma_{m}$ a single point, so by (i) every
$x\in(B/\delta)^{|A|}$ is paired in $R'$ with that point, i.e.\
$(B/\delta)^{|A|}\times B_{m}/\sigma_{m}\subseteq R'$, contradicting (iii).

By (i) and (iv), $R'$ is subdirect in $(\mathbf B/\delta)^{|A|}\times
\mathbf B_{m-1}/\sigma_{m}$; both factors are finite algebras with the single
basic operation $w$, so Convention~\ref{conv:standing} holds for them. $R'$ is
closed under substitution in its first $|A|$ coordinates, since every defining
condition of $P'$ is unary in $a_{i}$ or binary in $(a_{i},a_{j})$ or
$(a_{i},b)$; and $|B/\delta|\le|A|$. So Lemma~\ref{lem:full-fibre}, with
$F=B/\delta$ and $N=|A|$, supplies $d\in B_{m-1}/\sigma_{m}$ with
$(B/\delta)^{|A|}\times\{d\}\subseteq R'$.

Now apply Lemma~\ref{lem:central-relation} to the transpose
$R'^{-1}=\{(y,x):(x,y)\in R'\}\sdle\mathbf B_{m-1}/\sigma_{m}\times
(\mathbf B/\delta)^{|A|}$, whose associated set is
$C^{\star}=\{c\in B_{m-1}/\sigma_{m}:\{c\}\times(B/\delta)^{|A|}\subseteq
R'^{-1}\}$. Alternative (a) fails, since $d\in C^{\star}$. And $C^{\star}$ is
\emph{proper}: were it all of $B_{m-1}/\sigma_{m}$ we should have
$R'=(B/\delta)^{|A|}\times B_{m-1}/\sigma_{m}$, which contains
$(B/\delta)^{|A|}\times B_{m}/\sigma_{m}$ and contradicts (iii). So
alternative (b) makes $C^{\star}$ a proper nonempty central subuniverse of
$B_{m-1}/\sigma_{m}$, contradicting Definition~\ref{def:types}(iii) for the
step $B_{m}\subt[A]{\TD(\sigma_{m})}B_{m-1}$; and alternative (c) gives a
proper nonempty binary absorbing subuniverse of the \emph{power}
$(B/\delta)^{|A|}$, which Lemma~\ref{lem:power-to-base} moves down to
$\mathbf B/\delta$, contradicting (F3). Both obligations of
Caveat~\ref{haz:central-relation} are thereby discharged.

\emph{Subcase 2B: $(B/\delta)^{|A|}\subseteq S_{k,0}$ but
$(B/\delta)^{|A|}\not\subseteq S_{k,\ell}$.} Choose $n$ minimal with
$(B/\delta)^{|A|}\not\subseteq S_{k,n}$; then $n\ge1$ and
$(B/\delta)^{|A|}\subseteq S_{k,n-1}$. Fixing once and for all $b_{1}\in
B\cap C\subseteq B_{k}\cap C_{n}$ --- this is the one place where the
hypothesis $B\cap C\neq\varnothing$ is spent inside Case 2 --- its constant
tuple lies in $P_{k,n}$, so
\begin{equation}\label{eq:ig-2B-nonempty}
  \varnothing\neq S_{k,n}\cap(B/\delta)^{|A|}\subsetneq(B/\delta)^{|A|} .
\end{equation}
Split on $\mathcal T_{n}$.

\emph{Subsubcase 2B1: $\mathcal T_{n}\in\{\TBA,\TC\}$.} Verbatim 2A1, with the
pp-formula
$\bigwedge_{i,j}(\sigma\cap\omega)(x_{i},x_{j})\wedge\bigwedge_{i}B_{k}(x_{i})
\wedge\bigwedge_{i}X(x_{i})$, with $X$ instantiated at $C_{n-1}$ and at
$C_{n}$, and with \eqref{eq:ig-2B-nonempty} in place of
\eqref{eq:ig-2A-nonempty}. It contradicts (F3).

\emph{Subsubcase 2B2: $\mathcal T_{n}=\TS$.} Choose a nonempty $G\subseteq
C_{n}$ with $G\subt{\TBA,\TC}C_{n-1}$. Here $B$ and $C_{n-1}$ lie on different
chains, so Lemma~\ref{lem:strong-nonempty} does not apply and the induction is
used instead: the inductive hypothesis (s) for the pair $(C_{n-1},B)$, at
$(n-1)+k$, with $C_{n-1}\cap B\supseteq C\cap B\neq\varnothing$, gives
$G\cap B\neq\varnothing$. Now run 2A2 with $G$ in place of $C_{n}$: the
witness $b\in G\cap B$ lies in $B_{k}\cap C_{n-1}$ and makes the intersection
nonempty, $G\subseteq C_{n}$ makes it proper by
\eqref{eq:ig-2B-nonempty}, and Lemma~\ref{lem:power-to-base} contradicts (F3).

\emph{Subsubcase 2B3: $\mathcal T_{n}=\TD$,} with $C_{n}=C_{n-1}\cap E$ for a
block $E$ of $\omega_{n}$. Put
\[
  P'=\bigl\{(a_{1},\dots,a_{|A|},b): b\in C_{n-1},\ \forall i\;
    a_{i}\in B_{k}\cap C_{n-1},\ \forall i,j\;(a_{i},a_{j})\in\sigma\cap\omega,
    \ \forall i\;(a_{i},b)\in\omega_{n}\bigr\}
\]
and let $R'$ be its image in $(\mathbf A/\delta)^{|A|}\times\mathbf A/
\omega_{n}$. No saturation appears in the first $|A|$ coordinates this time,
because $a_{i}\in B_{k}=B$ already gives $a_{i}/\delta\in B/\delta$.

\emph{(i)} $\proj_{1,\dots,|A|}(R')=(B/\delta)^{|A|}$: for
$x\in(B/\delta)^{|A|}\subseteq S_{k,n-1}$ take witnesses $a_{i}\in B_{k}\cap
C_{n-1}$; by (F2) $(a_{i},a_{j})\in\omega\subseteq\omega_{n}$, so
$b=a_{1}\in C_{n-1}$ serves.
\emph{(ii)} $\proj_{|A|+1}(R')=(B_{k}\cap C_{n-1})/\omega_{n}$, by the
argument of 2A3(ii).
\emph{(iii)} $(B/\delta)^{|A|}\times C_{n}/\omega_{n}\not\subseteq R'$: with
$b_{1}\in B\cap C\subseteq B_{k}\cap C_{n}$ in place of $b_{0}$, the argument
of 2A3(iii) puts every witness in $C_{n-1}\cap E=C_{n}$ and so every
$x\in(B/\delta)^{|A|}$ in $S_{k,n}$, contradicting the choice of $n$.
\emph{(iv)} $\proj_{|A|+1}(R')=C_{n-1}/\omega_{n}$: the inductive hypothesis
(d) for the pair $(C_{n-1},B_{k})$ with the dividing congruence $\omega_{n}$
for $C_{n-1}\le A$, at $(n-1)+k$, leaves $|(B_{k}\cap C_{n-1})/\omega_{n}|=1$
or $(B_{k}\cap C_{n-1})/\omega_{n}=C_{n-1}/\omega_{n}$. In the first case the
sole class is $E$, since $b_{1}\in B_{k}\cap C_{n}\subseteq E$; so
$B_{k}\cap C_{n-1}\subseteq E$ and $B_{k}\cap C_{n-1}=B_{k}\cap C_{n-1}\cap
E=B_{k}\cap C_{n}$, whence $P_{k,n-1}=P_{k,n}$ and $S_{k,n-1}=S_{k,n}$,
contradicting the choice of $n$.

From here 2A3 runs word for word with $\omega_{n}$ for $\sigma_{m}$,
$C_{n-1}$ for $B_{m-1}$ and $C_{n}$ for $B_{m}$: a full fibre by
Lemma~\ref{lem:full-fibre}, a proper
nonempty $C^{\star}$, and Lemma~\ref{lem:central-relation} contradicting either
Definition~\ref{def:types}(iii) for $C_{n}\subt[A]{\TD(\omega_{n})}C_{n-1}$ or
(F3).

\emph{Subcase 2C: $(B/\delta)^{|A|}\subseteq S_{k,\ell}$.} For $b\in B$ the
constant tuple $(b/\delta,\dots,b/\delta)$ lies in $S_{k,\ell}$, so it has a
witness in $(B_{k}\cap C_{\ell})^{|A|}=(B\cap C)^{|A|}$, whose first entry $a$
satisfies $a/\delta=b/\delta$ with $a\in B\cap C$. Hence $B/\delta\subseteq
(B\cap C)/\delta$, and the reverse containment is trivial, so
$(B\cap C)/\delta=B/\delta$.
\end{proof}

\begin{hazard}[The four places this proof can be misread]\label{haz:intersection-draft}
\begin{enumerate}\alphlabels\tightlist
\item \emph{$S_{m,n}$ is an image, not an intersection.} A tuple lies in it
  when \emph{one} tuple of representatives satisfies all the conditions at
  once, and those conditions couple the representatives to each other. So
  $S_{m,n}$ is not a power of $(B_{m}\cap C_{n})/\delta$ cut down by $S$, and
  the membership tests in 2A3(iii), 2B3(iii) and 2C all work by producing or
  constraining a single witnessing tuple.
\item \emph{Lemmas~\ref{lem:ba-center-intersect} and \ref{lem:factor-type}
  deliver $\subte{}$, never $\subt{}$.} Part (s) therefore splits five times on
  whether a containment is proper, and in each case the equality branch
  \emph{proves the lemma} rather than contradicting anything: an absorbing
  subuniverse that is everything is no obstruction, but it does make the set on
  the left swallow the set one is trying to meet. Reading these lemmas as
  producing proper subuniverses collapses all five splits, and it is the most
  tempting error available here.
\item \emph{The second alternative of Lemma~\ref{lem:self-intersection-pc} has
  to be excluded by hand}, in 2A3(iv), and likewise its analogue in 2B3(iv).
  What kills it is that $B_{m}$ lies inside a single $\sigma_{m}$-block: the
  alternative would make the last coordinate of $R'$ constant and hence, with
  surjectivity of the first $|A|$, force exactly the containment that the
  minimal choice of $m$ forbids. In 2B3(iv) the same manoeuvre also identifies
  \emph{which} class the single one is --- the one defining $C_{n}$, witnessed
  by a point of $B\cap C$.
\item \emph{Every appeal to Lemma~\ref{lem:central-relation} owes two
  discharges}, per Caveat~\ref{haz:central-relation}: the centre must be shown
  proper --- here by the same non-containment that the choice of $m$ or $n$
  provides --- and the binary absorbing subuniverse of alternative (c) lives on
  a power and must be moved down by Lemma~\ref{lem:power-to-base}.
\end{enumerate}

Two dependencies remain outside the proof: it rests on
Lemma~\ref{lem:self-intersection-pc}, which is still an outline, and on the
imported Lemma~\ref{lem:power-to-base}. Minimality of the chains for $B$ and
$C$ is \emph{not} among them, and could not be: the inductive appeals in part
(s), 2B2 and 2B3 are to truncations of the chain for $C$, which need not be
minimal even when the original is.
\end{hazard}

\begin{corollary}\label{cor:intersection-good}
Let $R\sdle\mathbf A_{1}\times\mathbf A_{2}$, $B_{1}\lll A_{1}$,
$B_{2}\lll A_{2}$, and let $\sigma$ be a dividing congruence for
$B_{1}\lll A_{1}$. Then $\proj_{1}\bigl(R\cap(B_{1}\times B_{2})\bigr)/\sigma$
is empty, of size $1$, or equal to $B_{1}/\sigma$.
\end{corollary}

\begin{proof}
Let $f_{i}:\mathbf R\to\mathbf A_{i}$ be the projections. By
Lemma~\ref{lem:reverse-hom}(b), $f_{i}^{-1}(B_{i})\lll\mathbf R$; apply
Lemma~\ref{lem:intersection-good}(d) inside $\mathbf R$ to these two, with the
congruence $f_{1}^{-1}(\sigma)$, and read the conclusion back through $f_{1}$.
\end{proof}

\begin{hazard}\label{haz:cor-intersection}
The alternative ``empty'' is present in
Corollary~\ref{cor:intersection-good} and absent from
Lemma~\ref{lem:intersection-good}, which hypothesises $B\cap C\neq\varnothing$.
It is a real alternative, and every call site below has to exclude it
separately.
\end{hazard}

\begin{lemma}[Parallelogram property for $\TD$]\label{lem:parallelogram-D}
Let $B_{1},B_{2}\lll A$, $C_{1},C_{1}'\subt[A]{\TD(\sigma_{1})}B_{1}$ and
$C_{2},C_{2}'\subt[A]{\TD(\sigma_{2})}B_{2}$ with $C_{1}'\cap C_{2}$,
$C_{1}\cap C_{2}'$ and $C_{1}'\cap C_{2}'$ all nonempty. Then
$C_{1}\cap C_{2}\neq\varnothing$.
\end{lemma}

\begin{proof}
By Lemma~\ref{lem:intersection-good}(d), $(B_{1}\cap B_{2})/\sigma_{1}$ is of
size $1$ or all of $B_{1}/\sigma_{1}$. In the first case the single block is
the one defining $C_{1}$, because $C_{1}\cap C_{2}'\neq\varnothing$ puts a
point of $C_{1}$ in $B_{1}\cap B_{2}$; hence $B_{1}\cap B_{2}\subseteq C_{1}$
and $C_{1}\cap C_{2}=B_{1}\cap C_{2}\supseteq C_{1}'\cap C_{2}\neq\varnothing$.
The symmetric argument applies if $(B_{1}\cap B_{2})/\sigma_{2}$ has size $1$.

So assume $(B_{1}\cap B_{2})/\sigma_{1}=B_{1}/\sigma_{1}$ and
$(B_{1}\cap B_{2})/\sigma_{2}=B_{2}/\sigma_{2}$, and put
$S=\{(a/\sigma_{1},a/\sigma_{2}):a\in B_{1}\cap B_{2}\}$, which is then
subdirect in $(B_{1}/\sigma_{1})\times(B_{2}/\sigma_{2})$. Apply
Lemma~\ref{lem:intersection-good}(d) to each $\sigma_{1}$-class $C$ of $B_{1}$
against $B_{2}$ with the congruence $\sigma_{2}$; the hypothesis
$C\cap B_{2}\neq\varnothing$ holds because $S$ is subdirect. Each fibre of $S$
is therefore a single point or the whole of $B_{2}/\sigma_{2}$.

If some fibre is full, then either $S$ is full --- and then the fibre of
$C_{1}/\sigma_{1}$ contains $C_{2}/\sigma_{2}$, giving $C_{1}\cap C_{2}\neq
\varnothing$ --- or the centre of $S$ is proper and nonempty, and
Lemma~\ref{lem:central-relation} produces a proper nonempty binary absorbing or
central subuniverse of $B_{1}/\sigma_{1}$ or $B_{2}/\sigma_{2}$, contradicting
Definition~\ref{def:types}(iii). If every fibre is a single point, then the
fibre of $C_{1}'/\sigma_{1}$ is both $C_{2}/\sigma_{2}$ (witnessed by
$C_{1}'\cap C_{2}$) and $C_{2}'/\sigma_{2}$ (witnessed by $C_{1}'\cap C_{2}'$),
so $C_{2}=C_{2}'$ and $C_{1}\cap C_{2}=C_{1}\cap C_{2}'\neq\varnothing$.
\end{proof}

\subsection{Bridges}\label{sec:bridge-proofs}

Three of the statements below --- Lemmas~\ref{lem:nice-bridge-abelian},
\ref{lem:block-good-bridge} and \ref{lem:nontrivial-bridge} --- require their
bridge to be \emph{pair-reflexive}, and the hypothesis is not cosmetic in any
of them: Remark~\ref{rem:phi-twist} exhibits a bridge satisfying every other
hypothesis of the first and none of its conclusion. Lemma~\ref{lem:symmetrise}
is what makes the hypothesis free, since every bridge these statements are
applied to below is a $\delta^{\circ}$.

\begin{lemma}\label{lem:nice-bridge-abelian}
Let $\sigma$ be a congruence on $\mathbf A$ and $\delta$ a symmetric,
\emph{pair-reflexive} bridge from $\sigma$ to $\sigma$ with
$\proj_{1,2}(\delta)=\bcol\delta=A^{2}$. Then there is an abelian group
$(G;+,-)$ with $(A/\sigma;\delta/\sigma)\cong(G;x_{1}-x_{2}=x_{3}-x_{4})$.
\end{lemma}

\begin{proof}
Factor by $\sigma$ (Lemma~\ref{lem:irred-quotient}) and assume
$\sigma=0_{\mathbf A}$. By Lemma~\ref{lem:bridge-abelian} and
Lemma~\ref{lem:hobby-mckenzie}, $\mathbf A$ is affine, hence polynomially
equivalent to a $\mathbb Z_{p}$-module $\mathbf G$ and Maltsev via $x-y+z$.

View $\delta$ as a binary relation on $\mathbf A^{2}$ as in the proof of
Lemma~\ref{lem:bridge-abelian}. It is reflexive, by pair-reflexivity together
with $\proj_{1,2}(\delta)=A^{2}$, and symmetric by hypothesis. In a Maltsev
algebra a reflexive compatible binary relation is a congruence: if $\eta$ is
reflexive and compatible then for $(x,y),(y,z)\in\eta$ the Maltsev term applied
to $(x,y),(y,y),(y,z)$ gives $(x,z)\in\eta$. So $\delta$ is a congruence of
$\mathbf G^{2}$.

Clause (d) of Definition~\ref{def:bridge} says that a pair
$\bigl((x_{1},x_{2}),(x_{3},x_{4})\bigr)$ of $\delta$ has $x_{1}=x_{2}$ exactly
when $x_{3}=x_{4}$, so the diagonal of $\mathbf G^{2}$ is a $\delta$-block. A
congruence of $\mathbf G^{2}$ with the diagonal as a block is the kernel of
$(x,y)\mapsto x-y$ composed with a quotient of $\mathbf G$; since the diagonal
is exactly one block, no further collapsing occurs, and
$\delta=\{x_{1}-x_{2}=x_{3}-x_{4}\}$. Taking $G=\mathbf G$ finishes.
\end{proof}

\begin{remark}[Pair-reflexivity is not removable]\label{rem:phi-twist}
Without pair-reflexivity the correct conclusion is
$\delta=\{(x_{1},x_{2},x_{3},x_{4}):\varphi(x_{1}-x_{2})=x_{3}-x_{4}\}$ for an
automorphism $\varphi$ of $G$: factoring out the diagonal exhibits $\delta$ as
the preimage of a subgroup $K\le G\times G$ that meets $G\times 0$ and
$0\times G$ trivially and is subdirect, hence a graph. Symmetry says
$\varphi^{2}=\mathrm{id}$; pair-reflexivity says $\varphi=\mathrm{id}$.

Both $\varphi$ occur. On $\mathbf Z_{3}\in\mathcal V_{4}$ --- the operation
$x_{1}+x_{2}+x_{3}+x_{4}$ is idempotent since $4\equiv1$, and is a special WNU
--- take $\sigma=0$ and
\[
  \delta=\{(x_{1},x_{2},x_{3},x_{4})\in\mathbb Z_{3}^{4}:
    x_{1}-x_{2}+x_{3}-x_{4}=0\},
\]
which is $\varphi=-\mathrm{id}$. It is a subuniverse of $\mathbf A^{4}$, being
a subgroup; clause (d) holds because $x_{1}=x_{2}$ iff $x_{3}=x_{4}$;
$\proj_{1,2}(\delta)=\bcol\delta=A^{2}$; and it is symmetric. Yet no bijection
$f\colon\mathbb Z_{3}\to\mathbb Z_{3}$ carries $\delta$ to
$\{x_{1}-x_{2}=x_{3}-x_{4}\}$: such an $f$ would have to satisfy
$f(x_{1})-f(x_{2})=f(x_{3})-f(x_{4})$ whenever $x_{1}-x_{2}+x_{3}-x_{4}=0$, and
putting $(x_{1},x_{2},x_{3},x_{4})=(0,b,b,0)$ gives $2\bigl(f(0)-f(b)\bigr)=0$,
hence $f(0)=f(b)$ for every $b$ since $p$ is odd --- so $f$ is constant, not a
bijection. Since every abelian group of order $3$ is $\mathbb Z_{3}$, no $G$
works, and Lemma~\ref{lem:nice-bridge-abelian} genuinely fails without its
hypothesis.
\end{remark}

\begin{lemma}\label{lem:block-good-bridge}
Let $\delta$ be a bridge from $0_{\mathbf A}$ to $0_{\mathbf A}$ with
$\proj_{1,2}(\delta)\subseteq\bcol\delta$ and $\proj_{1,2}(\delta)$ linked.
Then $\mathbf A$ is BA and centre free.
\end{lemma}

\begin{proof}[Proof outline, with the delicate steps in full]
Replace $\delta$ by $\sigma:=\delta^{\circ}$ of Lemma~\ref{lem:symmetrise} and
put $\omega=\proj_{1,2}(\sigma)=\proj_{1,2}(\delta)$. Note first that
$\bcol\delta$ is \emph{reflexive}: clause (c) of Definition~\ref{def:bridge}
gives $0_{A}\subseteq\proj_{1,2}(\delta)$, and the hypothesis gives
$\proj_{1,2}(\delta)\subseteq\bcol\delta$. Hence $\delta$ is a reflexive
bridge, so Lemma~\ref{lem:symmetrise} applies and gives
$\omega\subseteq\bcol\delta\subseteq\bcol\sigma$ with $\bcol\sigma$
symmetric, which is what legitimises the ``without loss of generality'' below.
Induct on $|A|$.

Suppose $B\subt{T}\mathbf A$ with $T\in\{\TBA,\TC\}$. Since $\omega$ is linked,
$\omega$ meets $(A\setminus B)\times B$ or $B\times(A\setminus B)$; say the
first, and pick $(a,b)$ there.

\emph{Case 1: some subalgebra $D\lneq A$ contains $a$ and $b$.} Put
$\rho=(\omega\circ\omega^{-1})\cap D^{2}$, which is a subalgebra of
$\mathbf D^{2}$, reflexive because $\omega$ is reflexive and symmetric by
construction. Its transitive closure is again a subalgebra, hence a congruence
on $\mathbf D$; let $E$ be the block containing $a$ and $b$, a subuniverse.
Because $\omega$ is reflexive, every undirected $\omega$-edge is a $\rho$-edge
and every $\rho$-edge is a two-edge undirected $\omega$-path, so $\rho$ and the
undirected $\omega$-graph have the same connected components; hence
$\omega\cap E^{2}$ is linked, and $a,b$ lie in one component.
\emph{Note that $(\omega\cup\omega^{-1})\cap D^{2}$ would not serve: a union of
two subuniverses need not be closed.} Moreover
$\proj_{1,2}(\sigma\cap E^{4})=\omega\cap E^{2}$: the inclusion $\subseteq$ is
trivial and $\supseteq$ holds because pair-reflexivity puts
$(x_{1},x_{2},x_{1},x_{2})$ in $\sigma$ for every $(x_{1},x_{2})\in\omega$.
So the inductive hypothesis applies to $\sigma\cap E^{4}$ and makes $E$ BA and
centre free, while Lemma~\ref{lem:ba-center-intersect} makes $B\cap E$ a proper
nonempty subuniverse of $E$ of type $T$ --- proper because $a\in E\setminus B$,
nonempty because $b\in B\cap E$. Contradiction.

\emph{Case 2: no such $D$.} Then $\{a\}\circ\omega=A$, because $\{a\}$ is a
subuniverse by idempotence, hence so is $\{a\}\circ\omega$, and it contains
$a$ and $b$. Put $\xi(x_{1},x_{2})=\sigma(a,x_{1},x_{2},b)$; from
$(a,b,a,b),(a,a,b,b)\in\sigma$ both projections of $\xi$ contain $a$ and $b$,
hence equal $A$. Put $C=B\circ\xi$, which is nonempty because $B$ is and $\proj_{1}(\xi)=A$, so
$C\subte{T}\mathbf A$ by Lemma~\ref{lem:pp-propagation}, with $a\in C$ and
$b\notin C$ --- the latter
because $b\in C$ would give $(a,c,b,b)\in\sigma$ with $c\in B$, and clause (d)
would force $a=c\in B$.

If $T=\TBA$, applying the absorbing term to $(a,b,a,b)$ and $(a,a,b,b)$ gives
$(a,b_{1},b_{2},b)\in\sigma$ with $b_{1},b_{2}\in B$, so $C\cap B\neq
\varnothing$. If $T=\TC$, the three tuples $(a,b,a,b)$, $(a,a,a,a)$ and
$(a,a,b,b)$ witness
\[
  \sigma\cap(\{a\}\times A\times C\times B),\quad
  \sigma\cap(\{a\}\times C\times C\times A),\quad
  \sigma\cap(\{a\}\times C\times A\times B)\;\neq\;\varnothing ,
\]
which is the leave-one-out hypothesis of
Corollary~\ref{cor:no-central-critical} for the ternary relation
$\{(x_{2},x_{3},x_{4}):\sigma(a,x_{2},x_{3},x_{4})\}$ with $n=3$, all types
$\TC$ and all ambients $A$ --- legitimate because $\lll$ is reflexive.
\emph{The three tight boxes of these witnesses --- $\{a\}\times B\times C\times
B$, $\{a\}\times C\times C\times C$, $\{a\}\times C\times B\times B$ --- are
subsets of the displayed ones and do not match the corollary's shape; the
widening is what makes the appeal legitimate.} The corollary's conclusion is
therefore available:
$\sigma\cap(\{a\}\times C\times C\times B)\neq\varnothing$. Put
$C'=\proj_{2}$ of that set, nonempty because the set is. Then
$C'\subte{\TC}A$ by Lemma~\ref{lem:pp-propagation}, using
$\{a\}\circ\omega=A$, and hence
$C'\subte{\TC}C$ by Lemma~\ref{lem:ba-center-intersect}; and by clause (d)
either $a\notin C'$ or $C\cap B\neq\varnothing$.

So in every case we have $\varnothing\neq C'\subt{T}C\subt{T}A$ --- taking
$C'=B\cap C$ when $C\cap B\neq\varnothing$. Then $|C|>1$, and applying the
inductive hypothesis to $\sigma\cap C^{4}$ contradicts $C'\subt{T}C$: its
hypotheses hold because $\{a\}\circ\omega=A$ makes $\omega\cap C^{2}$ a star
centred at $a$, hence linked.
\end{proof}

\begin{longlemma}\label{lem:nontrivial-bridge}
Let $\sigma$ be an irreducible congruence on $\mathbf A$ and $\delta$ a
\emph{pair-reflexive} bridge from $\sigma$ to $\sigma$ with
$\bcol\delta\supsetneq\sigma$. Then
\begin{enumerate}\alphlabels\tightlist
\item $\cover\sigma$ is a congruence;
\item $B/\sigma$ is BA and centre free for every block $B$ of $\cover\sigma$;
\item if $\delta$ is also symmetric then there is a prime $p$ such that every
  block $B$ of $\cover\sigma$ satisfies
  $(B/\sigma;(\delta\cap B^{4})/\sigma)\cong(\mathbb Z_{p}^{n_{B}};
   x_{1}-x_{2}=x_{3}-x_{4})$ for some $n_{B}\ge 0$.
\end{enumerate}
\end{longlemma}

\begin{proof}
By Lemma~\ref{lem:irred-quotient} we may assume $\sigma=0_{\mathbf A}$. Since
$\bcol\delta$ is a subuniverse of $\mathbf A^{2}$ properly containing $\sigma$
and stable under it, minimality gives $\cover\sigma\subseteq\bcol\delta$; the
same argument applied to $\proj_{1,2}(\delta)$, which properly contains $\sigma$
by clause (c), gives $\cover\sigma\subseteq\proj_{1,2}(\delta)$.

Let $B$ be a block of $\LeftLinked(\cover\sigma)$ with $|B|>1$ and
put $\delta'=\delta\cap B^{4}\cap(\cover\sigma\times\cover\sigma)$. Then
$\proj_{1,2}(\delta')=\cover\sigma\cap B^{2}$: for
$(x_{1},x_{2})\in\cover\sigma\cap B^{2}$ pair-reflexivity puts
$(x_{1},x_{2},x_{1},x_{2})$ in $\delta$ --- legitimate because
$\cover\sigma\subseteq\proj_{1,2}(\delta)$ --- and that tuple lies in $B^{4}$
and in $\cover\sigma\times\cover\sigma$ --- \emph{this is where
pair-reflexivity is spent, and without it the projection can collapse to
$0_{B}$}. The resulting relation is linked precisely because $B$ is a block,
and is contained in $\bcol{\delta'}$ because $\cover\sigma\subseteq\bcol\delta$;
so Lemma~\ref{lem:block-good-bridge} makes $B$ BA and centre free. Then
$\cover\sigma\cap B^{2}$, being linked and subdirect on $B$, must be all of
$B^{2}$ by Corollary~\ref{lem:linked-implies-abs}. Blocks of size $1$ satisfy
$B^{2}\subseteq\cover\sigma$ trivially, so
$\cover\sigma=\LeftLinked(\cover\sigma)$ is an equivalence
relation and hence a congruence, giving (a); and (b) follows, one-element
algebras being vacuously BA and centre free.

For (c), fix a block $B$ of $\cover\sigma$ with $|B|\ge2$ and apply
Lemma~\ref{lem:nice-bridge-abelian} to $\delta\cap B^{4}$, whose hypotheses now
hold: $\proj_{1,2}(\delta\cap B^{4})=B^{2}$ by the same pair-reflexivity
computation, and $\bcol{\delta\cap B^{4}}=B^{2}$ because
$B^{2}\subseteq\cover\sigma\subseteq\bcol\delta$. This gives an abelian group
$G_{B}$ with $(B;\delta\cap B^{4})\cong(G_{B};x_{1}-x_{2}=x_{3}-x_{4})$.

It remains to make every $G_{B}$ elementary abelian for one common prime. Put
$\omega=\delta\cap(\cover\sigma\times\cover\sigma)$, which has the same
properties. From $x_{1}-x_{2}=x_{3}-x_{4}$ one pp-defines
$(k+1)x_{1}=k\,x_{2}+x_{3}$ for every $k$, by
\[
  \bigl((k+1)x_{1}=k\,x_{2}+x_{3}\bigr)=\exists x_{4}\;
  \bigl(k\,x_{1}=(k-1)x_{2}+x_{4}\bigr)\wedge(x_{1}-x_{2}=x_{3}-x_{4}),
\]
and hence, identifying $x_{3}$ with $x_{1}$, the relation
$q\,x_{1}=q\,x_{2}$ for every $q$. Let $S$ be what that pp-definition defines
when $x_{1}-x_{2}=x_{3}-x_{4}$ is replaced by $\omega$. If some $G_{B}$ had an
element of order $p_{1}^{2}$, or elements of coprime orders, or if two blocks
had elements of different prime orders, then for a suitable $q$ the relation $S$
would be a subuniverse of $\mathbf A^{2}$ with $S\supsetneq\sigma$ and
$\cover\sigma\not\subseteq S$, contradicting the minimality of $\cover\sigma$.
Hence every element of every $G_{B}$ has one and the same prime order $p$, and
$G_{B}\cong\mathbb Z_{p}^{n_{B}}$.
\end{proof}

\begin{hazard}[It is minimality, not irreducibility]\label{haz:minimality-not-irred}
The contradiction at the end of the proof is with the \emph{minimality of
$\cover\sigma$}, which is a consequence of irreducibility and not the same
statement. The distinction matters wherever $\cover\sigma$ is treated as data
attached to a proof of irreducibility (Caveat~\ref{haz:irreducible}).
\end{hazard}

\begin{proof}[Proof of Lemma~\ref{lem:linear-equiv}]
$(a)\Rightarrow(b)$: take $\delta=S$ of Definition~\ref{def:linear-congruence},
which by Lemma~\ref{lem:linear-bridge} is a bridge with
$\bcol S=\cover\sigma\supsetneq\sigma$.

$(b)\Rightarrow(a)$: given a bridge $\delta$ with
$\bcol\delta\supsetneq\sigma$ --- reflexive by
Lemma~\ref{lem:bridge-reflexive}, since the hypothesis contains
$\bcol\delta\supseteq\sigma$ --- form $\delta^{\circ}$ of
Lemma~\ref{lem:symmetrise}; it is symmetric and pair-reflexive, and
$\bcol{\delta^{\circ}}\supseteq\bcol\delta\supsetneq\sigma$. Now
Lemma~\ref{lem:nontrivial-bridge} gives items (b) and (c) of
Definition~\ref{def:linear-congruence}, with the witness
$S=\delta^{\circ}\cap(\cover\sigma\times\cover\sigma)$.
\end{proof}



\begin{longlemma}\label{lem:pc-inductive-step}
Let $\sigma$ be a congruence on $\mathbf A$ and $\delta$ a reflexive bridge
from $\sigma$ to $\sigma$ with
\begin{enumerate}\romanlabels\tightlist
\item $\delta(x_{1},x_{2},x_{3},x_{4})=\delta(x_{3},x_{4},x_{1},x_{2})$;
\item $(a,b,a,b),(b,a,b,a)\in\delta$ for every $(a,b)\in\proj_{1,2}(\delta)$;
\item $\RightLinked(\proj_{1,2,3}(\delta))=A^{2}$.
\end{enumerate}
Then $(a,a,b,b)\in\delta$ for some $a\neq b$.
\end{longlemma}

\begin{proof}
Factor by $\sigma$ and assume $\sigma=0_{\mathbf A}$; clause (c) of
Definition~\ref{def:bridge} then forces $|A|\ge2$. Write
$P=\proj_{1,2}(\delta)$, which is symmetric by (ii), and
$W=\proj_{1,2,3}(\delta)\sdle\mathbf P\times\mathbf A$. Induct on $|A|$.

\emph{Case 1: some $B\subt{T}\mathbf A$ with $|B|>1$, $T\in\{\TBA,\TC\}$.}
Put $W'=\proj_{1,2,3}(\delta\cap B^{4})$ and write
$W(x_{1},x_{2},x_{3})=\exists x_{4}\,\delta(x_{1},x_{2},x_{3},x_{4})\wedge
A(x_{1})\wedge A(x_{2})\wedge A(x_{3})\wedge A(x_{4})$ with $A$ the full unary
relation. Replacing $A$ by $B$ --- a \emph{single} application of
Lemma~\ref{lem:pp-propagation}, which replaces every occurrence, so no
transitivity of $\subte{\TC}$ is needed --- gives $W'\subte{T}W$. It is
nonempty, containing $(b,b,b)$ for every $b\in B$, and its projections are
$\proj_{1,2}(W')=P\cap B^{2}$, by (ii), and $\proj_{3}(W')=B$. So
Lemma~\ref{lem:absorb-linked-sub} makes $W'$ linked, which is (iii) for
$\delta\cap B^{4}$. Conditions (i) and (ii) are inherited and $\delta\cap B^{4}$
is reflexive. If it is a bridge, the inductive hypothesis applies and produces
the required pair inside $B$. If it is not, then $\proj_{1,2}(\delta\cap
B^{4})=0_{B}$, so by clause (d) every tuple of $\delta\cap B^{4}$ has the form
$(x,x,y,y)$; linkedness of $W'$ on $|B|\ge2$ points then forbids the perfect
matching, so some $(a,a,b,b)\in\delta$ with $a\neq b$.

\emph{Case 2: some $\{a\}\subt{T}\mathbf A$, $T\in\{\TBA,\TC\}$, and no such $B$
of size $>1$.} By (i) the relation $\proj_{1,3}(\delta)$ is symmetric, and by
(iii) it is linked, so $\{a\}\circ\proj_{1,3}(\delta)\neq\{a\}$; pick
$(a,b,c,d)\in\delta$ with $c\neq a$. If $a=b$ then $c=d$ by clause (d) and we
are done. If $\{a\}\circ\proj_{1,3}(\delta)\neq A$ then that set is a strong
subuniverse containing $a$ and $c$, so of size $>1$, which is Case 1.
Otherwise choose $(a,e,b,f)\in\delta$. By (ii), $(b,a,b,a)\in\delta$. Let $g$
be a ternary absorbing operation for $\{a\}$, available for $T=\TC$ by
Lemma~\ref{lem:central-ternary}. Applying $g$
to $(a,e,b,f)$, $(b,a,b,a)$ and $(a,a,a,a)$ gives $(a,a,g(b,b,a),a)\in\delta$,
whence $g(b,b,a)=a$ by clause (d); applying $g$ to $(b,a,b,a)$, $(b,a,b,a)$ and
$(a,e,b,f)$ then gives $(a,a,b,a)\in\delta$, whence $b=a$, a contradiction.

\emph{Case 3: $\mathbf A$ is BA and centre free.} Put
$C=\{(a,b)\in P:\{(a,b)\}\times A\subseteq W\}$. By
Lemma~\ref{lem:bacon-left}, applied to the transpose of $W$,
$C\subte{\TBA}\mathbf P$. If $C=P$ then in particular $(a,a)\in C$ for any $a$;
otherwise $C\subt{\TBA}P$ and Lemma~\ref{lem:absorbing-equality} gives
$C\cap 0_{\mathbf A}\neq\varnothing$. Either way $(a,a)\in C$ for some $a$, so
for every $c$ there is $d$ with $\delta(a,a,c,d)$, and clause (d) forces $d=c$.
Taking $c\neq a$ finishes the proof.
\end{proof}

\begin{lemma}[{\cite[Lem.~7.2]{ZhukJACM}}]\label{lem:absorbing-equality}
If $0_{\mathbf A}\subseteq\sigma\le\mathbf A^{2}$ and
$\omega\subt{\TBA}\sigma$, then $\omega\cap 0_{\mathbf A}\neq\varnothing$.
\end{lemma}

\begin{proof}[Proof of Lemma~\ref{lem:pc-bridges-trivial}]
Put $\xi=\delta\ast\delta^{-1}$, which is symmetric and pair-reflexive. Each of
$\RightLinked(\proj_{1,2,3}(\xi))$ and
$\RightLinked(\proj_{1,2,4}(\xi))$ is a congruence containing
$\sigma$, so by minimality of $\cover\sigma$ each is $\sigma$ or contains
$\cover\sigma$; this gives three cases.

If both equal $\sigma$, then for $(a,b,c,d)\in\xi$ the classes $c/\sigma$ and
$d/\sigma$ are determined by $a/\sigma,b/\sigma$; since $(a,b,a,b)\in\xi$ we get
$(a,c),(b,d)\in\sigma$, so $\xi$ is trivial. \emph{Transferring that
determinacy from $\xi=\delta\ast\delta^{-1}$ to $\delta$ itself is a counting
argument, and it is where the hypothesis
$\proj_{1,2}(\delta)=\proj_{3,4}(\delta)$ is spent:} on
$\cover\sigma/\sigma$, which is finite, $\delta$ induces a bipartite relation
whose left neighbourhoods are nonempty, pairwise disjoint --- that is what the
case gives --- and cover the right side; equal finite cardinalities force every
neighbourhood to be a singleton, so the induced relation is a bijection and
determinacy holds in both directions.

With that, introduce
\begin{align*}
  \zeta_{1}(x_{1},x_{2},x_{3},x_{4})=\;&
    \exists y\exists y'\exists z\exists z'
    \exists t_{1}\exists t_{2}\exists t_{3}\exists t_{4}\;\\
    &\delta(x_{1},y,z,t_{1})\wedge\delta(x_{2},y,z',t_{2})\wedge
     \delta(x_{3},y',z,t_{3})\wedge\delta(x_{4},y',z',t_{4}),\\
  \zeta_{2}(x_{1},x_{2},x_{3},x_{4})=\;&
    \exists y\exists z\;\;\delta(x_{1},y,z,x_{3})\wedge\delta(x_{2},y,z,x_{4}),
\end{align*}
and pick $(a,b,c,d)\in\delta$ with $(a,b)\in\cover\sigma\setminus\sigma$. From
$(a,b,c,d),(b,b,b,b)\in\delta$ one gets $(a,b,a,b)\in\zeta_{1}$ and hence
$\proj_{1,2}(\zeta_{1})\supseteq\cover\sigma$. Split.

\emph{If $\zeta_{1}$ is a bridge}, then $\bcol{\zeta_{1}}=\sigma$ --- forced,
since $\sigma$ is PC --- and $(a,a,c,c)\in\zeta_{1}$ for every
$(a,b,c,d)\in\delta$, so $\proj_{1,3}(\delta)=\sigma$.

\emph{If $\zeta_{1}$ is not a bridge}, there is $(a,a,b,c)\in\zeta_{1}$ with
$(b,c)\notin\sigma$. Take witnesses $y=d$, $y'=d'$, $z=e$, $z'=e'$,
$t_{i}=f_{i}$. Determinacy makes $(e,e')\in\sigma$, so
$(b,d',e,f_{3}),(c,d',e,f_{4})\in\delta$ and therefore
$\proj_{1,2}(\zeta_{2})\supsetneq\sigma$; determinacy in both directions then
makes $\zeta_{2}$ a bridge, whence $\bcol{\zeta_{2}}=\sigma$ and
$\proj_{1,4}(\delta)=\sigma$.

So $\proj_{1,3}(\delta)=\sigma$ or $\proj_{1,4}(\delta)=\sigma$, and by
the same argument with $x_{1},x_{2}$ interchanged, $\proj_{2,4}(\delta)=\sigma$
or $\proj_{2,3}(\delta)=\sigma$. \emph{Two of the four combinations must now be
excluded:} $\proj_{1,3}=\proj_{2,3}=\sigma$ would
give $\proj_{1,2}(\delta)\subseteq\sigma$, against
$\proj_{1,2}(\delta)=\cover\sigma\supsetneq\sigma$, and likewise for
$\proj_{1,4}=\proj_{2,4}$. The two surviving combinations are the conclusion,
the reverse inclusion following from stability under $\sigma$.

If $\RightLinked(\proj_{1,2,3}(\xi))\supseteq\cover\sigma$, choose
a block $B$ of it that is not a $\sigma$-block. Then $a,b,d\in B$ whenever
$c\in B$ and $(a,b,c,d)\in\xi$, so $\xi\cap(A^{2}\times B\times A)\subseteq
B^{4}$ and therefore $\proj_{1,2,3}(\xi\cap B^{4})=\proj_{1,2,3}(\xi)\cap
(A^{2}\times B)$, which is linked. \emph{If $\cover\sigma\cap B^{2}\supsetneq
\sigma\cap B^{2}$ then $\xi\cap B^{4}$ is a bridge and
Lemma~\ref{lem:pc-inductive-step} applies, contradicting PC-ness via
Lemma~\ref{lem:linear-equiv}. If not, the lemma cannot be cited --- it is not a
bridge --- and the case must be inlined: clause (d) then puts both pairs of
every tuple in $\sigma$, and linkedness on a $B$ that is not a $\sigma$-block
produces $(x_{1},x_{1},x_{3},x_{3})\in\xi$ with $(x_{1},x_{3})\notin\sigma$
directly.} The third case is the mirror image, using $\proj_{1,2,4}$.
\end{proof}

\begin{proof}[Proof of Lemma~\ref{lem:no-cross-bridge}]
Let $\sigma_{1},\sigma_{2}$ be congruences on $\mathbf A_{1},\mathbf A_{2}$;
both are irreducible, $\sigma_{1}$ because it is linear
(Definition~\ref{def:linear-congruence}(a)). Replace $\delta$ first by
$\delta\cap(\cover{\sigma_{1}}\times\cover{\sigma_{2}})$, which is again a
bridge with the same collapse, and then by the $\omega$ of
Lemma~\ref{lem:rectangularise}. By clause (d) of that lemma the containment in
$\cover{\sigma_{1}}\times\cover{\sigma_{2}}$ survives, and by clause (c)
neither linking equivalence changes, so the case analysis below is unaffected;
by clause (b), $\bcol\delta$ is now rectangular. Being a bridge inside
$\cover{\sigma_{1}}\times\cover{\sigma_{2}}$, it has
$\proj_{3,4}(\delta)\supsetneq\sigma_{2}$ and stable under $\sigma_{2}$, so
minimality of the cover (Proposition~\ref{prop:cover}) gives
\begin{equation}\label{eq:ncb-proj}
  \proj_{3,4}(\delta)=\cover{\sigma_{2}},
\end{equation}
and symmetrically $\proj_{1,2}(\delta)=\cover{\sigma_{1}}$.

\emph{The easy case.} Suppose $\RightLinked(\bcol\delta)\supsetneq\sigma_{2}$.
The composite $\delta^{-1}\ast\delta$ is a bridge from $\sigma_{2}$ to
$\sigma_{2}$ by Lemma~\ref{lem:bridge-comp} --- both congruences are
irreducible --- with collapse $\bcol\delta^{-1}\circ\bcol\delta$. Iterating,
the $k$-fold composite has collapse $(\bcol\delta^{-1}\circ\bcol\delta)^{k}$,
and for $k$ large this is $\RightLinked(\bcol\delta)\supsetneq\sigma_{2}$. So
Lemma~\ref{lem:linear-equiv} makes $\sigma_{2}$ linear.

So assume $\RightLinked(\bcol\delta)=\sigma_{2}$. Then
$\LeftLinked(\bcol\delta)\supseteq\cover{\sigma_{1}}$: otherwise minimality of
$\cover{\sigma_{1}}$ forces $\LeftLinked(\bcol\delta)=\sigma_{1}$, and
$\bcol\delta$ then induces a bijection
$\mathbf A_{1}/\sigma_{1}\cong\mathbf A_{2}/\sigma_{2}$ carrying
$\cover{\sigma_{1}}/\sigma_{1}$ to $\cover{\sigma_{2}}/\sigma_{2}$, so
$\sigma_{2}$ is linear because $\sigma_{1}$ is.

\emph{Normalising the first and third coordinates.} Put
$\delta'(x_{1},x_{2},x_{3},x_{4})=\delta(x_{1},x_{2},x_{3},x_{4})\wedge
\bcol\delta(x_{1},x_{3})$.

If $\delta'$ is a bridge, it has $\proj_{1,3}(\delta')=\bcol{\delta'}$ by
construction, and we take $\omega=\delta'$.

If $\delta'$ is not a bridge, then --- clauses (a), (b), (d) being inherited ---
clause (c) fails, so there is no $(a,b,c,d)\in\delta$ with
$(a,b)\notin\sigma_{1}$ and $(a,c)\in\bcol\delta$. Put
\[
  \delta''(x_{1},x_{2},x_{3},x_{4})=\exists z\;
    \bcol\delta(x_{1},x_{3})\wedge\delta(x_{2},z,x_{3},x_{4}).
\]
This is a bridge. For clause (d) in one direction: if
$(x_{1},x_{2})\in\sigma_{1}$ then $(x_{2},x_{3})\in\bcol\delta$, so the
failure above applied to $(x_{2},z,x_{3},x_{4})\in\delta$ forces
$(x_{2},z)\in\sigma_{1}$ and hence $(x_{3},x_{4})\in\sigma_{2}$. In the other
direction, if $(x_{3},x_{4})\in\sigma_{2}$ then $(x_{2},z)\in\sigma_{1}$ and
$(x_{2},x_{3})\in\bcol\delta$, whence $(x_{1},x_{2})\in\sigma_{1}$. For
clause (c): since $\delta'\neq\delta$ there is $(a,b,c,d)\in\delta$ with
$(a,c)\notin\bcol\delta$. Some $e$ has $(e,c)\in\bcol\delta$: by
\eqref{eq:ncb-proj} the reflexive pair $(c,c)\in\sigma_{2}\subseteq
\cover{\sigma_{2}}$ is $\proj_{3,4}$ of some $(e,e',c,c)\in\delta$; clause (d)
of Definition~\ref{def:bridge} puts $(e,e')\in\sigma_{1}$, and stability of the
first two variables under $\sigma_{1}$ --- clause (a) --- replaces $e'$ by $e$,
giving $(e,e,c,c)\in\delta$. That $e$ gives $(e,a,c,d)\in\delta''$ with
$(e,a)\notin\sigma_{1}$. Take $\omega=\delta''$.

Either way we have a bridge $\omega$ from $\sigma_{1}$ to $\sigma_{2}$ with
$\proj_{1,3}(\omega)=\bcol\omega=\bcol\delta$.

\emph{The contradiction.} Suppose $\sigma_{2}$ is \emph{not} linear, so it is
PC. Put
\[
  \xi_{1}(x_{1},x_{2},x_{3},x_{4})=\exists x_{5}\exists x_{6}\;
    \omega(x_{5},x_{6},x_{1},x_{2})\wedge\omega(x_{5},x_{6},x_{3},x_{4}),
\]
a bridge from $\sigma_{2}$ to $\sigma_{2}$; by
Lemma~\ref{lem:pc-bridges-trivial} it is
$\sigma_{2}(x_{1},x_{3})\wedge\sigma_{2}(x_{2},x_{4})$ --- the other
alternative is excluded by symmetry of $\xi_{1}$ in its two pairs. Put
\[
  \xi_{2}(x_{1},x_{2},x_{3},x_{4})=\exists x_{5}\exists x_{6}\;
    \omega(x_{5},x_{6},x_{1},x_{2})\wedge\omega(x_{6},x_{5},x_{3},x_{4}).
\]
Here $(x_{5},x_{1}),(x_{6},x_{3})\in\bcol\omega$ and
$(x_{5},x_{6})\in\cover{\sigma_{1}}\subseteq\LeftLinked(\bcol\omega)$, which
gives $(x_{1},x_{3})\in\sigma_{2}$; so $\xi_{2}$ too is
$\sigma_{2}(x_{1},x_{3})\wedge\sigma_{2}(x_{2},x_{4})$ by
Lemma~\ref{lem:pc-bridges-trivial}. Comparing the two,
$(b,a,c,d)\in\omega$ whenever $(a,b,c,d)\in\omega$.

Now put
\[
  \zeta(x_{1},x_{2},x_{3},x_{4})=\exists y_{1}\exists y_{2}\exists y_{3}\;
    \omega(y_{1},y_{2},x_{1},x_{2})\wedge\omega(y_{1},y_{3},x_{1},x_{3})
    \wedge\omega(y_{2},y_{3},x_{1},x_{4}).
\]
If $x_{1}=x_{2}$ then $y_{1}=y_{2}$ by the shape of $\xi_{1}$, and then
$x_{3}=x_{4}$. Take any $(c,d)\in\cover{\sigma_{2}}$; since
$\proj_{3,4}(\omega)=\cover{\sigma_{2}}$ there is
$(a,b)\in\cover{\sigma_{1}}$ with $(a,b,c,d)\in\omega$. Now
$(a,c)\in\bcol\omega$ --- which \emph{is} $(a,a,c,c)\in\omega$, by the
definition of the collapse --- and
$(a,b)\in\cover{\sigma_{1}}\subseteq\LeftLinked(\bcol\omega)$, so
rectangularity in the form of Lemma~\ref{lem:rectangularise}(b) gives
$(b,c)\in\bcol\omega$, that is $(b,b,c,c)\in\omega$. \emph{This is the one
step of the proof that consumes rectangularity.} Sending
$(x_{1},x_{2},x_{3},x_{4})\mapsto(c,d,c,d)$ with $(y_{1},y_{2},y_{3})=(a,b,a)$
witnesses $(c,d,c,d)\in\zeta$; sending it to $(c,c,d,d)$ with
$(y_{1},y_{2},y_{3})=(a,a,b)$ witnesses $(c,c,d,d)\in\zeta$. So
$\zeta(x_{1},x_{2},x_{3},x_{4})\wedge\zeta(x_{3},x_{4},x_{1},x_{2})$ is a
bridge from $\sigma_{2}$ to $\sigma_{2}$ whose collapse contains
$\cover{\sigma_{2}}\supsetneq\sigma_{2}$, so $\sigma_{2}$ is linear by
Lemma~\ref{lem:linear-equiv} --- contradicting the supposition.
\end{proof}

\begin{hazard}[Where rectangularity is spent]\label{haz:rectangularise}
The proof opens by normalising twice, and both steps were asserted in an
earlier draft rather than argued.

\emph{Rectangularity of $\bcol\delta$} is now
Lemma~\ref{lem:rectangularise}: composing $\delta$ with itself enough times
replaces $\bcol\delta$ by $\LeftLinked(\bcol\delta)\circ\bcol\delta$, which is
a union of full boxes, one per connected component, and leaves both linking
equivalences and the containment in
$\cover{\sigma_{1}}\times\cover{\sigma_{2}}$ alone. It is worth knowing that
the normalisation is spent \emph{once}, at $(b,b,c,c)\in\omega$ in the
construction of $\zeta$; nothing else in the proof uses it. Note also that this
is not the situation of Caveat~\ref{haz:reflexivise}, where the property turned
out to hold already: a genuine replacement of $\delta$ is being made here, and
the lemma is what licenses it.

\emph{Choosing $e$ with $(e,c)\in\bcol\delta$} is the second, and it is
three lines from $\proj_{3,4}(\delta)=\cover{\sigma_{2}}$, clause (d) of
Definition~\ref{def:bridge} and stability under $\sigma_{1}$. What made it look
harder than it is was the earlier draft's appeal to $\proj_{1}(\bcol\delta)$,
which is the wrong projection: what is needed is that $c$ lies in
$\proj_{2}(\bcol\delta)$, and the reflexive pair $(c,c)$ supplies it.
\end{hazard}

\begin{proof}[Proof of Lemma~\ref{lem:bridge-to-pc}]
(a) is Lemma~\ref{lem:no-cross-bridge} contrapositively: if $\sigma_{2}$ were
linear then so would $\sigma_{1}$ be, via $\delta^{-1}$.

For (b) and (c), $\delta^{-1}\ast\delta$ is a bridge from $\sigma_{1}$ to
$\sigma_{1}$, and $\sigma_{1}$ is PC, so by
Lemma~\ref{lem:pc-bridges-trivial} it is trivial:
$\sigma_{1}(x_{1},x_{3})\wedge\sigma_{1}(x_{2},x_{4})$ or
$\sigma_{1}(x_{1},x_{4})\wedge\sigma_{1}(x_{2},x_{3})$. Its collapse is
$\bcol\delta\circ\bcol\delta^{-1}$, which is therefore $\sigma_{1}$; so the
relation $\{(a/\sigma_{1},b/\sigma_{2}):(a,b)\in\bcol\delta\}$ has
left-neighbourhoods that are pairwise disjoint. Running the same argument on
$\delta\ast\delta^{-1}$ makes the right-neighbourhoods pairwise disjoint. Both
sides are covered, since $\proj_{1,2}(\delta)=\cover{\sigma_{1}}$ and
$\proj_{3,4}(\delta)=\cover{\sigma_{2}}$, so the induced relation is a
bijection, which is (c), and it is an isomorphism of algebras because
$\bcol\delta$ is a subuniverse --- that is (b).

For (d), put
\[
  \xi(x_{1},x_{2},x_{5},x_{6})=\exists x_{3}\exists x_{4}\;
    \delta(x_{1},x_{2},x_{3},x_{4})\wedge\bcol\delta(x_{5},x_{3})\wedge
    \bcol\delta(x_{6},x_{4}),
\]
a bridge from $\sigma_{1}$ to $\sigma_{1}$ by (c). Lemma~\ref{lem:pc-bridges-trivial}
gives $\xi=\sigma_{1}(x_{1},x_{5})\wedge\sigma_{1}(x_{2},x_{6})$ or
$\sigma_{1}(x_{1},x_{6})\wedge\sigma_{1}(x_{2},x_{5})$, and transporting each
alternative through the bijection of (c) gives the two alternatives of (d).
\end{proof}

\begin{hazard}[The third case]\label{haz:pc-bridges-hedges}
The third case is genuinely the mirror of the second: $(a,b,a,b)\in\xi$ places
$b$ in the block, and $\cover\sigma\subseteq\RightLinked$ places
$a$ and $c$. Writing it out is mechanical, and it is not written out here.
\end{hazard}

\subsection{Types interaction}\label{sec:types-interaction}

\begin{lemma}[{\cite[Lem.~8.19]{ZhukJACM}}]\label{lem:bridge-between-congruences}
Let $\omega,\sigma_{1},\sigma_{2}$ be congruences on $\mathbf A$ with
$\omega\cap\sigma_{1}=\omega\cap\sigma_{2}$ and
$\omega\setminus\sigma_{1}\neq\varnothing$. Then there is a bridge $\delta$
from $\sigma_{1}$ to $\sigma_{2}$ with $\bcol\delta=\sigma_{1}\circ\sigma_{2}$.
\end{lemma}

\begin{proof}
Put
$\delta(x_{1},x_{2},y_{1},y_{2})=\exists z_{1}\exists z_{2}\;
\sigma_{1}(x_{1},z_{1})\wedge\sigma_{1}(x_{2},z_{2})\wedge\omega(z_{1},z_{2})
\wedge\sigma_{2}(z_{1},y_{1})\wedge\sigma_{2}(z_{2},y_{2})$.
Clauses (a) and (b) are visible. For (d): if $(x_{1},x_{2})\in\sigma_{1}$ then
$(z_{1},z_{2})\in\sigma_{1}\cap\omega=\sigma_{2}\cap\omega$, so
$(y_{1},y_{2})\in\sigma_{2}$; \emph{and conversely}, if
$(y_{1},y_{2})\in\sigma_{2}$ then $(z_{1},z_{2})\in\omega\cap\sigma_{2}=
\omega\cap\sigma_{1}$, so $(x_{1},x_{2})\in\sigma_{1}$ --- the direction the
source omits, and the one that spends the other half of the hypothesis. For
(c): $\sigma_{1}\subseteq\proj_{1,2}(\delta)$ by taking $z_{i}=y_{i}=x_{1}$,
and choosing $(a,b)\in\omega\setminus\sigma_{1}$ gives $(a,b,a,b)\in\delta$
with $(a,b)\notin\sigma_{1}$ and, since
$\omega\setminus\sigma_{1}=\omega\setminus\sigma_{2}$, also
$(a,b)\notin\sigma_{2}$. Finally $\bcol\delta=\sigma_{1}\circ\sigma_{2}$: the
inclusion $\supseteq$ by taking $z_{1}=z_{2}$, using reflexivity of $\omega$,
and $\subseteq$ by reading off $(x,z_{1})\in\sigma_{1}$, $(z_{1},y)\in\sigma_{2}$.
\end{proof}

\begin{longlemma}\label{lem:two-stable}
Let $C_{1}\subt[A]{T_{1}(\sigma_{1})}B_{1}\lll A$ and
$C_{2}\subt[A]{T_{2}(\sigma_{2})}B_{2}\lll A$ with
$T_{1},T_{2}\in\{\TBA,\TC,\TS,\TD\}$, $C_{1}\cap B_{2}\neq\varnothing$,
$B_{1}\cap C_{2}\neq\varnothing$ and $C_{1}\cap C_{2}=\varnothing$. Then
$T_{1}=T_{2}\in\{\TBA,\TC,\TD\}$; and if $T_{1}=T_{2}=\TD$ there is a bridge
$\delta$ from $\sigma_{1}$ to $\sigma_{2}$ with
$\bcol\delta=\sigma_{1}\circ\sigma_{2}$.
\end{longlemma}

\begin{proof}
\emph{No $\TS$.} If $T_{1}=\TS$, pick $S_{1}\subseteq C_{1}$ with
$S_{1}\subt{\TBA,\TC}B_{1}$. Since $C_{2}\lll A$ --- the chain for $B_{2}$
extended by one $\TD$ step --- and $B_{1}\cap C_{2}\neq\varnothing$,
Lemma~\ref{lem:intersection-good}(s) gives $S_{1}\cap C_{2}\neq\varnothing$,
hence $C_{1}\cap C_{2}\neq\varnothing$. Symmetrically for $T_{2}$.

\emph{Both absorbing.} If $T_{1},T_{2}\in\{\TBA,\TC\}$ then by
Lemma~\ref{lem:ba-center-intersect} $C_{1}\cap B_{2}\subte{T_{1}}B_{1}\cap
B_{2}$ and $B_{1}\cap C_{2}\subte{T_{2}}B_{1}\cap B_{2}$; both are proper,
since $C_{1}\cap B_{2}=B_{1}\cap B_{2}$ would put $B_{1}\cap C_{2}$ inside
$C_{1}\cap C_{2}=\varnothing$. Their intersection is empty, so
Lemma~\ref{lem:possible-intersections} gives $T_{1}=T_{2}\in\{\TBA,\TC\}$.

\emph{Mixed.} Let $T_{1}\in\{\TBA,\TC\}$ and $T_{2}=\TD$; the mirror case
$T_{1}=\TD$, $T_{2}\in\{\TBA,\TC\}$ is identical with the indices swapped, the
hypotheses being symmetric. By
Lemma~\ref{lem:intersection-good}(d), $(B_{1}\cap B_{2})/\sigma_{2}$ has size
$1$ or is all of $B_{2}/\sigma_{2}$. In the first case the single block is the
one defining $C_{2}$, because $B_{1}\cap C_{2}\neq\varnothing$; then
$B_{1}\cap B_{2}=B_{1}\cap C_{2}$, so $C_{1}\cap B_{2}\subseteq C_{1}\cap
C_{2}=\varnothing$, a contradiction. In the second,
Lemmas~\ref{lem:ba-center-intersect} and \ref{lem:factor-type} give
$(C_{1}\cap B_{2})/\sigma_{2}\subte{T_{1}}B_{2}/\sigma_{2}$, and this is
\emph{proper}: equality would make the $\sigma_{2}$-block defining $C_{2}$ meet
$C_{1}$, i.e.\ $C_{1}\cap C_{2}\neq\varnothing$. So $B_{2}/\sigma_{2}$ has a
proper nonempty strong subuniverse, contradicting
Definition~\ref{def:types}(iii).

\emph{Both $\TD$.} Let $\omega$ be the intersection of the dividing congruences
of the chosen chains for $B_{1}\lll A$ and $B_{2}\lll A$. Apply
Lemma~\ref{lem:intersection-good}(d) to $B_{1}$ and $C_{2}$ with the congruence
$\sigma_{1}$ --- \emph{to $C_{2}$, not to $B_{2}$; against $B_{2}$ the step does
not follow.} Its second alternative, $(B_{1}\cap C_{2})/\sigma_{1}=B_{1}/
\sigma_{1}$, would put a point of $C_{2}$ in the $\sigma_{1}$-block defining
$C_{1}$; so $|(B_{1}\cap C_{2})/\sigma_{1}|=1$ and the ``moreover'' clause
gives $\sigma_{1}\supseteq\omega\cap\sigma_{2}$, because the chain witnessing
$C_{2}\lll A$ is the chain for $B_{2}$ extended by the step
$C_{2}\subt{\TD(\sigma_{2})}B_{2}$ (Caveat~\ref{haz:lll-chain}). Symmetrically
$\sigma_{2}\supseteq\omega\cap\sigma_{1}$, so
$\omega\cap\sigma_{1}=\omega\cap\sigma_{2}$. Finally, for $c_{1}\in C_{1}\cap
B_{2}$ and $c_{2}\in B_{1}\cap C_{2}$ we have $(c_{1},c_{2})\in\omega$, since
both lie in $B_{1}\cap B_{2}$ and each $B_{i}$ sits inside a single block of
every dividing congruence of its own chain; and $(c_{1},c_{2})\notin\sigma_{1}$,
since $c_{2}\in B_{1}$ and $c_{2}\,\sigma_{1}\,c_{1}\in C_{1}$ would give
$c_{2}\in C_{1}\cap C_{2}$. Now Lemma~\ref{lem:bridge-between-congruences}
applies.
\end{proof}

\subsection{Factorisation}\label{sec:factorisation}

\begin{lemma}\label{lem:center-pushed-in}
Let $R\sdle\mathbf A_{1}\times\mathbf A_{2}$ with
$R\cap(B_{1}\times B_{2})\neq\varnothing$, $B_{i}\lll A_{i}$, $\sigma$ a
congruence on $\mathbf A_{1}$ with $B_{1}/\sigma$ BA and centre free. If some
$c\in A_{2}$ has $(E\times\{c\})\cap R\neq\varnothing$ for every $E\in
B_{1}/\sigma$, then such a $c$ may be chosen in $B_{2}$.
\end{lemma}

\begin{proof}[Proof outline]
If $c$ may already be chosen in $B_{2}$ there is nothing to do; otherwise walk
down a chain for $B_{2}\lll A_{2}$ and let $B_{2}''\subt{T(\delta)}B_{2}'$ be
the first step at which the choice fails, so $B_{2}\subseteq B_{2}''$. Put
$S'$ and $S''$ for the sets of $\sigma$-class tuples of length $|A_{1}|$
witnessed by some $c$ in $B_{2}'$, resp.\ $B_{2}''$. A single good $c\in B_{2}'$
makes $S'=(B_{1}/\sigma)^{|A_{1}|}$, and $S''\neq\varnothing$ because
$R\cap(B_{1}\times B_{2})\neq\varnothing$ and $B_{2}\subseteq B_{2}''$ ---
\emph{this is where that hypothesis is spent}. For $T\in\{\TBA,\TC\}$,
Lemmas~\ref{lem:pp-propagation} and \ref{lem:power-to-base} then put a strong
subuniverse on $B_{1}/\sigma$, a contradiction; for $T=\TS$ the same follows
once $R\cap(B_{1}\times D_{2})\neq\varnothing$, which
Lemma~\ref{lem:intersection-good}(s) supplies inside $\mathbf R$.

For $T=\TD$ one builds
$S\le(B_{1}/\sigma)^{|A_{1}|}\times B_{2}'/\delta$ and finds a full column over
$d=c/\delta$, while $(B_{1}/\sigma)^{|A_{1}|}\times B_{2}''/\delta\not\subseteq
S$ by minimality. Corollary~\ref{cor:intersection-good} then computes
$\proj_{|A_{1}|+1}(S)$. It offers three alternatives and only the third gives
the contradiction. ``Empty'' is excluded by
$R\cap(B_{1}\times B_{2})\neq\varnothing$. ``Size one'' is excluded thus: the
single $\delta$-class would have to be $[c]_{\delta}$, since $c$ has an
$R$-partner in $B_{1}$; and $B_{2}''=B_{2}'\cap E$ for a $\delta$-block $E$, so
either $E=[c]_{\delta}$, putting $c$ in $B_{2}''$ against the minimal choice of
$B_{2}'$, or $E\cap[c]_{\delta}=\varnothing$, making
$R\cap(B_{1}\times B_{2}'')=\varnothing$ against $B_{2}\subseteq B_{2}''$. So
$S$ is a central relation, and Lemmas~\ref{lem:central-relation} and
\ref{lem:power-to-base} contradict the hypotheses.
\end{proof}

\begin{lemma}\label{lem:leftlinked-stay-full}
Under the hypotheses of Lemma~\ref{lem:center-pushed-in}, if
$(\LeftLinked(R)\cap B_{1}^{2})/\sigma=(B_{1}/\sigma)^{2}$ then
$\LeftLinked\bigl(R\cap(B_{1}\times B_{2})\bigr)/\sigma=
(B_{1}/\sigma)^{2}$.
\end{lemma}

\begin{proof}[Proof outline]
Write $R_{n}=(R\circ R^{-1})^{n}$, which is reflexive and symmetric because $R$
is subdirect, so $R_{2n}=R_{n}\circ R_{n}^{-1}$ --- \emph{this is why the
source restricts to $n=2^{k}$}. If $(R_{1}\cap B_{1}^{2})/\sigma$ is already
full, Lemma~\ref{lem:bacon-left} applied to
$S=\{(a/\sigma,b):a\in B_{1},(a,b)\in R\}$, subdirect over
$(B_{1}/\sigma)\times\proj_{2}(S)$ rather than over $(B_{1}/\sigma)\times A_{2}$,
produces the required $c$, which Lemma~\ref{lem:center-pushed-in} moves into
$B_{2}$; two classes are then linked through $c$. Otherwise take the largest
$n=2^{k}$ at which fullness fails, apply Lemma~\ref{lem:bacon-left} to $R_{n}$,
push $c$ into $B_{1}$, and conclude by Lemma~\ref{lem:central-relation} that
$(R_{n}\cap B_{1}^{2})/\sigma$ is full after all.
\end{proof}

\begin{longlemma}\label{lem:cut-or-nothing}
Let $\sigma$ be a dividing congruence for $B\lll A$, let $\delta$ be a
congruence on $\mathbf A$ with $(\delta\cap B^{2})/\sigma\neq B^{2}/\sigma$, and
let $\omega$ be the intersection of the dividing congruences of the chosen chain
for $B\lll A$. Then \textup{(1)} $\delta\cap B^{2}\subseteq\sigma\cap B^{2}$;
\textup{(2)} $\sigma\supseteq\delta\cap\omega$; \textup{(3)}
$(\delta\vee(\sigma\cap\omega))\cap B^{2}=\sigma\cap B^{2}$; \textup{(4)}
$(\delta\vee(\sigma\cap\omega))\cap\omega=\sigma\cap\omega$.
\end{longlemma}

\begin{proof}
Throughout we use $B^{2}\subseteq\omega$, which holds because each $B_{i}$ in
the chain lies inside a single block of its own dividing congruence, and $B$ is
contained in every $B_{i}$. \emph{It is used at least four times below.}

(2). The hypothesis gives $\sigma$-classes $C_{1},C_{2}$ meeting $B$ with
$\bigl((C_{1}\cap B)\times(C_{2}\cap B)\bigr)\cap\delta=\varnothing$. Then
$C_{1}\cap B\subt{\TD(\sigma)}B$, and Lemma~\ref{lem:self-intersection-pc} with
$m=k$ and $B_{k}=C_{1}\cap B$ gives
$\bigl((C_{1}\cap B)\circ\delta\cap B\bigr)/\sigma=\{C_{1}\}$, since the value
$B/\sigma$ is excluded by the choice of $C_{2}$; its second clause is
$\sigma\supseteq\delta\cap\omega$.

(1) follows, since $\delta\cap B^{2}\subseteq\delta\cap\omega\subseteq\sigma$.

(4). Put $\delta'=\delta\vee(\sigma\cap\omega)$, which is
$\LeftLinked(\delta\circ(\sigma\cap\omega))$. If
$(\delta'\cap B^{2})/\sigma$ were full then Lemma~\ref{lem:leftlinked-stay-full}
would make $\LeftLinked$ of the restricted relation full, so some
single linking step would cross $\sigma$-classes: there would be
$a_{1},a_{2},c\in B$ and $b_{1},b_{2}\in A$ with $(a_{i},b_{i})\in\delta$,
$(b_{i},c)\in\sigma\cap\omega$ and $(a_{1},a_{2})\notin\sigma$. But
$a_{1},a_{2},c\in B$ gives $(a_{i},c)\in\omega$, hence $(a_{i},b_{i})\in\omega$,
hence $(a_{i},b_{i})\in\sigma$ by (2), hence $(a_{1},a_{2})\in\sigma$ --- a
contradiction. So $(\delta'\cap B^{2})/\sigma\neq(B/\sigma)^{2}$, and applying
(2) to $\delta'$ gives $\sigma\supseteq\delta'\cap\omega$; with
$\delta'\supseteq\sigma\cap\omega$ this is (4).

(3) follows from (4) by intersecting with $B^{2}\subseteq\omega$.
\end{proof}

\begin{lemma}\label{lem:main-existence}
Let $B\lll A$ and let $\sigma$ be a congruence on $\mathbf A$ with
$|B/\sigma|>1$ and $B/\sigma$ BA and centre free. Then there is a dividing
congruence $\delta\supseteq\sigma$ for $B\le A$.
\end{lemma}

\begin{proof}
Choose $\delta\supseteq\sigma$ maximal with $|B/\delta|>1$; it exists because
$\sigma$ qualifies. Then $B/\delta$ is BA and centre free, since the preimage
under $B/\sigma\twoheadrightarrow B/\delta$ of a proper nonempty strong
subuniverse would be one of $B/\sigma$.

It remains to see that $\delta$ is irreducible. Suppose
$\delta=S_{1}\cap\dots\cap S_{k}$ with each $S_{i}\supsetneq\delta$ a
subuniverse of $\mathbf A^{2}$ \emph{stable under $\delta$}. Some $S_{i}$ has
$B^{2}\not\subseteq S_{i}$, else $B^{2}\subseteq\delta$. Now
$\LeftLinked(S_{i})$ is a congruence containing
$S_{i}\supsetneq\delta\supseteq\sigma$, so maximality of $\delta$ forces
$B^{2}\subseteq\LeftLinked(S_{i})$, i.e.\
$(\LeftLinked(S_{i})\cap B^{2})/\delta=(B/\delta)^{2}$. By
Lemma~\ref{lem:leftlinked-stay-full}, applied \emph{with the congruence
$\delta$}, $\LeftLinked(S_{i}\cap B^{2})/\delta=(B/\delta)^{2}$,
so $(S_{i}\cap B^{2})/\delta$ is linked. It is also \emph{proper}: stability
under $\delta$ makes $S_{i}$ a union of products of $\delta$-blocks, so
$(S_{i}\cap B^{2})/\delta=(B/\delta)^{2}$ would give $B^{2}\subseteq S_{i}$.
Corollary~\ref{lem:linked-implies-abs} then produces a proper nonempty strong
subuniverse of $B/\delta$, a contradiction.
\end{proof}

\begin{hazard}[Where each hypothesis is spent]\label{haz:main-existence}
Two things in that proof are easy to lose. The relation must be factored by
$\delta$ and not by $\sigma$ --- $(S_{i}\cap B^{2})/\sigma$ lives on $B/\sigma$
and the conclusion is about $B/\delta$, and it is $B/\delta$ that was shown BA
and centre free. And properness needs the words \emph{stable under $\delta$} in
the unfolding of irreducibility: that clause is what makes $S_{i}$ a union of
products of $\delta$-blocks, and it is the only part of
Definition~\ref{def:irreducible} doing work here. Without it,
$S_{i}\cap B^{2}\subsetneq B^{2}$ survives factoring only by accident.
\end{hazard}

\begin{lemma}\label{lem:factor-by-delta}
Let $\delta$ be a congruence on $\mathbf A$ and
$C\subt[A]{T(\sigma)}B\lll A$ with $T\in\{\TPC,\TL\}$. Then
$C/\delta=B/\delta$, or $C/\delta\subt{\TS}B/\delta$, or
$C/\delta\subt[A/\delta]{T}B/\delta$; and in the last case
$\delta\cap B^{2}\subseteq\sigma\cap B^{2}$.
\end{lemma}

\begin{proof}[Proof outline]
Put $R=\{(a/\sigma,a/\delta):a\in B\}$ and note
$R\circ R^{-1}=(\delta\cap B^{2})/\sigma$, so
$\LeftLinked(R)=(B/\sigma)^{2}$ exactly when
$(\delta\cap B^{2})/\sigma=B^{2}/\sigma$. In that case
Lemma~\ref{lem:bacon-left} yields $U\subte{\TBA,\TC}B/\delta$ with
$(B/\sigma)\times U\subseteq R$, and $U\subseteq C/\delta$, giving the first two
alternatives.

Otherwise Lemma~\ref{lem:cut-or-nothing}(4) gives
$\sigma\cap\omega=\delta'\cap\omega$ for $\delta'=\delta\vee(\sigma\cap\omega)$,
and $B/\delta'\cong B/\sigma$ by (3), so $B/\delta'$ is BA and centre free and
Lemma~\ref{lem:main-existence} produces a dividing congruence
$\delta''\supseteq\delta'$ for $B$. Applying
Lemma~\ref{lem:cut-or-nothing}(2) to $\delta''$ --- legitimate because
$\sigma\cap B^{2}\subseteq\delta''$ and $|B/\delta''|>1$ give
$(\delta''\cap B^{2})/\sigma\neq B^{2}/\sigma$ --- yields
$\delta''\cap\omega=\sigma\cap\omega$, hence $\delta''\cap B^{2}=\sigma\cap
B^{2}$, so the $\delta''$-classes and $\sigma$-classes agree on $B$ and
$C/\delta\subt[A/\delta]{\mathcal T(\delta''/\delta)}B/\delta$. Finally
$\mathcal T=T$: $\omega\setminus\delta''\neq\varnothing$, since
$\omega\supseteq B^{2}$ and $|B/\sigma|>1$, so
Lemma~\ref{lem:bridge-between-congruences} gives a bridge from $\delta''$ to
$\sigma$, and Lemma~\ref{lem:no-cross-bridge} in both directions makes the two
congruences of the same type. The last clause holds because
$\delta\cap B^{2}\subseteq\delta'\cap B^{2}=\sigma\cap B^{2}$.
\end{proof}

\subsection{The headline properties}\label{sec:headline}

\begin{proof}[Proof of Lemma~\ref{lem:ubiquity}]
If $\mathbf B$ has a proper nonempty binary absorbing or central subuniverse,
take it. Otherwise $\mathbf B$ is BA and centre free and $|B|>1$, so
Lemma~\ref{lem:main-existence} with $\sigma=0_{\mathbf A}$ gives a dividing
congruence $\delta$ for $B$; take $C=B\cap E$ for any block $E$ \emph{of
$\delta$} with $B\cap E\neq B$, which exists because $|B/\delta|>1$. Its type is
$\TL$ or $\TPC$ according to $\delta$.
\end{proof}

\begin{proof}[Proof of Lemma~\ref{lem:intersect-all} and
  Corollary~\ref{cor:intersect-all-weak}]
(t). For $T\in\{\TBA,\TC\}$ this is Lemma~\ref{lem:ba-center-intersect}. For
$T=\TS$, take $E\subseteq C$ with $E\subt{\TBA,\TC}B$; if $C\cap D=\varnothing$
there is nothing to prove, and otherwise $B\cap D\neq\varnothing$, so
Lemma~\ref{lem:intersection-good}(s) gives $E\cap D\neq\varnothing$ and
Lemma~\ref{lem:ba-center-intersect} gives $E\cap D\subt{\TBA,\TC}B\cap D$. For
$T\in\{\TPC,\TL\}$, Lemma~\ref{lem:intersection-good}(d) makes $(B\cap D)/\sigma$
empty, a point, or all of $B/\sigma$; the three cases give $C\cap D=\varnothing$,
$C\cap D\in\{\varnothing,B\cap D\}$, and $C\cap D\subt{T(\sigma)}B\cap D$
respectively.

(i). Apply (t) to each step of a chain for $B\lll^{A}A$, obtaining
$B\cap D\dlll^{A}D$. The corollary follows by composing with $D\lll A$.
\end{proof}

\begin{hazard}[The superscript in (i) is needed]\label{haz:intersect-all-strength}
Corollary~\ref{cor:intersect-all-weak} is weaker than
Lemma~\ref{lem:intersect-all}(i), and three later arguments need the stronger
form: Corollary~\ref{cor:prop-relations}(b) and (m), which conclude
$R''\dlll^{R}R'$, and Lemma~\ref{lem:multitype-chain}. Stating only the
corollary --- as is natural, since it is the symmetric-looking form --- makes
all three unprovable.
\end{hazard}

\begin{proof}[Proof of Lemma~\ref{lem:reverse-hom}]
(bt) is proved directly, from the definitions; the reflection corollary of
\S\ref{sec:properties} is \emph{derived} from this lemma and so may not be used
here. For
$T\in\{\TBA,\TC\}$ the witnessing term transfers: $f$ is a homomorphism, so
$f\bigl(t(\bar a)\bigr)=t\bigl(f(\bar a)\bigr)$, and a tuple with all but one
entry in $f^{-1}(C')$ and the remaining entry in $f^{-1}(B')$ maps to one with
all but one entry in $C'$ and the remaining entry in $B'$; the centrality
clause pulls back the same way, since $f\bigl(\Sg(X)\bigr)=\Sg\bigl(f(X)\bigr)$
for surjective $f$. For $T=\TS$ the witness $D'\subseteq C'$ pulls back to
$f^{-1}(D')\subseteq f^{-1}(C')$ and the previous case applies. For
$T\in\{\TD,\TL,\TPC\}$, $C'=B'\cap E$ for a block $E$ of $\sigma$, so
$f^{-1}(C')=f^{-1}(B')\cap f^{-1}(E)$ with $f^{-1}(E)$ a block of
$f^{-1}(\sigma)$; condition (i) of Definition~\ref{def:types} transfers because
$f^{-1}\bigl(\cover\sigma\bigr)=\cover{f^{-1}(\sigma)}$ by
Lemma~\ref{lem:irred-quotient}, and condition (iii) because
$f^{-1}(B')/f^{-1}(\sigma)\cong B'/\sigma$.

(b) follows by iterating (bt) along a chain, and (bm) from (bt) by
intersecting.
\end{proof}

\begin{proof}[Proof of Lemma~\ref{lem:propagation}]
The backward clauses are Lemma~\ref{lem:reverse-hom}.

(fs) is Lemma~\ref{lem:factor-type}. (ft) splits by type: for
$T\in\{\TBA,\TC,\TS\}$ it follows from (fs); for $T\in\{\TPC,\TL\}$ it is
Lemma~\ref{lem:factor-by-delta}; and for $T=\TD$, which that lemma does not
cover, split the witnessing congruence into its linear and PC cases and apply
the lemma to each, the two conclusions being the same three alternatives. Then
(f) follows by iteration, all three outcomes of (ft) staying inside the type
list of Definition~\ref{def:lll}.

(fm). Let $C=C_{1}\cap\dots\cap C_{t}$ with $C_{i}\subt[A]{T(\sigma_{i})}B$ and
induct on $t$. Write $\delta$ for the kernel of $f$, so that $f(X)$ may be read
as $X/\delta$ throughout. For $t=1$ this is (ft), and \emph{the hypothesis that
$f(B)$ is $\TS$-free is exactly what kills its middle alternative}, which
$\MT T$ cannot absorb. If $C_{i}/\delta=B/\delta$ for some $i$, say $i=t$, put
$D_{j}=C_{j}\cap C_{t}$; by Lemma~\ref{lem:intersect-all}(t)
$D_{j}\subte[A]{T(\sigma_{j})}C_{t}$, and the inductive hypothesis applied to
$C\subte{\MT T}C_{t}$ gives the claim. Otherwise every
$C_{i}/\delta\subt[A/\delta]{T}B/\delta$, and
$\delta\cap B^{2}\subseteq\sigma_{i}\cap B^{2}$ for every $i$ by the last clause
of Lemma~\ref{lem:factor-by-delta}. That is what makes the images meet
correctly: $(C_{i}\circ\delta)\cap B=C_{i}$, so a $\delta$-class $E$ lying in
every $C_{i}/\delta$ has $E\cap B\subseteq C_{1}\cap\dots\cap C_{t}$, whence
\[
  (C_{1}\cap\dots\cap C_{t})/\delta=C_{1}/\delta\cap\dots\cap C_{t}/\delta ,
\]
the inclusion $\subseteq$ being trivial. This is the one place where the failure
of images to commute with intersections (\S\ref{sec:legislation}, item
\ref{leg:quotient}) is repaired, and the saturation hypothesis
$\delta\cap B^{2}\subseteq\sigma_{i}\cap B^{2}$ is precisely what buys it.
\end{proof}

\begin{longtheorem}[Stable intersection, restated]
Let $C_{i}\subt[A]{T_{i}(\sigma_{i})}B_{i}\lll A$ with
$T_{i}\in\{\TBA,\TC,\TS,\TL,\TPC\}$ for $i\in[n]$, $n\ge2$, and suppose
$\bigcap_{i}C_{i}=\varnothing$ while $B_{j}\cap\bigcap_{i\neq j}C_{i}\neq
\varnothing$ for every $j$. Then all $T_{i}$ are $\TBA$; or all are $\TL$ and
for all $k,\ell$ there is a bridge $\delta$ from $\sigma_{k}$ to $\sigma_{\ell}$
with $\bcol\delta=\sigma_{k}\circ\sigma_{\ell}$; or $n=2$ and
$T_{1}=T_{2}=\TC$; or $n=2$, $T_{1}=T_{2}=\TPC$ and $\sigma_{1}=\sigma_{2}$.
\end{longtheorem}

\begin{proof}[Proof of Theorem~\ref{thm:stable-intersection}]
For $n=2$ combine Lemmas~\ref{lem:two-stable} and \ref{lem:no-cross-bridge}.
For general $n$ put $B_{2}'=B_{2}\cap C_{3}\cap\dots\cap C_{n}$ and
$C_{2}'=C_{2}\cap B_{2}'$; then $C_{2}'\subt{T_{2}(\sigma_{2})}B_{2}'\lll A$ by
Lemma~\ref{lem:intersect-all}, properly because $C_{2}'=B_{2}'$ would make
$C_{1}\cap B_{2}'$ empty. The three hypotheses of the pair case hold, so
$T_{1}=T_{2}$ and, for the dividing types, the bridge exists.

If $T_{1}=\TPC$ then $\sigma_{1}\circ\sigma_{2}=\sigma_{1}=\sigma_{2}$. Indeed,
composing the bridge with its transpose gives a bridge from $\sigma_{1}$ to
itself, and if $\sigma_{1}\circ\sigma_{2}\supsetneq\sigma_{1}$ its collapse
properly contains $\sigma_{1}$, making $\sigma_{1}$ linear by
Lemma~\ref{lem:linear-equiv}. That leaves the possibility of proper nesting,
$\sigma_{1}\circ\sigma_{2}=\sigma_{1}$ with $\sigma_{1}\neq\sigma_{2}$, which
is excluded by Lemma~\ref{lem:bridge-to-pc}: a bridge between PC congruences
makes $\mathbf A/\sigma_{1}\cong\mathbf A/\sigma_{2}$ with the induced relation
on blocks bijective, and two congruences on one finite algebra with equal
quotient sizes and $\sigma_{2}\subseteq\sigma_{1}$ are equal. With $\sigma_{1}=\sigma_{2}=:\sigma$, the blocks $E_{1},E_{2}$ defining
$C_{1},C_{2}$ are \emph{distinct} --- equal blocks would put the point of
$B_{1}\cap C_{2}'$ in $C_{1}\cap C_{2}'$ --- so $C_{1}\cap C_{2}=\varnothing$;
applying this to every pair and then the hypothesis at $j=3$ forces $n=2$.

For $T_{1}=\TC$ the argument just given has no analogue --- there is no
congruence and no pair of blocks --- and $n=2$ is instead
Lemma~\ref{lem:no-central-critical}, whose two hypotheses are exactly (ii) and
(iii) here.
\end{proof}

\begin{hazard}[What the type-$\TC$ case costs]\label{haz:d4}
The type-$\TC$ argument is the one place in this section that reaches outside
the material developed here, to \cite[Lem.~6.11]{ZhukStrong}: \emph{for
$n\ge3$ and $E_{i}$ central in $\mathbf D$, there is no
$(E_{1},\dots,E_{n})$-essential relation $R\le D^{n}$}. Two notes for anyone
checking it. Its printed proof cites \emph{itself} three times where it means
the preceding lemma, so a self-contained rendering must supply both. And its
ingredients are exactly the Barto--Kazda criterion
(Lemma~\ref{lem:barto-kazda}) and ``central implies ternary absorbing''
(Lemma~\ref{lem:central-ternary}), both of which are already available above.

Nonemptiness and properness of each $E_{i}$ are established before the
essential-relation theorem is applied, and both are needed: essentiality is a
statement about proper nonempty subsets.
\end{hazard}

\begin{proof}[Proof of Corollary~\ref{cor:stable-intersection}]
Pull back along the projections $f_{i}:\mathbf R\to\mathbf A_{i}$ using
Lemma~\ref{lem:reverse-hom}(b),(bt), apply
Theorem~\ref{thm:stable-intersection} inside $\mathbf R$, and translate
$\bcol\delta=\sigma_{k}'\circ\sigma_{\ell}'$ back as
$\sigma_{k}\circ\proj_{k,\ell}(R)\circ\sigma_{\ell}$.
\end{proof}

\begin{proof}[Proof of Lemma~\ref{lem:multitype-chain}]
With $C=\bigcap_{i\le n}C_{i}$ and $D_{j}=\bigcap_{i\le j}C_{i}$,
Lemma~\ref{lem:intersect-all}(t) gives $D_{j+1}\subte[A]{T}D_{j}$, using
$D_{j}\lll A$ from part (i); collapsing the equalities gives the chain.
\end{proof}

\begin{proof}[Proof of Lemma~\ref{lem:anchored-intersection}]
By Lemma~\ref{lem:multitype-chain} fix a chain
$C_{2}=F_{s}\subt[A_{2}]{T}\dots\subt[A_{2}]{T}F_{0}=B_{2}$. If $s=0$ then
$C_{2}=B_{2}$ and the second part of (i) is the conclusion. Each truncation
of the chain witnesses $F_{j}\lll^{A_{2}}B_{2}$, and concatenating it with
the chain for $B_{2}\lll A_{2}$ gives $F_{j}\lll A_{2}$ for every $j$.

Suppose for contradiction that $W\cap(C_{1}\times C_{2})=\varnothing$. Along
the chain, $W\cap(C_{1}\times F_{0})\neq\varnothing$ by (i) and
$W\cap(C_{1}\times F_{s})=\varnothing$, so there is a least $j$ with
$W\cap(C_{1}\times F_{j})=\varnothing$; fix it, so that
$W\cap(C_{1}\times F_{j-1})\neq\varnothing$. Write the step as
$F_{j}\subt[A_{2}]{T(\xi)}F_{j-1}$ and let $T'$ be the class of $\xi$ ---
linear or PC, exhaustively, by Definition~\ref{def:linear-congruence} --- so
that $F_{j}\subt[A_{2}]{T'(\xi)}F_{j-1}$; for $T\in\{\TL,\TPC\}$ the class
is $T'=T$, and for $T=\TD$ it is whichever the step's congruence has.

Apply Corollary~\ref{cor:stable-intersection} to $W$ with the two-member
family
\[
  C_{1}\subt[A_{1}]{T_{1}(\sigma_{1})}B_{1}\lll A_{1},
  \qquad
  F_{j}\subt[A_{2}]{T'(\xi)}F_{j-1}\lll A_{2}.
\]
Its hypotheses hold: $W$ is subdirect; the joint intersection
$W\cap(C_{1}\times F_{j})$ is empty by the choice of $j$; and both
leave-one-out configurations are nonempty ---
$W\cap(B_{1}\times F_{j})\supseteq W\cap(B_{1}\times C_{2})\neq\varnothing$
by (i), since $C_{2}=F_{s}\subseteq F_{j}$, and
$W\cap(C_{1}\times F_{j-1})\neq\varnothing$ by the minimality of $j$. So one
of the four alternatives holds of the pair $\{T_{1},T'\}$. None can:
$T'\in\{\TL,\TPC\}$, so (ba) and (c) fail; $\TS$ occurs in no alternative;
and (l), (pc) require $T_{1}=T'\in\{\TL,\TPC\}$, which (ii) excludes ---
for $T_{1}=\TL$ it puts $T=\TPC$, making every step of the chain of class
$\TPC$, and symmetrically for $T_{1}=\TPC$. The contradiction refutes the
assumption, so $W\cap(C_{1}\times C_{2})\neq\varnothing$.
\end{proof}

\begin{proof}[Proof of Lemma~\ref{lem:linear-on-top}]
By Lemma~\ref{lem:linear-equiv} there is a bridge $\delta$ from $\sigma$ to
$\sigma$ with $\bcol\delta\supsetneq\sigma$, reflexive by
Lemma~\ref{lem:bridge-reflexive}; replace it by $\delta^{\circ}$, symmetric and
pair-reflexive.
Lemma~\ref{lem:nontrivial-bridge}(c) applies, and $\cover\sigma=A^{2}$ is a
congruence with the single block $A$, so $A/\sigma\cong\mathbb Z_{p}^{m}$ as a
group. Now $m=0$ gives $\sigma=A^{2}$, impossible for an irreducible congruence
(\S\ref{sec:legislation}, item \ref{leg:emptyfamily}). And $m\ge2$ contradicts
irreducibility: $A/\sigma$ is affine, so the basic operation acts as
$\sum_{i}c_{i}x_{i}$ with $\sum_{i}c_{i}=1$, and the coordinate congruences
$\theta_{i}=\{(x,y):x_{i}=y_{i}\}$ are linear subspaces of
$(\mathbb Z_{p}^{m})^{2}$, hence subuniverses of $(\mathbf A/\sigma)^{2}$, each
properly containing $0_{\mathbf A/\sigma}$, with
$\bigcap_{i}\theta_{i}=0_{\mathbf A/\sigma}$ --- so
$0_{\mathbf A/\sigma}$ is reducible, and by Lemma~\ref{lem:irred-quotient} so is
$\sigma$. Hence $m=1$, and the group isomorphism is an isomorphism of algebras
by Proposition~\ref{prop:p-divides}.
\end{proof}

\begin{lemma}\label{lem:multiply-all-linear}
Let $C_{1}\subt[A]{\MT T}B_{1}\lll A$ with $T\in\{\TPC,\TL,\TD\}$, let
$B_{2}\lll A$, $C_{1}\cap B_{2}=\varnothing$ and $B_{1}\cap B_{2}\neq
\varnothing$. Then $\bigl(C_{1}\circ(\omega_{1}\cap\dots\cap\omega_{s})\bigr)
\cap B_{2}=\varnothing$, where $\omega_{1},\dots,\omega_{s}$ are all congruences
of type $T$ on $\mathbf A$ with $\cover{\omega_{i}}\supseteq B_{1}^{2}$.
\end{lemma}

\begin{proof}[Proof of Case~1, and of Case~2 from Claim~\ref{cl:multiply-step}]
The induction is downward on $|B_{1}|$: the inductive hypothesis is the lemma
for every larger first argument, and it starts at $B_{1}=A$.

By Definition~\ref{def:multitypes} write $C_{1}=\bigcap_{i\in[t]}C_{1}^{i}$
with $t\ge1$ and $C_{1}^{i}\subt[A]{T(\sigma_{i})}B_{1}$ --- of type $T$, not
merely of type $\TD$, which is what makes each $\sigma_{i}$ a congruence of type
$T$ with $\cover{\sigma_{i}}\supseteq B_{1}^{2}$ and hence one of the
$\omega_{j}$. Put $K_{i}=C_{1}^{i}\circ\sigma_{i}$ and
$\Omega=\omega_{1}\cap\dots\cap\omega_{s}$. Definition~\ref{def:types}(ii)
gives $C_{1}^{i}=K_{i}\cap B_{1}$, and
Corollary~\ref{cor:prop-times-cong}(e),(f) gives
$K_{i}\subt[A]{T(\sigma_{i})}B_{1}\circ\sigma_{i}\lll A$, so in particular
$K_{i}\lll A$. Two inclusions are used, and both are immediate:
\begin{equation}\label{eq:mal-incl}
  C_{1}\circ\Omega\;\subseteq\;\bigcap_{i\in[t]}K_{i},
  \qquad\text{and}\qquad
  X\subseteq Y\;\Longrightarrow\;X\circ\Omega\subseteq Y\circ\Omega .
\end{equation}
The first holds because $C_{1}\subseteq C_{1}^{i}$ and
$\Omega\subseteq\sigma_{i}$ for each $i$, the latter because $\sigma_{i}$ is
one of the $\omega_{j}$.

Put $D=\bigcap_{i\in[t]}K_{i}\cap B_{2}$, and note
$D\cap B_{1}=\bigcap_{i}(K_{i}\cap B_{1})\cap B_{2}=C_{1}\cap B_{2}=
\varnothing$.

\emph{Case 1: $D=\varnothing$.} Then by \eqref{eq:mal-incl}
$(C_{1}\circ\Omega)\cap B_{2}\subseteq\bigcap_{i}K_{i}\cap B_{2}=\varnothing$,
which is the conclusion. This is the whole of the case $B_{1}=A$, since there
$D=D\cap B_{1}=\varnothing$.

\emph{Case 2: $D\neq\varnothing$.} Each $K_{i}\lll A$ and $B_{2}\lll A$, so
Corollary~\ref{cor:intersect-all-weak}, applied $t$ times, gives
$D\dlll A$; as $D\neq\varnothing$, $D\lll A$. Fix the chain
$A=X_{0}\supsetneq X_{1}\supsetneq\dots\supsetneq X_{r}=B_{1}$ witnessing
$B_{1}\lll A$. Then $D\cap X_{0}=D\neq\varnothing$ and $D\cap X_{r}=
D\cap B_{1}=\varnothing$, so there is a least $p\in[r]$ with
$D\cap X_{p}=\varnothing$. Write $B_{1}'=X_{p}$ and $B_{1}''=X_{p-1}$, so that
\begin{equation}\label{eq:mal-step}
  B_{1}\subseteq B_{1}'\subt[A]{}B_{1}''\lll A,\qquad
  D\cap B_{1}'=\varnothing,\qquad D\cap B_{1}''\neq\varnothing .
\end{equation}

\begin{claim}\label{cl:multiply-step}
That step is of multi-type $T$: $B_{1}'\subt[A]{\MT T}B_{1}''$.
\end{claim}

Granting the claim, the inductive hypothesis applies to
$B_{1}'\subt[A]{\MT T}B_{1}''\lll A$ and $D\lll A$, whose hypotheses
$B_{1}'\cap D=\varnothing$ and $B_{1}''\cap D\neq\varnothing$ are
\eqref{eq:mal-step}, and which sits at a strictly larger first argument since
$B_{1}\subseteq B_{1}'\subsetneq B_{1}''$. It returns
$(B_{1}'\circ\Omega'')\cap D=\varnothing$, where $\Omega''$ is the intersection
of all type-$T$ congruences $\omega$ with $\cover{\omega}\supseteq
(B_{1}'')^{2}$. Every such $\omega$ has $\cover{\omega}\supseteq B_{1}^{2}$,
so that family is a \emph{subfamily} of $\{\omega_{1},\dots,\omega_{s}\}$ and
$\Omega''\supseteq\Omega$; hence $(B_{1}'\circ\Omega)\cap D=\varnothing$ as
well. Finally, by \eqref{eq:mal-incl} and $C_{1}\subseteq B_{1}\subseteq
B_{1}'$,
\[
  (C_{1}\circ\Omega)\cap B_{2}\;\subseteq\;
  (B_{1}'\circ\Omega)\cap\bigcap_{i\in[t]}K_{i}\cap B_{2}\;=\;
  (B_{1}'\circ\Omega)\cap D\;=\;\varnothing . \qedhere
\]
\end{proof}

\begin{hazard}[Claim~\ref{cl:multiply-step} is not proved, and it is the whole
  difficulty]\label{haz:multiply-citation}
Everything above is a proof except the claim, and the claim carries two
separate obligations. Neither is discharged here, so
Lemma~\ref{lem:multiply-all-linear} is \emph{not} available; see
\S\ref{sec:status}.

\emph{(a) The step is not of type $\TBA$, $\TC$ or $\TS$.} For $\TS$ this is
Lemma~\ref{lem:intersection-good}(s): $B_{1}'$ then contains some
$G\subt{\TBA,\TC}B_{1}''$, and (s) applied to $B_{1}''$ and $D$ --- which meet,
by \eqref{eq:mal-step} --- gives $G\cap D\neq\varnothing$, contradicting
$B_{1}'\cap D=\varnothing$. For $\TBA$ and $\TC$ separately the same route is
\emph{not} available, because $\subt{\TBA,\TC}$ in (s) and in
Lemma~\ref{lem:strong-nonempty} is a conjunction
(Caveat~\ref{haz:ba-central-meets}) and a single absorbing step supplies only
one of the two types. Lemma~\ref{lem:intersect-all}(t) gives
$B_{1}'\cap D\dsubte{\TBA}B_{1}''\cap D$, whose dot is exactly the possibility
to be excluded.

\emph{(b) When $T\in\{\TL,\TPC\}$, a type-$\TD$ step is not enough.} The
inductive hypothesis needs $B_{1}'\subt{\MT T}B_{1}''$ for the \emph{same} $T$.
A $\TD$ step carries an irreducible congruence, which is linear or PC by
Caveat~\ref{haz:linear-def}, but nothing so far forces it to match $T$. This
obligation is invisible while the components are written as $\TD$, which is why
the typing in the first paragraph of the proof matters.

Obligation (b) is the more serious: it is not a suppressed step but a missing
hypothesis in the induction, and closing it plausibly means strengthening the
inductive statement to carry the types of the intervening steps rather than
patching the claim.
\end{hazard}

\begin{proof}[Proof of Lemma~\ref{lem:maximal-mult}]
Write $\bar X$ for $X/\sigma$ and $\bar{\mathbf A}=\mathbf A/\sigma$. Note
first that for any $X\subseteq A$, $(X\circ\sigma)\cap B_{2}=\varnothing$ says
exactly that no $\sigma$-class meets both $X$ and $B_{2}$, i.e.\ that
$\bar X\cap\bar B_{2}=\varnothing$. So the hypotheses read
$\bar C_{1}\cap\bar B_{2}=\varnothing$ and, by
\ref{hyp:maximal-mult}, $\bar P\cap\bar B_{2}=\varnothing$; and
$\bar B_{1}\cap\bar B_{2}\neq\varnothing$ because
$B_{1}\cap B_{2}\neq\varnothing$.

\emph{Step 1: images commute with intersection on $F$.} For $j\in F$ the factor
$\bar C_{1}^{j}$ lands in the type-$T$ alternative of
Lemma~\ref{lem:propagation}(ft), so the ``moreover'' clause of
Lemma~\ref{lem:factor-by-delta} gives
$\sigma\cap B_{1}^{2}\subseteq\sigma_{j}\cap B_{1}^{2}$. Hence
\[
  \bigcap_{j\in F}\bar C_{1}^{j}\;=\;\bar P .
\]
Indeed $\subseteq$ is the only direction needing an argument. Let
$[a]\in\bigcap_{j\in F}\bar C_{1}^{j}$; since each
$\bar C_{1}^{j}\subseteq\bar B_{1}$ we may pick $a\in B_{1}$. For each
$j\in F$ there is $a_{j}\in C_{1}^{j}$ with $(a,a_{j})\in\sigma$; as
$a,a_{j}\in B_{1}$ this puts $(a,a_{j})\in\sigma\cap B_{1}^{2}\subseteq
\sigma_{j}$, so $a$ lies in the same $\sigma_{j}$-block as $a_{j}$, and
$C_{1}^{j}=B_{1}\cap E_{j}$ for that block gives $a\in C_{1}^{j}$. So
$a\in P$.

\emph{This is what \S\ref{sec:legislation}, item \ref{leg:quotient} warns
about, and it is exactly where the warning does not bite}: among the type-$T$
factors, and only among them, the saturation identity holds and the image of
the intersection is the intersection of the images.

\emph{Step 2: a proper multi-type reduction modulo $\sigma$.} Each
$\bar C_{1}^{j}\subt[\bar A]{T}\bar B_{1}$ for $j\in F$, and $\bar P$ is
nonempty because $P\supseteq C_{1}\neq\varnothing$. So by Step 1
\[
  \bar P\subte[\bar A]{\MT T}\bar B_{1},
\]
and the containment is \emph{proper}, since
$\bar P\cap\bar B_{2}=\varnothing$ while
$\bar B_{1}\cap\bar B_{2}\neq\varnothing$. Also $\bar B_{1}\lll\bar A$ and
$\bar B_{2}\lll\bar A$ by Corollary~\ref{cor:prop-quotient}(f).

\emph{Step 3: multiply by all type-$T$ congruences.} Apply
Lemma~\ref{lem:multiply-all-linear} inside $\bar{\mathbf A}$ to
$\bar P\subt[\bar A]{\MT T}\bar B_{1}\lll\bar A$ and
$\bar B_{2}\lll\bar A$: it returns the congruences
$\delta_{1},\dots,\delta_{s}$ of type $T$ on $\bar{\mathbf A}$ with
$\cover{\delta_{i}}\supseteq\bar B_{1}^{2}$, and
$\bigl(\bar P\circ\bigcap_{i}\delta_{i}\bigr)\cap\bar B_{2}=\varnothing$.
Note the ambient here is still $\bar B_{1}$, which is what the conclusion's
condition $\cover{\omega_{i}}\supseteq B_{1}^{2}$ requires.

\emph{Step 4: back to $\mathbf A$.} Let $\omega_{i}$ be the preimage of
$\delta_{i}$ under $\mathbf A\twoheadrightarrow\bar{\mathbf A}$, so
$\omega_{i}\supseteq\sigma$. By Lemma~\ref{lem:irred-quotient} the
correspondence $\omega\mapsto\omega/\sigma$ preserves irreducibility,
linearity and PC-ness and carries covers to covers, so each $\omega_{i}$ is of
type $T$ with $\cover{\omega_{i}}\supseteq B_{1}^{2}$. Pulling Step 3 back,
$\bigl(P\circ\bigcap_{i}\omega_{i}\bigr)\cap B_{2}=\varnothing$, and
$C_{1}\subseteq P$ gives
$\bigl(C_{1}\circ\bigcap_{i}\omega_{i}\bigr)\cap B_{2}=\varnothing$.

Finally $\bigcap_{i}\omega_{i}\supseteq\sigma$, and it has the property that
$\sigma$ was chosen maximal for; so $\sigma=\bigcap_{i}\omega_{i}$.
\end{proof}

\begin{hazard}[What Hypothesis~\ref{hyp:maximal-mult} costs, and three failed
  routes]\label{haz:maximal-mult}
What the other hypotheses give is
$\bar C_{1}\cap\bar B_{2}=\varnothing$, and $C_{1}\subseteq P$, so
$\bar C_{1}\subseteq\bar P$: Hypothesis~\ref{hyp:maximal-mult} asks for
emptiness of the \emph{larger} set, and is therefore strictly stronger. Three
routes to it have been tried and none works.

\emph{Drop the non-type-$T$ factors one at a time.} If $j_{0}\notin F$ and
$C_{1}'=\bigcap_{j\neq j_{0}}C_{1}^{j}$, one would like
$\bar C_{1}'\cap\bar B_{2}=\varnothing$ and then induct. But recovering that
from $\bar C_{1}\cap\bar B_{2}=\varnothing$ needs
$\bar C_{1}=\bar C_{1}'\cap\bar C_{1}^{j_{0}}$, which is the
image-versus-intersection identity again --- available for type-$T$ factors by
Step~1 of the proof above and \emph{not} for the others, since the saturation
clause of Lemma~\ref{lem:factor-by-delta} is exactly what the $\subt{\TS}$
alternative withholds.

\emph{Pass to a minimal subfamily.} Choosing $F'$ minimal with
$\bigcap_{j\in F'}\bar C_{1}^{j}\cap\bar B_{2}=\varnothing$ would make
minimality supply the leave-one-out hypothesis of
Theorem~\ref{thm:stable-intersection}, whose four conclusions have no room for
$\TS$ --- so no member of $F'$ would be $\TS$-typed. But no such $F'$ need
exist: the whole family gives
$\bigl(\bigcap_{j}C_{1}^{j}\bigr)/\sigma\cap\bar B_{2}=\varnothing$, whereas
$\bigl(\bigcap_{j}C_{1}^{j}\bigr)/\sigma\subseteq\bigcap_{j}\bar C_{1}^{j}$
with the inclusion generally strict (\S\ref{sec:legislation}, item
\ref{leg:quotient}). This is the route an earlier draft of this document took,
and it is invalid.

\emph{Descend into a strong subuniverse.} The $\subt{\TS}$ factors do at least
never help: a factor with $\bar C_{1}^{j}\subt{\TS}\bar B_{1}$ contains a
subuniverse $\bar D$ of $\bar B_{1}$ that is simultaneously binary absorbing
and central, and Lemma~\ref{lem:strong-nonempty} makes $\bar D$ meet
$\bar B_{1}\cap\bar B_{2}$, so $\bar C_{1}^{j}\cap\bar B_{2}\neq\varnothing$.
But the configuration reproduces itself on $\bar D$ --- $\bar D\lll\bar A$,
$\bar C_{1}\cap\bar D\lll\bar D$, $\bar B_{2}\cap\bar D\lll\bar D$, still
disjoint, with $\bar D\cap\bar B_{2}\neq\varnothing$ --- so no iteration of
that fact closes anything, and the descent shrinks the ambient while the
conclusion's congruence family is indexed by the ambient.

So the hypothesis stays, and the debt is recorded in \S\ref{sec:status}.

A fourth route is open, and unlike the three above it has not failed --- it
has not been attempted. Lemma~\ref{lem:maximal-mult} has exactly one consumer,
the typing step of Lemma~\ref{lem:parallelogram-crucial}, and the older proof
of the dichotomy establishes the parallelogram property of a crucial
constraint \emph{without any statement of this shape}: by recursion on a
measure over reductions, restricting to the minimal linear reduction
containing a solution of the weakened instance, with a fit-in lemma for
one-of-four chains through a subdirect relation doing the intersection work.
No maximal congruence appears. If Lemma~\ref{lem:parallelogram-crucial} is
proved along those lines, its typing conclusion recovered separately ---
irreducibility from Lemma~\ref{lem:crucial-irreducible}, the type from the
bridge machinery --- then the sole call site of
Lemma~\ref{lem:maximal-mult} disappears, and
Hypothesis~\ref{hyp:maximal-mult} with it. The transport cost is not small:
the fit-in lemma has to be restated in this document's vocabulary, and the
recursion needs the reduction machinery of \S\ref{sec:main-statements}. But it
converts an open problem into writing, and it is the first route with that
property.
\end{hazard}

\clearpage
%!TeX root=csp-proof.tex
\section{The main statements}\label{sec:main-statements}

Everything converges here.
Theorem~\ref{thm:main-inductive} is the single induction that carries the whole
proof; the three theorems the algorithm consumes are corollaries of it.

\subsection{The three claims the algorithm needs}

The algorithm of \S\ref{sec:algorithm} is justified by exactly three facts.
Informally:

\begin{description}\tightlist
\item[\textup{(IC1)}]\label{ic:existence} every domain of size $\ge 2$ has a
  strong subset, or an equivalence $\sigma$ with
  $D_{x}/\sigma\cong\mathbb Z_{p}$;
\item[\textup{(IC2)}]\label{ic:safe} reducing a variable's domain to a strong
  subset of it does not destroy solvability, for a consistent enough instance;
\item[\textup{(IC3)}]\label{ic:codim} for a linked, consistent enough,
  strong-subset-free instance, the set of ``linear cosets'' containing a
  solution is empty, full, or a codimension-one affine subspace.
\end{description}

(IC1) is Lemma~\ref{lem:ubiquity} together with
Lemma~\ref{lem:linear-on-top}. (IC2) is
Theorem~\ref{thm:reductions-safe}. (IC3) is
Theorem~\ref{thm:codim-one}.

\subsection{The main induction}

\begin{longtheorem}[Main inductive claim]\label{thm:main-inductive}
Let $\mathcal I$ be an irreducible cycle-consistent instance and let $D^{(1)}$
be a $1$-consistent reduction for $\mathcal I$ with $D^{(1)}\lll D$.
\begin{enumerate}\tightlist
\item[\textup{(1)}] If $\mathcal I$ is crucial in $D^{(1)}$ then
  \textup{(1a)} holds, and \textup{(1b)} or \textup{(1c)} holds:
  \begin{enumerate}\tightlist
  \item[\textup{(1a)}] every constraint of $\mathcal I$ has the parallelogram
    property;
  \item[\textup{(1b)}] $\mathcal I$ is a connected linear-type instance whose
    solution set is subdirect;
  \item[\textup{(1c)}] there is an expanded covering $\mathcal J$ of
    $\mathcal I$, crucial in $D^{(1)}$, with a linked connected subinstance
    $\Upsilon$ whose solution set is not subdirect.
  \end{enumerate}
\item[\textup{(2)}] If $D^{(2)}\subte{T}D^{(1)}$ is a $1$-consistent reduction
  for $\mathcal I$ with $T\in\{\TBA,\TC\}$ --- and, when $T=\TBA$, uniform:
  $D^{(2)}\subte{\TBA(t)}D^{(1)}$ for a single binary term $t$
  (Caveat~\ref{haz:uniform-witness}) --- and $\mathcal I^{(1)}$ has a solution,
  then $\mathcal I^{(2)}$ has a solution.
\end{enumerate}
\end{longtheorem}

\begin{hazard}[The induction, stated precisely]\label{haz:the-induction}
``Induction on the size of $D^{(1)}$'' hides the entire well-foundedness
argument, and it must be made precise, because \emph{the two halves of the
conclusion are proved by mutual induction and each invokes the other}.

The intended measure is
\[
  \mu(D^{(1)})=\sum_{x\in\Var(\mathcal I)}\bigl|D^{(1)}_{x}\bigr|\in\mathbb N,
\]
and the statement being proved by strong induction on $\mu$ is the
conjunction of (1) and (2), \emph{universally quantified over $\mathcal I$}:
\[
  P(k):\quad
  \forall\,\mathcal I,\ \forall\,D^{(1)}\ \text{with}\ \mu(D^{(1)})\le k,\
  \textup{(1)}\ \text{and}\ \textup{(2)}\ \text{hold}.
\]
The quantifier over $\mathcal I$ must be \emph{inside}, because every use of the
inductive hypothesis changes the instance: (2) applies it to a weakening
$\mathcal I'$ of $\mathcal I^{(2)}$; Case~1 of (1) applies it to $\mathcal I$
with $D^{(2)}$; Case~2 applies it to weakenings and to expanded coverings. An
induction that fixed $\mathcal I$ and inducted on $D^{(1)}$ alone does not go
through.

Three further points.
\begin{enumerate}\alphlabels\tightlist
\item \emph{The two halves are established in a fixed order at each level.}
  Part (2) at a given measure uses (1) only at strictly smaller measure,
  whereas part (1) at that measure uses \emph{(2) at the same measure} ---
  Case~1 below does exactly that. So the scheme is: strong induction on
  $k=\mu(D^{(1)})$; at each level establish (2) for all instances, then (1) for
  all instances. That is acyclic, and it is why part (2) is proved first.
\item \emph{Which appeals decrease.} (2) invokes the hypothesis at
  $D^{(2)}\subsetneq D^{(1)}$. The containment is strict in the only case that
  needs an argument: $D^{(2)}\subte{T}D^{(1)}$ permits equality, but if
  $D^{(2)}=D^{(1)}$ then $\mathcal I^{(2)}=\mathcal I^{(1)}$ and (2) is its own
  hypothesis. Case~1 of (1) likewise. Case~2 invokes it at
  $D^{(B,\bot)}\subseteq D^{(B,\top)}$, which is $D^{(1)}$ with the $z$-domain
  cut to $B\lneq D^{(1)}_{z}$, the properness being that of
  $B\subt[D_{z}]{T(\sigma)}D^{(1)}_{z}$. In every case the decrease is at a
  single variable and by at least one element. The uniform statement is:
  \emph{if $D'\subte{T}D$ with $T$ any type and $D'\neq D$, then
  $\mu(D')<\mu(D)$}.
\item \emph{The measure does survive, because no appeal is made at an
  expanded covering.} This looks like the weak point of the scheme and is not.
  An expanded covering has \emph{more} variables than $\mathcal I$, and the
  reduction of a covering is the extension of $D^{(1)}$ along the parent map
  (Proposition~\ref{prop:covering-basics}(h)), which repeats the parent
  domains; its $\mu$ is \emph{larger}, and no observation about one parent
  variable being cut compares the two sums, since arbitrarily many children can
  outweigh the decrease. So if the induction did appeal at a covering, $\mu$
  would not serve.

  It does not. Every appeal above is to a \emph{weakening} of the current
  instance --- $\Var$ can only shrink --- paired with a reduction strictly
  smaller in $\mu$: part (2) at $D^{(2)}$; Case~1 of (1) at $D^{(2)}$ twice;
  the proof of (1a) at the minimal $\MT T$ reduction; and the three appeals in
  the $\mathcal B_{0}\neq\varnothing$ branch at $D^{(B,\bot)}$, whose members
  $\mathcal I'$ come from $\Omega$, every member of which is a weakening of
  $\mathcal I$ by condition~1 of Construction~\ref{constr:omega}. Expanded
  coverings enter only as \emph{outputs} --- the $\mathcal J$ of (1c), and the
  auxiliary $\bot(\cdot)$, $\Delta(\cdot)$ and $\Theta$ built from them --- and
  an output carries no well-foundedness obligation.

  What does apply this theorem to a covering is
  Theorem~\ref{thm:pc-safe}, which is a separate statement proved
  \emph{after} it and consuming it as a black box; the appeal there is an
  ordinary application of a proved theorem, not a step in this induction.
  Confusing the two is the natural mistake, and it is what an earlier draft of
  this caveat made.
\end{enumerate}
\end{hazard}

\begin{proof}[Proof of \textup{(2)}]
Suppose $\mathcal I^{(2)}$ has no solution. Weaken $\mathcal I^{(2)}$ to an
instance $\mathcal I'$ crucial in $D^{(2)}$ (Remark~\ref{rem:get-crucial}). By
the inductive hypothesis applied to $\mathcal I'$ and $D^{(2)}$, $\mathcal I'$
satisfies (1a) and (1b) or (1c).

\emph{Case \textup{(1c)}.} There is an expanded covering $\mathcal J$ of
$\mathcal I'$, crucial in $D^{(2)}$, with a linked connected subinstance
$\Upsilon$. Let $x$ be a variable of a constraint $C\in\Upsilon$. By
Lemma~\ref{lem:connected}(p), $\ConOne(C,x)$ is a perfect linear congruence;
choose $\zeta\subseteq\mathbf D_{x}\times\mathbf D_{x}\times\mathbf Z_{p}$ with
$(y_{1},y_{2},0)\in\zeta\iff(y_{1},y_{2})\in\ConOne(C,x)$ and
$\proj_{1,2}(\zeta)=\cover{\ConOne(C,x)}$.

Replace the variable $x$ of $C$ in $\mathcal J$ by a fresh $x'$ and add the
constraint $\zeta(x,x',z)$ with $z$ fresh of domain $\mathbf Z_{p}$; call the
result $\Theta$, and extend $D^{(1)},D^{(2)}$ by $D^{(1)}_{x'}=D^{(2)}_{x'}
=D_{x'}$. Let $E_{i}$ be the set of $b\in\mathbf Z_{p}$ such that $\Theta$ has
a solution in $D^{(i)}$ with $z=b$. By
Lemma~\ref{lem:ba-center-implies}, $E_{2}\dsubte{T}E_{1}$.

Since $\mathcal J$ is crucial in $D^{(2)}$, $\Theta^{(2)}$ has a solution but
none with $z=0$; since $\mathcal I^{(1)}$ has a solution, so does
$\mathcal J^{(1)}$, hence $\Theta^{(1)}$ has one with $z=0$. Thus
$0\in E_{1}$, $0\notin E_{2}$, and $E_{1}$ has at least two elements. As
$\mathbf Z_{p}$ has no proper subalgebra of size $>1$, $E_{1}=\mathbf Z_{p}$;
then $E_{2}\subt{T}\mathbf Z_{p}$ with $T\in\{\TBA,\TC\}$, contradicting
Lemma~\ref{lem:no-abs-in-linear}.

\emph{Case \textup{(1b)}.} Corollary~\ref{cor:same-type} gives a
contradiction directly.
\end{proof}

\begin{hazard}[Three gaps in the case \textup{(1c)} argument]\label{haz:case1c}
\begin{enumerate}\alphlabels\tightlist
\item ``$E_{1}$ contains at least two different elements, $0\in E_{1}$ and
  $0\notin E_{2}$, hence $E_{1}=\mathbf Z_{p}$'' uses that $E_{1}$ is a
  subuniverse of $\mathbf Z_{p}$ --- true because $E_{1}=\Theta^{(1)}(z)$ is a
  relation defined by an instance (Definition~\ref{def:instance-relation}) ---
  \emph{and} that the subuniverses of $\mathbf Z_{p}$ are the singletons and
  the whole set. The latter needs $w^{\mathbf Z_{p}}$ to generate: from $a\ne b$
  one reaches every element. Prove it; it is one line but it is the crux.
\item $E_{2}\dsubte{T}E_{1}$ is applied with the \emph{dotted} conclusion of
  Lemma~\ref{lem:ba-center-implies}, and the argument then treats $E_{2}$ as a
  proper subuniverse. The case $E_{2}=\varnothing$ must be handled: it is
  excluded because $\Theta^{(2)}$ has a solution. Say so.
\item The choice of $\zeta$ is by Definition~\ref{def:perfect-linear}, which
  supplies $\zeta$ with $\proj_{1,2}\zeta=\cover\sigma$ and
  $(a_{1},a_{2},b)\in\zeta\Rightarrow((a_{1},a_{2})\in\sigma\iff b=0)$. The
  proof uses the \emph{biconditional} $(y_{1},y_{2},0)\in\zeta\iff(y_{1},y_{2})
  \in\ConOne(C,x)$, which is stronger: it needs, for each $(y_1,y_2)\in\sigma$,
  that $(y_{1},y_{2},0)\in\zeta$. This follows from $\proj_{1,2}\zeta
  =\cover\sigma\supseteq\sigma$ plus the implication, but the step is
  suppressed.
\end{enumerate}
\end{hazard}

\begin{proof}[Proof of \textup{(1)}, sketch of the structure]
First, $|D^{(1)}_{x}|>1$ for some $x$: otherwise $\mathcal I^{(1)}$ has all
domains singletons and $1$-consistency makes it solvable, contradicting
cruciality.

\emph{Case 1: some $D^{(1)}_{x}$ has a nontrivial $\TBA$ or central
subuniverse.} Lemma~\ref{lem:find-consistent} yields a $1$-consistent reduction
$D^{(2)}\subte{T}D^{(1)}$, $T\in\{\TBA,\TC\}$. One checks $\mathcal I$ is
crucial in $D^{(2)}$: for a weakening $\mathcal J$ of a constraint,
$\mathcal J^{(1)}$ has a solution, so by the inductive hypothesis (part (2) at
$D^{(2)}$) $\mathcal J^{(2)}$ has one. Applying the inductive hypothesis
(part (1)) at $D^{(2)}$ gives the conclusion.

\emph{Case 2: otherwise.} Lemma~\ref{lem:ubiquity} gives
$E\subt[D_{z}]{T(\sigma)}D^{(1)}_{z}$ for some $z$ and $T\in\{\TPC,\TL\}$.
Property (1a) is obtained by choosing, for each $x$, a minimal
$D^{(2)}_{x}\subte{\MT T}D^{(1)}_{x}$ containing the value at $x$ of a solution
produced by weakening a fixed constraint $C$; Lemma~\ref{lem:multitype-chain}
gives $D^{(2)}\lll^{D}D^{(1)}$ and
Lemma~\ref{lem:minimal-consistent} splits into two subcases, in one of which the
inductive hypothesis applies and in the other of which
Lemma~\ref{lem:parallelogram-crucial} applies directly.

For (1b)/(1c) one forms, for each $B$ in the family
$\mathcal B=\{B:B\subt[D_{z}]{T(\sigma)}D^{(1)}_{z}\}$, the reduction
$D^{(B,\top)}$ that equals $D^{(1)}$ off $z$ and $B$ at $z$, and the maximal
$1$-consistent reduction $D^{(B,\bot)}\subseteq D^{(B,\top)}$ (possibly empty).
Lemma~\ref{lem:tree-coverings} supplies tree-coverings $\Upsilon_{B,x}$
computing $D^{(B,\bot)}_{x}$, and their conjunction $\Upsilon_{x}$ works
uniformly in $B$. Writing $\mathcal B_{0}$ for the $B$ with $D^{(B,\bot)}$
nonempty, the argument splits:
\begin{itemize}\tightlist
\item $\mathcal B_{0}=\varnothing$: the type must be $\TL$, and
  Lemma~\ref{lem:bridge-from-instance} applied to a suitably weakened
  tree-covering produces, for any two constraints sharing a variable, a
  reflexive bridge between their $\ConOne$'s. Hence $\mathcal I$ is connected,
  and (1b) or (1c) holds according as its solution set is subdirect.
\item $\mathcal B_{0}\neq\varnothing$: one constructs a family $\Omega$ of
  weakenings of $\mathcal I$ that is \emph{critically} unable to solve all of
  $\mathcal B_{0}$ (conditions 1--4 of
  Construction~\ref{constr:omega}) and argues by whether some member fails
  (1b).
\end{itemize}
\end{proof}

\begin{hazard}[What $\Omega$'s construction turns on]\label{haz:omega}
Condition 3 is a minimality condition, and its statement is where the work is.
Read as ``if we replace any instance in $\Omega$ by \emph{all} weaker instances
then 2 is not satisfied'' it inherits the defect of
Caveat~\ref{haz:weakening}: the replacement need not make progress, and the
iteration that is supposed to produce $\Omega$ has no measure. Read as in Proposition~\ref{constr:omega} --- one member replaced by the
finitely many $\mathcal J[C\Uparrow]$, and termination measured on the
constraints rather than on the solution set --- both problems go away at once,
and the derivation of condition 4 becomes exact rather than approximate,
because what condition 3 then hands over is verbatim
Definition~\ref{def:crucial}.

Two things about the construction are worth keeping in view.

\begin{enumerate}\alphlabels\tightlist
\item \emph{The counting argument is used twice, one level apart.}
  Lemma~\ref{lem:weakening-terminates} orders \emph{instances} by counting
  their constraints by $\nu$-value; Proposition~\ref{constr:omega} orders
  \emph{sets of instances} by counting their members. Both are the multiset
  extension of a well-founded order, written as a lexicographic order on a
  product of copies of $\mathbb N$ indexed by a finite set, which is the same
  argument in a form whose well-foundedness is visible without citing
  Dershowitz--Manna.
\item \emph{Condition 4's witness depends on $\mathcal J$.} The $B$ produced
  varies with the member, so condition 4 is a $\Pi\Sigma$ statement: ``for
  every $\mathcal J$ there is $B$'' may not be strengthened to a single $B$
  serving all of $\Omega$, and the consumers in \S\ref{sec:main-statements}
  pick a fresh $B$ for each $\mathcal J$ accordingly.
\item \emph{A second descending sequence sits inside the first.} Subcase~1 of
  the $\mathcal B_{0}\neq\varnothing$ branch builds
  $\mathcal B_{1}\supseteq\mathcal B_{2}\supseteq\dots$ by repeatedly applying
  the inductive hypothesis, terminating because $\mathcal B_{1}$ is finite and
  the sequence is strictly decreasing. That is a different and much easier
  well-foundedness argument than
  Proposition~\ref{constr:omega}'s, and the two are worth keeping apart:
  nested constructions of this shape are where informal proofs most often hide
  a circularity, and neither of these does.
\end{enumerate}
\end{hazard}

\begin{longproposition}[Construction of $\Omega$]\label{constr:omega}
Let $\mathcal I$ be crucial in $D^{(1)}$ with $\mathcal B_{0}\neq\varnothing$,
and for a weakening $\mathcal J$ of $\mathcal I$ put
$\Sol(\mathcal J)=\{B\in\mathcal B_{0}:\mathcal J^{(B,\bot)}\text{ has a
solution}\}$. Then there is a finite set $\Omega$ of instances with:
\begin{enumerate}\tightlist
\item every member is a weakening of $\mathcal I$;
\item $\bigcap_{\mathcal J\in\Omega}\Sol(\mathcal J)=\varnothing$;
\item for every $\mathcal J\in\Omega$, the set
  $\bigl(\Omega\setminus\{\mathcal J\}\bigr)\cup
   \{\mathcal J[C\Uparrow]:C\in\mathcal J\}$
  does \emph{not} satisfy 2;
\item for every $\mathcal J\in\Omega$ there is $B\in\mathcal B_{0}$ with
  (a) $\mathcal J$ crucial in $D^{(B,\bot)}$ and (b) $B\in\Sol(\mathcal J')$
  for every $\mathcal J'\in\Omega\setminus\{\mathcal J\}$.
\end{enumerate}
\end{longproposition}

\begin{proof}
\emph{Monotonicity.} If $\mathcal J'$ is weaker than $\mathcal J$ then every
solution of $\mathcal J^{(B,\bot)}$ restricts to one of
$\mathcal J'^{(B,\bot)}$, so $\Sol(\mathcal J)\subseteq\Sol(\mathcal J')$.

\emph{Conditions 1--3.} Start with $\Omega=\{\mathcal I\}$. Condition 1 is
clear, and condition 2 holds because $D^{(B,\bot)}\subseteq D^{(1)}$ and
$\mathcal I^{(1)}$ has no solution, so $\Sol(\mathcal I)=\varnothing$. While
condition 3 fails at some $\mathcal J\in\Omega$ --- that is, while the
displayed set still satisfies 2 --- replace $\Omega$ by that set. Condition 1
survives, since a single weakening step applied to a weakening of $\mathcal I$
is again one; condition 2 survives by the choice of the step.

\emph{Termination.} Each $\mathcal J[C\Uparrow]$ is one weakening step from
$\mathcal J$, so by Lemma~\ref{lem:weakening-terminates}(c) it is strictly
below $\mathcal J$ in the order that lemma provides on the set $W$ of
weakenings of $\mathcal I$; and $W$ is finite by
Lemma~\ref{lem:weakening-terminates}(b). Fix a linear order on $W$ extending
it, and for a finite $\Omega\subseteq W$ let $c(\Omega)\in\mathbb N^{W}$ count
how many members $\Omega$ has at each element. One replacement removes a member
and inserts members strictly below it, so the coordinate there drops by one,
the coordinates above are untouched, and only coordinates below grow. Ordering
$\mathbb N^{W}$ lexicographically from the largest element of $W$ downwards ---
again a lexicographic order on a finite product of copies of $\mathbb N$ ---
every replacement strictly decreases $c(\Omega)$, so the loop halts at an
$\Omega$ satisfying 1--3.

This is the same counting argument twice over: once inside
Lemma~\ref{lem:weakening-terminates} to order instances by their constraints,
and once here to order sets of instances by their members.

\emph{Condition 4.} Fix $\mathcal J\in\Omega$. By condition 3 the displayed set
fails 2, so some $B\in\mathcal B_{0}$ lies in all of its members' $\Sol$. Then
$B\in\Sol(\mathcal J')$ for every $\mathcal J'\in\Omega\setminus
\{\mathcal J\}$, which is 4(b); and $B\notin\Sol(\mathcal J)$, since otherwise
$B$ would lie in $\bigcap_{\mathcal K\in\Omega}\Sol(\mathcal K)$, empty by
condition 2. So $\mathcal J^{(B,\bot)}$ has no solution while
$\mathcal J[C\Uparrow]^{(B,\bot)}$ has one for every constraint $C$ of
$\mathcal J$ --- which is Definition~\ref{def:crucial}, once no constraint of
$\mathcal J$ has a dummy variable.

None does. If $C$ had one, its projection $C'$ off that variable is weaker than
$C$ with the same violation set, and $C'$ is among the constraints of
$\mathcal J[C\Uparrow]$; so every solution of $\mathcal J[C\Uparrow]$ satisfies
$C'$ and hence $C$, giving $\Sol(\mathcal J[C\Uparrow])=\Sol(\mathcal J)$ and
so $B\in\Sol(\mathcal J)$. Finally $\mathcal J$ has at least one constraint,
since a weakening step replaces a constraint by a nonempty family rather than
deleting one, and $\mathcal I$ is crucial.
\end{proof}

\subsection{The three consumed theorems}\label{sec:codim-one}

\begin{theorem}[Reductions are safe]\label{thm:reductions-safe}
Let $\Theta$ be a cycle-consistent irreducible instance and
$B\subt[D_{x}]{T}D_{x}$ with $T\in\{\TBA,\TC,\TPC\}$. Then $\Theta$ has a
solution if and only if $\Theta$ has a solution with $x\in B$.
\end{theorem}

\begin{proof}
For $T=\TPC$ this is Theorem~\ref{thm:pc-safe}. For $T\in\{\TBA,\TC\}$,
Lemma~\ref{lem:find-consistent} gives a $1$-consistent reduction
$D^{(1)}\subte{T}D$ with $D^{(1)}_{x}\subseteq B$; by
Theorem~\ref{thm:main-inductive}(2), $\Theta^{(1)}$ has a solution, hence
$\Theta$ has one with $x\in B$.
\end{proof}

\begin{theorem}[PC reductions do not kill all solutions]\label{thm:pc-safe}
If $\mathcal I$ is cycle-consistent and irreducible, $B\subt[D_{y}]{\TPC(\sigma)}D_{y}$
for some $y\in\Var(\mathcal I)$, and $\mathcal I$ has a solution, then
$\mathcal I$ has a solution with $y\in B$.
\end{theorem}

\begin{lemma}[Subuniverses of these products are cosets]\label{lem:subuniverse-coset}
Let $q_{1},\dots,q_{r}$ be primes with $\mathbf Z_{q_{j}}\in\Vn$ and let
$\varnothing\neq L\le\mathbf Z_{q_{1}}\times\dots\times\mathbf Z_{q_{r}}$. Then
$L=g+M$ for some $g\in L$ and some subgroup $M$ of the product.
\end{lemma}

\begin{proof}
By Proposition~\ref{prop:p-divides}, $q_{j}\mid n-1$ for every $j$ and $w$ acts
on each factor as $x_{1}+\dots+x_{n}$; so $w$ acts on the product as the
$n$-fold sum, and $(n-1)h=0$ for every $h$ in the product. Fix $g\in L$ and put
$M=L-g$. If $\bar y_{1},\dots,\bar y_{n}\in M$ then
$w(\bar y_{1}+g,\dots,\bar y_{n}+g)=\sum_{j}\bar y_{j}+ng=
\bigl(\sum_{j}\bar y_{j}\bigr)+g$, which lies in $L$; so $M$ is closed under
the $n$-fold sum. Now $0\in M$, and for $y,z\in M$ taking
$(y,z,0,\dots,0)$ --- legitimate since $n\ge3$ --- gives $y+z\in M$. A finite
nonempty subset of a group closed under addition is a subgroup.
\end{proof}

\begin{longtheorem}[Codimension one]\label{thm:codim-one}
Suppose
\begin{enumerate}\romanlabels\tightlist
\item $\mathcal I$ is a linked cycle-consistent irreducible instance with
  $\Var(\mathcal I)=\{x_{1},\dots,x_{m}\}$;
\item $\mathbf D_{x_{i}}$ is $\TS$-free for every $i$;
\item $\mathcal I[C\Uparrow]$ has subdirect solution set for every constraint
  $C$ of $\mathcal I$;
\item $\sigma_{x_{i}}$ is the intersection of all linear congruences $\sigma$
  on $\mathbf D_{x_{i}}$ with $\cover{\sigma}=D_{x_{i}}^{2}$;
\item $L_{x_{i}}=\mathbf D_{x_{i}}/\sigma_{x_{i}}$;
\item $\phi\colon\mathbf Z_{q_{1}}\times\dots\times\mathbf Z_{q_{k}}\to
  L_{x_{1}}\times\dots\times L_{x_{m}}$ is a homomorphism, $q_{i}$ prime;
\item $\mathcal I[C\Uparrow]$ has a solution in $\phi(\bar a)$ for every
  constraint $C$ of $\mathcal I$ and every
  $\bar a\in\mathbf Z_{q_{1}}\times\dots\times\mathbf Z_{q_{k}}$ --- that is,
  $\mathcal I$ is crucial in every $\phi(\bar a)$
  (Definition~\ref{def:crucial}).
\end{enumerate}
Then
$\Delta=\{\bar a:\mathcal I\text{ has a solution in }\phi(\bar a)\}$ is empty,
full, or a coset of the kernel of a surjective homomorphism
$\lambda\colon\mathbf Z_{q_{1}}\times\dots\times\mathbf Z_{q_{k}}\to
\mathbf Z_{p}$ --- the solution set of the single equation
$\lambda(\bar a)=d$. Such a $\lambda$ vanishes on every factor with
$q_{i}\neq p$, so the equation involves only the coordinates over
$\mathbb F_{p}$.
\end{longtheorem}

\begin{proof}[Proof]
If $\Delta$ is full we are done. Otherwise pick $\bar b\notin\Delta$ and regard
$\phi(\bar b)$ as a reduction $D^{(1)}$ for $\mathcal I$; condition (vii) says
$\mathcal I$ is crucial in $D^{(1)}$.

\emph{Step 1: some $\ConOne(C,x)$ is a perfect linear congruence.} By
Lemma~\ref{lem:minimal-consistent}, either $C_{0}^{(1)}=\varnothing$ for some
constraint $C_{0}$, or $D^{(1)}$ is $1$-consistent.
\begin{itemize}\tightlist
\item If $C_{0}^{(1)}=\varnothing$: cruciality forces $\mathcal I=\{C_{0}\}$,
  say $C_{0}=R(y_{1},\dots,y_{t})$. By
  Lemma~\ref{lem:parallelogram-crucial}, $R$ has the parallelogram property and
  $\ConOne(R,1)$ is linear with $\cover{\ConOne(R,1)}=D_{y_{1}}^{2}$. By
  Lemma~\ref{lem:linear-on-top}, $\mathbf D_{y_{1}}/\delta\cong\mathbf Z_{p}$;
  with $\psi$ the quotient map, $\zeta=\{(a_{1},a_{2},b):\psi(a_{1})-\psi(a_{2})=b\}$
  witnesses perfection.
\item If $D^{(1)}$ is $1$-consistent: by Theorem~\ref{thm:main-inductive}(1)
  every constraint has the parallelogram property and (1b) or (1c) holds. In
  either case Lemma~\ref{lem:connected}(p) applies --- under (1c) to the
  linked connected subinstance, whose containing a constraint of $\mathcal I$
  is forced by condition (iii); under (1b) to $\mathcal I$ itself.
\end{itemize}

\emph{Step 2: add a $\mathbf Z_{p}$ coordinate.} Let $\zeta$ witness
perfection of $\ConOne(C,x)$. Add a variable $z$ with domain $\mathbf Z_{p}$,
replace $x$ in $C$ by a fresh $x'$, and add $\zeta(x,x',z)$, giving
$\mathcal I'$. Put
\[
  L=\{(\bar a,b):\mathcal I'\text{ has a solution with }z=b\text{ in }\phi(\bar a)\}
   \le\mathbf Z_{q_{1}}\times\dots\times\mathbf Z_{q_{k}}\times\mathbf Z_{p},
\]
and write $G=\mathbf Z_{q_{1}}\times\dots\times\mathbf Z_{q_{k}}$ and
$\pi\colon G\times\mathbf Z_{p}\to G$ for the projection. Condition (vii) gives
$\pi(L)=G$; in particular $L\neq\varnothing$. And $(\bar b,0)\notin L$, since a
solution of $\mathcal I'$ with $z=0$ in $\phi(\bar b)$ would, by the defining
property of $\zeta$, yield a solution of $\mathcal I$ in $\phi(\bar b)$, and
$\bar b\notin\Delta$.

\emph{Step 3: $L$ is a graph.} By Lemma~\ref{lem:subuniverse-coset},
$L=(\bar g,c)+M$ for some $(\bar g,c)\in L$ and some subgroup $M\le
G\times\mathbf Z_{p}$; and $\pi(M)=\pi(L)-\bar g=G$. The subgroup
$M\cap(\{0\}\times\mathbf Z_{p})$ is a subgroup of $\mathbf Z_{p}$, so it is
trivial or the whole of $\mathbf Z_{p}$. If it were the whole of
$\mathbf Z_{p}$ then $L$ would be closed under changing its last coordinate
arbitrarily, and $\pi(L)=G$ would then give $(\bar b,0)\in L$, which Step 2
excludes. So $M\cap(\{0\}\times\mathbf Z_{p})=0$, $\pi|_{M}$ is an isomorphism
$M\to G$, and $L$ is the graph of the affine map
\[
  f\colon G\to\mathbf Z_{p},\qquad
  f(\bar a)=c+\lambda(\bar a-\bar g),
\]
where $\lambda\colon G\to\mathbf Z_{p}$ is the homomorphism reading off the last
coordinate of $(\pi|_{M})^{-1}$.

\emph{Step 4: read off $\Delta$.} By the defining property of $\zeta$, an
$\bar a$ lies in $\Delta$ exactly when $\mathcal I'$ has a solution with $z=0$
in $\phi(\bar a)$, that is, exactly when $(\bar a,0)\in L$; so
$\Delta=f^{-1}(0)$. If $\lambda=0$ then $f$ is constant and $\Delta$ is empty
or full according as $c-\lambda(\bar g)\neq0$ or $=0$. Otherwise $\lambda$ is
onto $\mathbf Z_{p}$, since its image is a nontrivial subgroup of
$\mathbf Z_{p}$, and $\Delta$ is the coset of $\ker\lambda$ on which $f$
vanishes --- the solution set of one equation $\lambda(\bar a)=d$. Finally
$\lambda$ kills every factor $\mathbf Z_{q_{i}}$ with $q_{i}\neq p$: an element
of that factor has order dividing $q_{i}$, and its image has order dividing
$\gcd(q_{i},p)=1$. So the equation involves only the coordinates over
$\mathbb F_{p}$.
\end{proof}

\begin{hazard}[Why Step 3 is not one line]\label{haz:step2-linear-algebra}
An earlier draft compressed Steps 2--4 into ``$L$ has dimension $k$ and is the
solution set of one linear equation''. Two things make that false as stated.

$L$ is a nonempty subuniverse of an idempotent affine algebra, hence in general
an \emph{affine coset} and not a subgroup; Lemma~\ref{lem:subuniverse-coset} is
what supplies the translation, and it is where $n\ge3$ and
$q\mid n-1$ are spent. And a product of prime fields with \emph{different}
primes is not one vector space, so ``dimension $k$'' and ``codimension $1$''
have no meaning until the primary components are separated. What survives is
not a dimension count but the statement that $\Delta$ is a coset of the kernel
of a surjection onto $\mathbf Z_{p}$; that kernel has index $p$, which is the
only sense in which the codimension is one. The footnote in
\cite{ZhukSimplified} about ``$J$ may contain only variables on the same
field'' is the same issue surfacing in the algorithm, and \textup{(p3)} of
\S\ref{sec:algorithm} is organised prime by prime for exactly this reason.

What Step 2 still takes from elsewhere is $\pi(L)=G$, which is condition (vii)
read through the construction of $\mathcal I'$, and the property of $\zeta$
that relates solutions of $\mathcal I'$ with $z=0$ to solutions of
$\mathcal I$. Both belong to the half of \cite{ZhukSimplified} this document
has not yet read adversarially; see \S\ref{sec:status}.
\end{hazard}

\subsection{The supporting lemmas}

The lemmas invoked above, in dependency order.

\begin{lemma}[{\cite[Lem.~6.1]{ZhukJACM}}]\label{lem:expanded-consistency}
An expanded covering of a cycle-consistent irreducible instance is
cycle-consistent and irreducible.
\end{lemma}

\begin{lemma}\label{lem:minimal-consistent}
Suppose $D^{(1)}$ is a $1$-consistent reduction for $\mathcal I$, every
$D^{(1)}_{x}$ is $\TS$-free, $T\in\{\TPC,\TL,\TD\}$, $D^{(1)}\lll D$, and
$D^{(2)}_{x}\subte{\MT T}D^{(1)}_{x}$ is a minimal $\MT T$ subuniverse for
every $x$. Then either $C^{(2)}=\varnothing$ for some constraint $C$, or
$\mathcal I^{(2)}$ is $1$-consistent.
\end{lemma}

\begin{lemma}\label{lem:crucial-irreducible}
Let $R(x_{1},\dots,x_{n})$ be a \emph{rectangular} constraint of a
$1$-consistent instance $\mathcal I$, crucial in a reduction $D^{(\top)}$. Then
$\ConOne(R,i)$ is an irreducible congruence for every $i\in[n]$.
\end{lemma}

\begin{proof}
Take $i=1$; the other coordinates are the same. Suppose
$\ConOne(R,1)=\omega_{1}\cap\omega_{2}$ with
$\ConOne(R,1)\lneq\omega_{j}\le\mathbf D_{x_{1}}^{2}$ for $j=1,2$, and put
\[
  R_{j}(x_{1},\dots,x_{n})=\exists y\;
    \bigl(R(y,x_{2},\dots,x_{n})\wedge\omega_{j}(y,x_{1})\bigr).
\]
Each $R_{j}$ is a constraint on the same scope with
$R_{j}\supseteq R$, and the containment is proper because
$\omega_{j}\supsetneq\ConOne(R,1)$; so each $R_{j}$ is weaker than $R$ and is
among the constraints of $\mathcal I[R\Uparrow]$. And
$R_{1}\cap R_{2}=R$: a tuple in both has, for each $j$, a witness $y_{j}$ with
$R(y_{j},x_{2},\dots,x_{n})$ and $(y_{j},x_{1})\in\omega_{j}$, and rectangularity
of the first variable makes any two such witnesses $\ConOne(R,1)$-related to
each other, so $(y_{1},x_{1})\in\omega_{1}\cap\omega_{2}=\ConOne(R,1)$ and
$x_{1}$ may replace $y_{1}$.

Hence every solution of $\mathcal I[R\Uparrow]^{(\top)}$ satisfies both $R_{1}$
and $R_{2}$, so it satisfies $R$ and solves $\mathcal I^{(\top)}$ --- which has
no solution. That contradicts cruciality.
\end{proof}

\begin{hazard}[Rectangularity is a hypothesis here, not a
  conclusion]\label{haz:crucial-irreducible}
An earlier draft of this lemma read ``a constraint crucial in a reduction of a
cycle-consistent irreducible instance is rectangular, and its $\ConOne$'s are
irreducible''. That runs two statements together and gets the second one's
logic backwards: rectangularity is what the argument \emph{uses}, to know that
two witnesses for the same tuple are $\ConOne(R,1)$-related. Rectangularity of
crucial constraints is a separate fact, and it comes from the parallelogram
property (Lemma~\ref{lem:parallelogram-crucial}), not from cruciality alone.

The proof is also the reason Definition~\ref{def:crucial} quantifies over
$\mathcal I[R\Uparrow]$ rather than over single weaker constraints: what it
contradicts is that $R_{1}$ and $R_{2}$ can be imposed \emph{together}, and
each of them alone is a weakening that does gain a solution.
\end{hazard}

\begin{lemma}[A bridge from a relation]\label{lem:bridge-from-relation}
Let $R\sdle\mathbf A_{1}\times\dots\times\mathbf A_{n}$ have its first and last
variables rectangular, and suppose there are
$(b_{1},a_{2},\dots,a_{n})\in R$ and $(a_{1},\dots,a_{n-1},b_{n})\in R$ with
$(a_{1},a_{2},\dots,a_{n})\notin R$. Then
\[
  \delta(x_{1},x_{2},y_{1},y_{2})=\exists z_{2}\dots\exists z_{n-1}\;
    R(x_{1},z_{2},\dots,z_{n-1},y_{1})\wedge
    R(x_{2},z_{2},\dots,z_{n-1},y_{2})
\]
is a bridge from $\ConOne(R,1)$ to $\ConOne(R,n)$ with
$\bcol\delta=\proj_{1,n}(R)$.
\end{lemma}

\begin{proof}
Clauses (a) and (b) of Definition~\ref{def:bridge} are immediate from the
definition of $\ConOne$. For clause (d), rectangularity of the first and last
variables says exactly that $(x_{1},x_{2})\in\ConOne(R,1)$ if and only if
$(y_{1},y_{2})\in\ConOne(R,n)$ for tuples sharing the middle block
$z_{2},\dots,z_{n-1}$. Clause (c) is witnessed by
$(b_{1},a_{1},a_{n},b_{n})\in\delta$ --- take
$z_{i}=a_{i}$, so that the two conjuncts are the two tuples of the hypothesis
--- together with $(b_{1},a_{1})\notin\ConOne(R,1)$, which is what
$(a_{1},\dots,a_{n})\notin R$ says. Finally $\bcol\delta$ is $\proj_{1,n}(R)$,
since $(x,x,y,y)\in\delta$ asks for a single tuple of $R$ with first entry $x$
and last entry $y$.
\end{proof}

\begin{lemma}[Connected instances]\label{lem:connected}
Let $\mathcal I$ be a cycle-consistent connected instance. Then
\begin{enumerate}\tightlist
\item[\textup{(a)}] any two constraints with a common variable are adjacent;
\item[\textup{(b)}] for constraints $C_{1},C_{2}$, variables
  $x_{i}\in\Var(C_{i})$, and any path from $x_{1}$ to $x_{2}$, there is a
  bridge $\delta$ from $\ConOne(C_{1},x_{1})$ to $\ConOne(C_{2},x_{2})$ with
  $\bcol{\delta}$ containing all pairs connected by that path;
\item[\textup{(p)}] if $\mathcal I$ is linked then $\ConOne(C,x)$ is a perfect
  linear congruence for every $C\in\mathcal I$, $x\in\Var(C)$.
\end{enumerate}
\end{lemma}

\begin{hazard}\label{haz:connected-p}
Clause (p) is the bridge between the bridge machinery and the linear algebra;
it is what makes Theorem~\ref{thm:codim-one} possible. Its proof is (b)
followed by Lemma~\ref{lem:perfect-from-linked}: linkedness of the instance
makes $\bcol{\delta}$ linked as a relation, which upgrades the bridge. The two
senses of ``linked'' (Caveat~\ref{haz:linked}) meet exactly here, and the step
where one becomes the other should be isolated.
\end{hazard}

\begin{lemma}[{\cite[Lem.~5.6]{ZhukStrong}}]\label{lem:tree-coverings}
For a $1$-consistent instance and a reduction, there is a tree-covering whose
projection onto a given variable computes the maximal $1$-consistent
subreduction at that variable.
\end{lemma}

\begin{lemma}\label{lem:find-consistent}
Let $D^{(1)}$ be a $1$-consistent reduction of a cycle-consistent instance
$\mathcal I$ with $D^{(1)}\lll D$, and $B\subt[D_{x}]{T}D^{(1)}_{x}$ with
$T\in\{\TBA,\TC,\TPC\}$. Then there is a nonempty $1$-consistent reduction
$D^{(2)}\lll D^{(1)}$ with $D^{(2)}_{x}\subseteq B$. Moreover, if
$T\in\{\TBA,\TC\}$ then $D^{(2)}\subte{T}D^{(1)}$; and if $T=\TPC$ and every
$D^{(1)}_{y}$ is $\TS$-free then $D^{(2)}\subte{\TMPC}D^{(1)}$.
\end{lemma}

\begin{lemma}\label{lem:parallelogram-crucial}
Suppose $D^{(1)}$ is a $1$-consistent reduction for the one-constraint instance
$R(x_{1},\dots,x_{n})$, that $T\in\{\TL,\TPC,\TD\}$, that
$D^{(2)}\subte[D]{\MT T}D^{(1)}\lll D$, and that $R$ is crucial in $D^{(2)}$.
Then $R$ has the parallelogram property, and for every $i\in[n]$ the congruence
$\ConOne(R,i)$ is of type $T$ with
$\cover{\ConOne(R,i)}\supseteq(D^{(1)}_{x_{i}})^{2}$. If moreover $T=\TPC$ then
$n=2$.
\end{lemma}

\begin{longlemma}[Bridges from a subdirect PC/linear instance]\label{lem:bridge-from-instance}
Suppose
\begin{enumerate}\romanlabels\tightlist
\item $\mathcal I$ has subdirect solution set;
\item $D^{(1)}$ is a reduction with $D^{(1)}_{x}\lll D_{x}$ for every $x$;
\item $C\in\mathcal I$ is a constraint of type $T\in\{\TPC,\TL\}$;
\item $B\subt[D_{z}]{\mathcal T(\xi)}D^{(1)}_{z}$ for some variable $z$, with
  $\mathcal T\in\{\TBA,\TC,\TS,\TPC,\TL\}$;
\item if $T=\TPC$ then $\mathcal T\in\{\TPC,\TL\}$;
\item $\mathcal I^{(1)}$ has a solution;
\item $\mathcal I^{(1)}$ has no solution with $z\in B$;
\item $\mathcal I[C\Uparrow]$ has a solution in $D^{(1)}$ with $z\in B$.
\end{enumerate}
Then $\mathcal T=T$, and for every variable $x$ of $C$ there is a bridge
$\delta$ from $\xi$ to $\ConOne(C,x)$ with
$\bcol{\delta}\supseteq\mathcal I(z,x)$.
\end{longlemma}

\begin{hazard}[Stated only]\label{haz:bridge-from-instance}
Lemma~\ref{lem:bridge-from-instance} carries eight numbered hypotheses and is
\emph{stated without proof here}. It is the only source of bridges in
\S\ref{sec:main-statements}, and every use of connectedness routes through it;
so are Lemmas~\ref{lem:minimal-consistent}, \ref{lem:crucial-irreducible},
\ref{lem:bridge-from-relation}, \ref{lem:find-consistent},
\ref{lem:parallelogram-crucial}, Corollary~\ref{cor:same-type} and clauses
(a) and (b) of Lemma~\ref{lem:connected}. These are gaps in the logical chain
of this document, not imported results, and \S\ref{sec:status} records them as
such.

Its conclusion has two independent parts --- ``$\mathcal T=T$'' and ``there is a
bridge with $\bcol\delta\supseteq\mathcal I(z,x)$'' --- and different call
sites use different parts. The dependency order is
Theorem~\ref{thm:stable-intersection}, which it consumes, then this lemma, then
Theorem~\ref{thm:main-inductive}, which consumes it.
\end{hazard}

\begin{corollary}\label{cor:same-type}
In the situation of Theorem~\ref{thm:main-inductive}(2) with $\mathcal I'$
satisfying (1b), the type $\mathcal T$ of the reduction and the type of the
constraints must agree, which contradicts $\mathcal T\in\{\TBA,\TC\}$.
\end{corollary}

\clearpage
%!TeX root=csp-proof.tex
\section{The algorithm}\label{sec:algorithm}

The algorithm is \cite{ZhukJACM}'s. It is stated here because target
\textup{(T1)} of \S\ref{sec:target} is a proposition \emph{about} it and
cannot be written down otherwise.

\subsection{\textsc{Solve}}

\begin{verbatim}
function Solve(I):
    repeat
        I := ForceConsistency(I)
        if I = false: return "No"
        if some D_x has a strong subset B:
            I := ReduceDomain(I, x, B)
    until nothing changed
    return SolveLinear(I)
\end{verbatim}

\textsc{ForceConsistency} enforces cycle-consistency and irreducibility,
returning \texttt{false} if some domain empties. The reduction step is licensed
by Theorem~\ref{thm:reductions-safe}: restricting to a strong subset cannot
destroy solvability. When no strong subset remains, Lemma~\ref{lem:ubiquity}
together with Lemma~\ref{lem:linear-on-top} gives, for each domain of size at
least $2$, a congruence $\sigma_{x}$ with $D_{x}/\sigma_{x}\cong\mathbb Z_{q_{1}}
\times\dots\times\mathbb Z_{q_{n_{x}}}$; choose $\sigma_{x}$ minimal. That case
is \textsc{SolveLinear}.

\subsection{\textsc{SolveLinear}}

Write $\phi\colon\mathbb Z_{p_{1}}\times\dots\times\mathbb Z_{p_{k}}\to
D_{x_{1}}/\sigma_{x_{1}}\times\dots\times D_{x_{m}}/\sigma_{x_{m}}$ for a linear
map and
\[
  \phi^{-1}(\mathcal I)=\{\bar a:\mathcal I\text{ has a solution in }\phi(\bar a)\}.
\]
Computing $\phi^{-1}(\mathcal I)$ would solve $\mathcal I$; the algorithm cannot
do that, but it can do four things:

\begin{description}\tightlist
\item[\textup{(p0)}] decide $\bar a\in\phi^{-1}(\mathcal I)$, by reducing each
  domain to the corresponding block and recursing on a smaller domain;
\item[\textup{(p1)}] decide $\phi^{-1}(\mathcal I)=\mathbb Z_{p_{1}}\times\dots
  \times\mathbb Z_{p_{k}}$, by testing $\bar 0$ first and then the $k$ unit
  vectors with \textup{(p0)}. The set $\phi^{-1}(\mathcal I)$ is in general an
  affine \emph{coset}, not a subgroup; testing $\bar 0$ first is what makes it
  one, after which membership of the coordinate generators implies fullness;
\item[\textup{(p2)}] compute $\phi^{-1}(\mathcal I)$ when $\mathcal I$ is not
  linked, by splitting into linked pieces, recursing, and taking the union;
\item[\textup{(p3)}] compute $\phi^{-1}(\mathcal I)$ when it is empty or of
  codimension one, by finding $\bar a\notin\phi^{-1}(\mathcal I)$, then for each
  $i$ a $b_{i}$ with $\bar a[i\mapsto b_{i}]\in\phi^{-1}(\mathcal I)$ if one
  exists, and reading off the equation
  $\sum_{i\in J}(y_{i}-a_{i})/(b_{i}-a_{i})=1$ over the $i$ for which $b_{i}$
  exists. The displayed quotient is meaningful only within one prime: the
  codimension-one equation has a single target prime $p$, and only coordinates
  over $\mathbb F_{p}$ can carry a nonzero coefficient, so the computation is
  organised prime by prime (Caveat~\ref{haz:step2-linear-algebra}).
\end{description}

\begin{verbatim}
function SolveLinear(I):
    m := dim(D_{x_1}/s_1 x ... x D_{x_n}/s_n)
    phi := a bijective linear map Z_{p_1}x...xZ_{p_m} -> that product
    while phi^-1(I) != Z_{p_1}x...xZ_{p_m}:              // (p1)
        I' := I
        for C in I':                            // drop what can be dropped
            if phi^-1(I' \ {C}) != Z_{p_1}x...xZ_{p_m}:   // (p1)
                I' := I' \ {C}
        F := phi^-1(I')                                // (p2) or (p3)
        if F = {}: return "No"
        m := dim F ;  psi := a bijective linear map onto F ;  phi := phi o psi
    return "Yes"
\end{verbatim}

The inner loop makes $\mathcal I'$ minimal with $\phi^{-1}(\mathcal I')$ not
full. At that point Theorem~\ref{thm:codim-one} applies --- its hypothesis (vii)
is exactly minimality --- so $\phi^{-1}(\mathcal I')$ is empty or of codimension
one, and \textup{(p3)} computes it; or $\mathcal I'$ is not linked and
\textup{(p2)} does. Since $\phi^{-1}(\mathcal I)\subseteq\phi^{-1}(\mathcal I')$,
the image of $\phi$ shrinks while still containing all solutions. The dimension
strictly decreases, so the loop terminates.

\subsection{What \textup{(T1)} has to say}

\begin{hazard}[Two obstacles between the pseudocode and a function]
\label{haz:algorithm-obstacles}
\begin{enumerate}\alphlabels\tightlist
\item \emph{Termination is not structural.} \textsc{Solve} recurses through
  \textup{(p0)} on strictly smaller domains, and \textsc{SolveLinear} loops on
  a strictly decreasing dimension. Neither is a structural recursion; both need
  an explicit well-founded measure, and the two interleave ---
  \textsc{SolveLinear} calls \textsc{Solve} on a smaller domain, which calls
  \textsc{SolveLinear} again. The measure is lexicographic:
  $\bigl(\sum_{x}|D_{x}|,\ \dim\phi\bigr)$.
\item \emph{Several steps quantify over objects that are only finitely many by
  accident.} ``If some $D_{x}$ has a strong subset $B$'' quantifies over subsets
  of $D_{x}$, over congruences on $D_{x}$, and --- through absorption --- over
  \emph{term operations}, of which there are infinitely many. For the types the
  algorithm actually searches for, the search is finite for a reason that needs
  no general arity bound: $\TBA$ has a binary witness by definition, and a
  central subuniverse has a ternary one by Lemma~\ref{lem:central-ternary}.
\end{enumerate}
\end{hazard}

\subsection{\textsc{Solve} as a function}

The pseudocode is not a function: it makes unconstrained choices, and its
recursion is not structural. Both are fixed here, so that
Theorem~\ref{thm:t1} has a subject.

\begin{definition}[The call set and its measure]\label{def:solve-measure}
A \emph{call} is either $\mathsf S(\mathcal I)$, for $\mathcal I$ an instance
of $\CSP(\Gamma)$ with domains $D_{x}$, or $\mathsf L(\mathcal I,\phi)$, where
in addition $\phi$ is an injective linear map
$\mathbb Z_{p_{1}}\times\dots\times\mathbb Z_{p_{k}}\to\prod_{x}
D_{x}/\sigma_{x}$ with the $\sigma_{x}$ the minimal congruences of
\S\ref{sec:algorithm}. Put $N(\mathcal I)=\sum_{x}|D_{x}|$ and
\[
  \mu\bigl(\mathsf S(\mathcal I)\bigr)=
    \bigl(N(\mathcal I),\,N(\mathcal I)+1\bigr),
  \qquad
  \mu\bigl(\mathsf L(\mathcal I,\phi)\bigr)=
    \bigl(N(\mathcal I),\,\dim\phi\bigr),
\]
ordered lexicographically on $\mathbb N^{2}$. Since $\dim\phi\le N(\mathcal I)$
always, a call $\mathsf S(\mathcal I)$ is strictly above every
$\mathsf L(\mathcal I,\phi)$ with the same instance.
\end{definition}

\begin{definition}[$\textsc{Solve}_{\mathrm{alg}}$]\label{def:solve-alg}
Fix once and for all a linear order on the finite set of pairs $(x,B)$ with
$B\subseteq D_{x}$, and one on the finite set of candidate maps $\phi$; ``the
least'' below means least in these orders, which is what turns each choice in
the pseudocode into a value. Define $\textsc{Solve}_{\mathrm{alg}}$ on calls by
well-founded recursion on $\mu$:

\begin{description}\tightlist
\item[$\mathsf S(\mathcal I)$.] Put $\mathcal I'=
  \textsc{ForceConsistency}(\mathcal I)$. If some domain of $\mathcal I'$ is
  empty, return \texttt{false}. If some $D'_{x}$ has a proper nonempty subset
  $B$ with $B\subt{T}D'_{x}$ for some $T\in\{\TBA,\TC,\TD\}$, take the least
  such $(x,B)$ and return
  $\textsc{Solve}_{\mathrm{alg}}\bigl(\mathsf S(\textsc{ReduceDomain}(\mathcal
  I',x,B))\bigr)$. Otherwise return $\textsc{Solve}_{\mathrm{alg}}
  \bigl(\mathsf L(\mathcal I',\phi_{0})\bigr)$ for the least admissible
  $\phi_{0}$.
\item[$\mathsf L(\mathcal I,\phi)$.] Run the body of \textsc{SolveLinear},
  reading each test $\bar a\in\phi^{-1}(\mathcal I)$ of \textup{(p0)} as the
  value of $\textsc{Solve}_{\mathrm{alg}}$ at $\mathsf S$ of the instance
  $\mathcal I$ with each domain cut to the block $\phi(\bar a)_{x}$, and each
  iteration of the loop as the value at $\mathsf L(\mathcal I,\phi\circ\psi)$.
\end{description}
\end{definition}

\begin{lemma}[The recursion is well founded]\label{lem:solve-terminates}
Every recursive call made by Definition~\ref{def:solve-alg} is at a strictly
smaller $\mu$, so $\textsc{Solve}_{\mathrm{alg}}$ is a total function into
$\{\texttt{true},\texttt{false}\}$.
\end{lemma}

\begin{proof}
Four kinds of call. \textsc{ForceConsistency} only shrinks domains, so it does
not increase $N$; the reduction step replaces some $D'_{x}$ by a proper subset,
strictly decreasing $N$. The call from $\mathsf S(\mathcal I')$ to
$\mathsf L(\mathcal I',\phi_{0})$ leaves $N$ fixed and drops the second
component, by the inequality in Definition~\ref{def:solve-measure}. The
\textup{(p0)} calls cut every domain to a block of $\sigma_{x}$; each
$\sigma_{x}$ is proper at some variable, since $\mathsf L$ is reached only when
some domain has size at least $2$, so $N$ strictly decreases. The loop calls
leave $N$ fixed and strictly decrease $\dim\phi$, as
\S\ref{sec:algorithm} records.
\end{proof}

\begin{theorem}[Target \textup{(T1)}]\label{thm:t1}
For every instance $\mathcal I$ of $\CSP(\Gamma)$,
\[
  \textsc{Solve}_{\mathrm{alg}}\bigl(\mathsf S(\mathcal I)\bigr)=\texttt{true}
  \iff \mathcal I\ \text{has a solution}.
\]
\end{theorem}

\begin{hazard}[What Definition~\ref{def:solve-alg} still owes]\label{haz:t1-names-solve}
Theorem~\ref{thm:t1} is about \emph{this} function, and that is the whole of
its content: the weaker reading --- ``there exists a Boolean-valued function
deciding satisfiability'' --- follows for a fixed finite template by exhaustive
search and says nothing. Three things are needed before the definition is
airtight, and none is deep.

\begin{enumerate}\alphlabels\tightlist
\item \textsc{ForceConsistency} and \textsc{ReduceDomain} are used as given
  total functions. The first must be exhibited as a fixpoint iteration that
  terminates because it only shrinks domains; the second is immediate.
\item The search in $\mathsf S$ ranges over subsets $B$ and over the witnesses
  making $B\subt{T}D'_{x}$. It is finite for the reason in
  Caveat~\ref{haz:algorithm-obstacles}(b) --- $\TBA$ has a binary witness by
  definition and $\TC$ a ternary one by Lemma~\ref{lem:central-ternary}, and
  $\TD$ quantifies over congruences --- but the bound has to be stated for the
  ``least such $(x,B)$'' to be well defined.
\item Only $\TBA$, $\TC$ and $\TD$ are searched for, not $\TS$; that is
  legitimate because $\subt{\TS}$ is witnessed by a subset carrying the other
  two types, but the correctness proof must use
  Theorem~\ref{thm:reductions-safe} in the form that covers exactly the types
  searched.
\end{enumerate}

Theorem~\ref{thm:t1} itself is stated and not proved here; see
\S\ref{sec:status}.
\end{hazard}

Given that, Theorem~\ref{thm:t1} mentions no machine, is provable from
\S\ref{sec:main-statements}, and is the honest content of ``the algorithm
works'': that one particular, highly non-obvious algorithm --- the one whose
existence is the dichotomy --- is correct.

\begin{theorem}[Target \textup{(T2)}; stated, not proved]\label{thm:t2}
For each finite $\Gamma$ there are a program computing
$\textsc{Solve}_{\mathrm{alg}}$ and a
polynomial $p_{\Gamma}$ bounding its running time on instances of
$\CSP(\Gamma)$ in the size of the instance.
\end{theorem}

\begin{hazard}[\textup{(T2)} is not a wrapper]\label{haz:t2-scope}
Theorem~\ref{thm:t2} is the only statement of the tractable half that needs a
cost model, and it is not a corollary of Theorem~\ref{thm:t1}: it requires
$\textsc{Solve}_{\mathrm{alg}}$ to be exhibited as a program with a proved step
bound. Note
also that the statement is non-uniform. Because $\Gamma$ is fixed, the searches
over subsets of $A$, over congruences on $A$, and over term operations of
bounded arity are constant-time --- but ``constant'' means compiled away into a
program whose \emph{size} depends on $|A|$. So $\textsc{Solve}_{\mathrm{alg}}$
is a family with a polynomial for each $\Gamma$, and a bound
uniform in $\Gamma$ would be a different statement, and a false one.
\end{hazard}

\clearpage
%!TeX root=csp-proof.tex
\section{The hard half}\label{sec:hardness}

If no weak near-unanimity operation preserves $\Gamma$, then $\CSP(\Gamma)$ is
\textup{NP}-complete. This half is older, shorter, and better understood than
the tractable half, and it factors cleanly into algebra and complexity.

\subsection{The chain}

\begin{description}\tightlist
\item[\textup{(1)}] \emph{Reduction to a core, and to an idempotent algebra.}
  One may assume $\Gamma$ is a rigid core containing all singleton unary
  relations, so that $\Pol(\Gamma)$ is idempotent, at the cost of a
  logspace reduction \cite[``Cores and Idempotent Reducts'']{BradyNotes}.
  This is Caveat~\ref{haz:core}.
\item[\textup{(2)}] \emph{Maróti--McKenzie.} A finite idempotent algebra has a
  Taylor term if and only if it has a weak near-unanimity term
  \cite{MarotiMcKenzie}.
\item[\textup{(3)}] \emph{No Taylor term $\Rightarrow$ two disjoint
  subalgebras.} If a finite idempotent $\mathbf A$ has no Taylor term then
  there are nonempty $B,C\le\mathbf A$ with $B\cap C=\varnothing$ and
  \[
    N=(B\cup C)^{3}\setminus(B^{3}\cup C^{3})\;\le\;\mathbf A^{3}
  \]
  \cite[Cor.~1.5.11]{BradyNotes}. This replaces the usual formulation --- ``the
  variety generated contains a two-element algebra whose only operations are
  projections, and from there $\Gamma$ pp-interprets every finite language'' ---
  whose ``and'' hides an \textup{HSP} representation in a finite power and its
  translation into an interpretation. The cited corollary does that work and
  lands on a concrete relation, from which Lemma~\ref{lem:nae-interp} builds
  the interpretation in four lines.
\item[\textup{(3$'$)}] \emph{Inv--Pol.} $N$ is a subuniverse of $\mathbf A^{3}$
  and so is preserved by every polymorphism of $\Gamma$; Theorem~\ref{thm:galois}
  turns that into pp-definability from $\Gamma$. This is the one substantial
  import of the step, and only the hard direction of it is used.
\item[\textup{(4)}] \emph{Bulatov--Jeavons--Krokhin.} A primitive positive
  interpretation of $\Delta$ in $\Gamma$ yields a polynomial-time many-one
  reduction $\CSP(\Delta)\le_{p}\CSP(\Gamma)$ \cite{BJK}.
\item[\textup{(5)}] \emph{Cook--Levin.} $\textsc{Nae-$3$-Sat}$ is
  \textup{NP}-hard.
\end{description}

Steps (1)--(3) are algebra; step (4) is a syntactic construction plus a
complexity-theoretic wrapper; step (5) is pure complexity theory.

\subsection{The layered targets}

\begin{lemma}[The interpretation, explicitly]\label{lem:nae-interp}
Let $\mathbf A$ be an algebra, $B,C\le\mathbf A$ nonempty and disjoint, and
suppose $N=(B\cup C)^{3}\setminus(B^{3}\cup C^{3})$ is a subuniverse of
$\mathbf A^{3}$. Then $\{N\}$ pp-interprets $\{\textsc{Nae}\}$ on $\{0,1\}$,
with $d=1$.
\end{lemma}

\begin{proof}
A triple lies in $N$ exactly when its entries lie in $B\cup C$ and are not all
on the same side of the partition; in particular, for a repeated entry,
\begin{equation}\label{eq:nae-opposite}
  (x,z,z)\in N\iff x,z\in B\cup C\ \text{and $x,z$ lie on opposite sides.}
\end{equation}
Take $F=\proj_{1}(N)$, which is pp-definable from $N$ and equals $B\cup C$:
each $b\in B$ is the first entry of $(b,c,c)\in N$ for any $c\in C$, and
symmetrically. Let $\pi\colon F\to\{0,1\}$ send $B$ to $0$ and $C$ to $1$; it is
surjective because $B$ and $C$ are nonempty. The $\pi$-preimage of
$\textsc{Nae}$ is $N$ itself, by the description above. The $\pi$-preimage of
equality on $\{0,1\}$ is $B^{2}\cup C^{2}$, and by \eqref{eq:nae-opposite}
\[
  B^{2}\cup C^{2}=\{(x,y):\exists z\;N(x,z,z)\wedge N(y,z,z)\} :
\]
a $z$ opposite to both $x$ and $y$ exists exactly when $x$ and $y$ are on the
same side, since there are only two sides and both are nonempty.
\end{proof}

\begin{theorem}[Target \textup{(H0)}]\label{thm:h0}
Let $\mathbf A$ be a finite idempotent algebra with no weak near-unanimity term
operation. Then $\Inv(\mathbf A)$ primitively positively interprets the
constraint language $\{\textsc{Nae}\}$ on $\{0,1\}$, where
$\textsc{Nae}=\{0,1\}^{3}\setminus\{(0,0,0),(1,1,1)\}$.
\end{theorem}

\begin{proof}
By Lemma~\ref{lem:special-wnu}'s source \cite{MarotiMcKenzie}, a finite
idempotent algebra has a Taylor term if and only if it has a weak
near-unanimity term; so $\mathbf A$ has no Taylor term. Step \textup{(3)}
supplies nonempty disjoint $B,C\le\mathbf A$ with
$N=(B\cup C)^{3}\setminus(B^{3}\cup C^{3})\le\mathbf A^{3}$, so
$N\in\Inv(\mathbf A)$, and Lemma~\ref{lem:nae-interp} builds the
interpretation. Every relation it uses is pp-definable from $N$ and hence lies
in $\Inv(\mathbf A)$.
\end{proof}

\begin{theorem}[Target \textup{(H1)}; stated, not proved]\label{thm:h1}
If $\Gamma$ primitively positively interprets $\Delta$ then
$\CSP(\Delta)\le_{p}\CSP(\Gamma)$. Consequently, under \textup{(H0)},
$\CSP(\Gamma)$ is \textup{NP}-hard.
\end{theorem}

\begin{hazard}[Hardness is not completeness]\label{haz:np-membership}
Theorem~\ref{thm:h1} concludes \textup{NP}-\emph{hardness}. The word
\emph{complete} in Theorem~\ref{thm:dichotomy-core} additionally needs
membership, which is Theorem~\ref{thm:np-membership}. Small, but it needs the
cost model, so it belongs with \textup{(H1)} and not with \textup{(H0)}.
\end{hazard}

\begin{theorem}[\textup{NP} membership; stated, not proved]\label{thm:np-membership}
For each finite $\Gamma$, $\CSP(\Gamma)\in\textup{NP}$.
\end{theorem}

\begin{proof}[Proof sketch]
A solution is a certificate: an assignment of a value of $A$ to each variable,
of size $O(v\log|A|)$ for an instance with $v$ variables, hence linear in the
instance. Verifying it evaluates each constraint once, and for fixed $\Gamma$
that is a table lookup of constant cost. Writing this out needs the encoding
and the machine, which is why it is stated with \textup{(H1)} rather than
proved here.
\end{proof}

With Theorems~\ref{thm:h1} and \ref{thm:np-membership} together, the hard half
of Theorem~\ref{thm:dichotomy-core} reads \textup{NP}-\emph{complete} as
claimed.

Theorem~\ref{thm:h0} mentions no machine. Theorem~\ref{thm:h1} is where the
complexity layer enters, and it is small --- the reduction is a syntactic
substitution of formulas --- but it cannot be stated at all without a precise
notion of polynomial-time many-one reduction.

\subsection{Definitions}

\begin{definition}[pp-formula, pp-definability]\label{def:pp}
A \emph{primitive positive formula} over $\Gamma$ is
$\exists y_{1}\dots\exists y_{s}\;\bigwedge_{j}R_{j}(\bar z_{j})$ with
$R_{j}\in\Gamma\cup\{=\}$. A relation is \emph{pp-definable} from $\Gamma$ if
some pp-formula defines it.
\end{definition}

\begin{definition}[pp-interpretation]\label{def:pp-interp}
$\Gamma$ on $A$ \emph{pp-interprets} $\Delta$ on $B$ if there are $d\ge 1$, a
pp-definable $F\subseteq A^{d}$, and a surjection $\pi\colon F\to B$ such that
the $\pi$-preimage of every relation of $\Delta$, and of the equality relation
on $B$, is pp-definable from $\Gamma$.
\end{definition}

\begin{theorem}[Geiger; Bodnarchuk--Kalu\v znin--Kotov--Romov]\label{thm:galois}
For a finite $A$, a relation is pp-definable from $\Gamma$ if and only if it is
preserved by every polymorphism of $\Gamma$.
\end{theorem}

\begin{hazard}[The Galois connection is the expensive import]\label{haz:galois}
Theorem~\ref{thm:galois} is what converts the algebraic hypothesis (``no WNU
polymorphism'') into a syntactic conclusion (``pp-defines a hard relation''),
and it is the only genuinely substantial import of \S\ref{sec:hardness}. The
easy direction is a one-line calculation. The hard direction --- every relation
preserved by all polymorphisms is pp-definable --- goes through the free algebra
on $|A|^{k}$ generators and an indicator-instance construction; it is roughly
four printed pages and is the single largest item in the \textup{(H0)} budget.
\end{hazard}

\begin{remark}[Why the hard half is not the hard half]\label{rem:hardness-easier}
Despite its name, \textup{(H0)} is far smaller than \textup{(T0)}: perhaps
twenty printed pages against a hundred. The asymmetry is intrinsic. Showing a
problem is hard requires exhibiting one construction; showing it is easy
requires an algorithm that works on every instance, and the correctness proof
must survive every configuration the instance can present.
\end{remark}

\clearpage
%!TeX root=csp.tex
\section{Defects in the sources}\label{sec:defects}

A blueprint earns its keep by finding the places where the source cannot be
transcribed. This section collects them. They are not criticisms of
\cite{ZhukSimplified}, which is unusually careful prose --- across nine thousand
lines of \TeX{} it says ``obvious'' three times and ``clearly'' never. They are
the residue that any informal exposition leaves, and each one is a decision a
formalization has to make before it can write the corresponding statement.

Items marked \textbf{blocking} must be resolved before the affected statement
can be written in Lean at all. Items marked \textbf{convention} are places
where the source is consistent but leaves the reader to supply a reading.

\subsection{Blocking}

\begin{enumerate}\tightlist

\item[\textbf{D1}] \textbf{Lemma~\ref{lem:ba-center-implies} is false as
  printed.} \cite[Lemma~19]{ZhukSimplified} drops two hypotheses present in its
  cited source: subdirectness of $R$ in the first coordinate, and a common
  absorbing term in the $\TBA$ case. Counterexample and repair in
  Hazard~\ref{haz:ba-center-implies}. Used three times in \S3 and three times
  in \S4 of the source.

\item[\textbf{D2}] \textbf{Corollary~\ref{cor:prop-quotient} is never proved.}
  It is stated in \S2.3, used four times inside \S5 and once in \S3, and
  neither proved nor restated anywhere. It is Lemma~\ref{lem:propagation}
  specialised to the canonical surjection $\mathbf A\to\mathbf A/\delta$, but
  the source does not say so, and the specialisation is not quite mechanical:
  clause (m)'s hypothesis ``$B/\delta$ is $\TS$-free'' has to be matched against
  clause (fm)'s ``$f(B)$ is $\TS$-free''.

\item[\textbf{D3}] \textbf{There is a citation cycle in \S5.} Tracing the
  proofs: Lemma~9 $\Rightarrow$ Lemma~7 $\Rightarrow$ Lemma~86 $\Rightarrow$
  Lemma~85 $\Rightarrow$ Corollary~22 $\Rightarrow$ Theorem~21 $\Rightarrow$
  Lemma~8 $\Rightarrow$ Lemma~9, closed where Lemma~85's $T=\TC$ case invokes
  Corollary~22. Nothing in \S5.4 or \S5.7 can be formalized until the cycle is
  broken. The break appears to require a mixed essentiality lemma for three
  distinct central subuniverses, which is not in the paper.

\item[\textbf{D4}] \textbf{Theorem~\ref{thm:stable-intersection}(c) asserts
  $m=2$ without argument.} The proof reduces to pairs and shows the types agree;
  for $\TPC$, $\sigma_{1}=\sigma_{2}$ forces $m=2$. Nothing rules out $m\ge 3$
  with all types $\TC$. A Lean statement copying the source is unprovable as it
  stands. Same fix as D3.

\item[\textbf{D5}] \textbf{Theorem~\ref{thm:codim-one} hypothesis (ii) is too
  weak for its proof.} The source says ``$\mathbf D_{x_{i}}$ is $\TS$-free'';
  the proof, and the informal claim it formalises, need ``no nontrivial binary
  absorbing or central subuniverse on $\mathbf D_{x_{i}}$'' --- which is the
  \emph{commented-out} gloss sitting immediately below the hypothesis in the
  source. $\TS$-freeness is strictly weaker (Hazard~\ref{haz:S-type}).
  Strengthen the hypothesis.

\item[\textbf{D6}] \textbf{Lemma~\ref{lem:connected}(a) applies
  Lemma~\ref{lem:bridge-from-relation} without one of its hypotheses.} The
  missing hypothesis can fail; a witness is
  $R=\{(x,z,y)\in\mathbb Z_{4}\times\mathbb Z_{2}\times\mathbb Z_{4}:
  x\equiv z\equiv y \pmod 2\}$. In \cite{ZhukJACM} the needed tuples came from
  \emph{criticality}, a notion \cite{ZhukSimplified} dropped. The fix is to add
  ``every constraint relation is critical'' to the definition of a connected
  instance --- available, since crucial implies critical --- at the cost of
  reimporting about a page of the older machinery.

\item[\textbf{D7}] \textbf{Two gaps in Theorem~\ref{thm:main-inductive}.}
  In Case~1 the inductive hypothesis yields (1c) relative to $D^{(2)}$ while the
  goal is (1c) relative to $D^{(1)}$; the source writes only ``we derive the
  required conditions''. And both endgames conclude ``(1b) if the solution set
  is subdirect, (1c) otherwise'', but (1c) demands a \emph{linked} connected
  subinstance where only connectedness was established --- linkedness has to be
  extracted from irreducibility, which needs a fragmentation descent and a
  separate treatment of the empty-solution-set case.

\item[\textbf{D8}] \textbf{A lemma the proof needs is commented out of the
  source.} \texttt{LEMMinimalContainingIsMinimal} --- that an inclusion-minimal
  $\MT T$-subset containing a given element is minimal among $\MT T$-subsets ---
  is commented out along with its citation, yet is genuinely used in Case~2 of
  Theorem~\ref{thm:main-inductive}. It is true and must be reproved.

\item[\textbf{D9}] \textbf{The algorithm's pseudocode does not match the theorem
  it invokes.} \textsc{SolveLinear} \emph{removes} constraints, whereas
  Theorem~\ref{thm:codim-one}(vii) is about \emph{weakening} them; and
  primitive \textup{(p2)} returns a union of affine subspaces, which the next
  line treats as affine. \cite{ZhukSimplified} labels its pseudocode a sketch
  and points to \cite{ZhukJACM} for the precise algorithm, so this is a defect
  of the sketch rather than of the proof --- but target \textup{(T1)} is a
  statement about the algorithm, so it must be repaired from \cite{ZhukJACM}
  before \textup{(T1)} can be stated.

\item[\textbf{D10}] \textbf{The arity in Lemma~\ref{lem:special-wnu} is suspect.} The
  imported statement gives a special WNU of arity $n^{n!}$; the period of
  $y\mapsto w(x,\dots,x,y)$ need not divide $n^{n!}$ --- take $n=3$ and an
  element of order $5$. The remedy is to weaken the statement to ``\emph{some}
  arity $N$'', which is all the dichotomy proof uses. The precise arity matters
  only for \S4 of the source, which is cut.

\end{enumerate}

\subsection{Convention}

\begin{enumerate}\tightlist

\item[\textbf{C1}] $\mathbf Z_{p}\in\Vn$ requires $p\mid n-1$
  (Hazard~\ref{haz:Zp}).
\item[\textbf{C2}] $\cover\sigma$ is a tolerance, not a congruence
  (Hazard~\ref{haz:irreducible}). Typing it as a congruence collapses the
  linear/PC distinction silently.
\item[\textbf{C3}] The empty set is a subuniverse of every idempotent algebra,
  and is vacuously absorbing and vacuously central. Read syntactically, the
  clause defining $\subt{\TS}$ is therefore universally true. The global
  convention ``we do not allow empty subuniverses'' repairs it, but the clause
  does not inherit that convention on its own
  (Convention~\ref{conv:empty}).
\item[\textbf{C4}] Whether $\sigma=A^{2}$ counts as irreducible is genuinely
  ambiguous --- it turns on whether the empty intersection is allowed --- and
  $\cover\sigma$ does not exist in that case. \cite{ZhukSimplified} drops the
  word ``proper'' that \cite{ZhukJACM} has. Legislating that the empty family
  \emph{is} allowed makes irreducibility imply properness, which is the
  convenient choice and is what the Lean development adopts.
\item[\textbf{C5}] $\lll$ is data, not merely a proposition: several \S5 proofs
  speak of ``the congruences coming from $C\lll A$'' and of minimal chain
  length (Hazard~\ref{haz:lll-chain}).
\item[\textbf{C6}] ``Dimension'' of $\prod\mathbb Z_{q_{i}}$ with distinct
  primes is never defined, and there is no field over which it is a vector
  space. Adopt composition length and prove additivity
  (Hazard~\ref{haz:step2-linear-algebra}).
\item[\textbf{C7}] Two inequivalent definitions of ``linked instance'' coexist:
  global connectivity of the value graph, in Informal Claim (IC3); and the
  per-variable condition in \S3.1. The second does not imply the first. Use the
  formal one and carry non-fragmentation separately
  (Hazard~\ref{haz:linked}).
\item[\textbf{C8}] $B/\sigma$ is the image of $B$ in $A/\sigma$, not the
  quotient of $B$; and $(C_{1}\cap\dots\cap C_{t})/\delta=\bigcap C_{i}/\delta$
  is false in general and is proved on the fly inside
  Lemma~\ref{lem:propagation}(fm) (Hazard~\ref{haz:quotient-rel}).
\item[\textbf{C9}] Three sentences beginning ``Notice that'' are used as
  lemmas and never proved: that the relation $S$ of
  Definition~\ref{def:linear-congruence}(c) is a bridge with
  $\bcol S=\proj_{1,2}(S)=\proj_{3,4}(S)=\cover\sigma$; that $\sigma$ is
  irreducible/linear/PC exactly when $0_{\mathbf A/\sigma}$ is; and that the two
  formulations of $\TS$-freeness agree. Add
  Lemma~\ref{lem:rect-basics} to the same list.
\item[\textbf{C10}] One citation silently drops a case: the source's use of
  \cite[Thm.~6.15]{ZhukStrong} omits a third alternative --- a nontrivial
  \emph{projective} subuniverse --- present in the original. Its elimination is
  legitimate under the standing Taylor hypothesis, but the source does not say
  so, and a formalization must state it as a hidden import and check it against
  what the predecessor development actually proves.

\end{enumerate}

\begin{remark}[A retracted item, and what replaced it]\label{rem:d10-retracted}
An earlier draft of this section carried an eleventh blocking item: that the
recursion-depth bound $|A|+|\Gamma|$ of \cite[Lem.~5.2]{ZhukJACM} was ``not
established for the \textsc{CheckTuple} $\to$ \textsc{SolveNonlinked} path'',
and that ``without the repair the algorithm is not polynomial''. That was
wrong, and it is withdrawn.

Lemma~5.2 bounds the depth by arguing that every recursion site except
\textsc{CheckWeakerInstance} ``reduce[s] all domains of size greater than 1''.
At the \textsc{SolveNonlinked} site this needs the instance to be neither
linked nor fragmented --- otherwise the linked-component congruence can be
trivial on some domain, that domain does not shrink, and the bound degrades
from the constant $|A|$ to $O(|A|\cdot n)$, which is fatal, since the total work
is $\mathrm{poly}(N)^{\text{depth}}$. The source does state the hypothesis and
does discharge it. It is simply not in one place:
\begin{itemize}\tightlist
\item the precondition ``an instance that is not linked and not fragmented''
  is in the prose introducing the function, and \emph{not} in the
  \texttt{Input:} line of its pseudocode;
\item it is discharged at the call site by ``$\Theta_{0}$ is not fragmented
  because it is crucial'' --- itself a one-line argument the source does not
  give: a fragmented unsatisfiable instance has an unsatisfiable part, and no
  constraint of the \emph{other} part is then crucial;
\item the step from ``not fragmented'' to ``\emph{every} domain splits'' ---
  propagate the partition along a path, which exists by non-fragmentedness and
  is value-preserving by $1$-consistency --- appears once, inside the proof of a
  \emph{different} lemma, and is never cited here.
\end{itemize}
So the mathematics is present and the theorem is not in doubt. What is absent is
any single place where the six recursion sites of Lemma~5.2 are discharged, and
that is a blueprint's job rather than a defect of the source. The same pattern
holds for the $\Gamma$-order argument, whose case~(3) needs the projections to
be unchanged under weakening, which holds because the instance is
$1$-consistent when \textsc{CheckWeakerInstance} runs --- true, and unstated.

\emph{What does remain open}, and is now a sharp question rather than a flag:
\textsc{CheckTuple} calls \textsc{SolveNonlinked} not on $\Theta_{0}$ but on
$\Theta_{0}\wedge(x_{1}\in D_{1}')\wedge\dots\wedge(x_{n}\in D_{n}')$, and
restricting domains can only \emph{destroy} non-linkedness. Each $D_{i}'$ is a
union of $\sigma_{i}$-blocks, where $\sigma_{i}$ is the minimal linear
congruence; write $\tau$ for the linked-component congruence. If
$\tau\subseteq\sigma$ then the restriction deletes whole $\tau$-blocks and
non-linkedness survives. Whether $\tau\subseteq\sigma$ at this point is not
stated in the source, and I have not checked it. That is the next thing to
verify, and it is the kind of question a blueprint exists to surface.
\end{remark}

\begin{remark}[What this list is evidence of]\label{rem:defects-meaning}
Nine blocking items in a fifty-page paper is not a high rate; it is roughly what
one should expect, and it is the reason blueprints exist. Note the
distribution: one is a genuine misstatement (D1), one is an omission (D2), two
are gaps in the mathematics that need new arguments (D3, D4), three are
hypotheses that are too weak or missing (D5, D6, D8), and two are artefacts of
the interface with the older paper (D7, D9). None of them threatens the
theorem. All of them cost time.

Remark~\ref{rem:d10-retracted} is the counterweight, and the more instructive
entry. A tenth item was listed, survived a review round, and was published ---
and was wrong, because it was taken from a reader who worked from the pseudocode
while the hypothesis it needed was in the surrounding prose. A list like this
one is exactly as reliable as the reading behind each entry, and the failure
mode is not carelessness but reading the formal-looking part of a paper and not
the sentence before it.
\end{remark}

\clearpage
%!TeX root=csp.tex
\section{The complexity layer}\label{sec:complexity-layer}

Targets \textup{(T2)} and \textup{(H1)} are stated here and not proved. The
point of stating them is to be precise about what is \emph{not} being claimed.

\begin{definition}[What a cost model would have to supply]\label{def:cost-model}
To state \textup{(T2)} one needs: an encoding of instances of $\CSP(\Gamma)$ as
strings; a machine model with a step count; a predicate ``$f$ is computed in
time bounded by a polynomial in the input length''; and closure of that
predicate under composition. To state \textup{(H1)} one needs, in addition, a
polynomial-time many-one reduction relation $\le_{p}$ and the class
\textup{NP}.
\end{definition}

Mathlib supplies the machine (\texttt{Turing.TM2}), an encoding
(\texttt{Computability.Encoding}), and a time-bounded computation predicate
(\texttt{Turing.TM2ComputableInPolyTime}). It does not supply a complexity
\emph{class}, \textup{P}, \textup{NP}, \textup{NP}-completeness, $\le_{p}$, or
\textsc{Sat}; and the composition lemma
\texttt{TM2ComputableInPolyTime.comp} is an open \texttt{proof\_wanted} in
Mathlib itself. Moreover \texttt{TM2ComputableInPolyTime} applies to total
functions $\alpha\to\beta$, never to languages, so even the existing predicate
would need reshaping.

\begin{remark}\label{rem:complexity-scope}
Building this layer is a substantial project --- comparable in size to the
algebraic core, and with a large formal-verification literature of its own ---
and it contributes none of the mathematics of the dichotomy. Attempting it as
part of this project would be a scoping error. It is exactly the kind of
component that should be built once, in Mathlib, by people interested in
complexity theory, and then consumed.
\end{remark}

\section{Formalization}\label{sec:formalization}

\subsection{What is reused}\label{sec:reuse}

The predecessor development \texttt{ZhukLean} is a complete, \texttt{sorry}-free
Lean~4 formalization of Zhuk's centre theorem, built on
\texttt{Mathlib.ModelTheory}. It is consumed here as a \texttt{lake} dependency,
not copied. It supplies:

\begin{itemize}\tightlist
\item products of structures (\texttt{piStructure}, \texttt{prodStructure},
  coordinatewise realization, reindexing and evaluation homomorphisms) --- the
  one part of the universal-algebra vocabulary Mathlib lacks;
\item absorption in the tuple form (\texttt{Witnesses}, \texttt{Absorbs},
  \texttt{BinAbsorbs}), Taylor terms (\texttt{IsTaylorOn}), and central
  absorption (\texttt{CentrallyAbsorbs});
\item the relational description of absorption by essential relations
  (Barto--Kazda), which is what converts absorption at some arity into
  absorption at arity~$3$;
\item Zhuk's centre theorem itself (\texttt{zhuk\_center}) and the ternary
  collapse (\texttt{exists\_ternary\_witnesses}) --- which is precisely
  Lemma~\ref{lem:central-ternary}, cited without proof in
  \cite{ZhukSimplified}.
\end{itemize}

Three of \cite{ZhukSimplified}'s sixteen black-box imports are therefore already
discharged.

\subsection{What is built here}\label{sec:built}

\begin{center}
\begin{tabular}{@{}llr@{}}
\hline
Module & Contents & Status \\
\hline
\texttt{CSP/Product} & products of structures & complete \\
\texttt{CSP/Congruence} & \texttt{Congruence L M}, order, \texttt{sInf}, quotient & complete \\
\texttt{CSP/Irreducible} & irreducibility, existence and uniqueness of $\cover\sigma$ & complete \\
\texttt{CSP/Bridge} & bridges, collapse, composition, linear/PC & 2 imports open \\
\texttt{CSP/Types} & the six types, multi-types, $\lll$ & complete \\
\texttt{CSP/Instance} & instances, reductions, consistency, cruciality & complete \\
\texttt{CSP/Main} & the main statements & targets \\
\hline
\end{tabular}
\end{center}

Everything below \texttt{CSP/Main} is \texttt{sorry}-free. The two open items in
\texttt{CSP/Bridge} are Lemma~\ref{lem:bridge-comp} and its collapse identity,
both imported from \cite{ZhukJACM}.

\subsection{Design decisions}\label{sec:design}

\begin{description}
\item[Build on \texttt{Mathlib.ModelTheory}, keep the language generic.]
  A bespoke ``finite algebra with one WNU operation'' type would bundle the
  standing hypotheses as fields and avoid threading them, which is a real cost
  in the predecessor development. But it would discard
  \texttt{Substructures.lean} --- 985 lines of exactly the $\Sg$ API needed,
  including \texttt{mem\_closure\_iff\_exists\_term} with the variable type
  equal to the generating set, an interface no hand-rolled clone matches --- and
  it would discard the predecessor's 1600 lines wholesale. Congruences and
  quotients are also \emph{cheaper} in \texttt{ModelTheory}, not dearer.
  Instantiating to $\Vn$ is a thin layer on top.

\item[$\lll$ is \texttt{Relation.ReflTransGen}.]
  Its induction principle is the one the proofs use.

\item[Dotted relations are defined as disjunctions.]
  \texttt{DotLll C B := C = }$\varnothing$\texttt{ }$\vee$\texttt{ Lll C B}, so the case split is forced at every
  use site rather than resolved by the reader.

\item[The bridge criterion is the definition of linearity.]
  Zhuk defines a linear congruence by a per-block isomorphism to
  $\mathbb Z_{p}^{m}$ and derives the bridge criterion. We reverse this: the
  criterion is first-order over finite data and is what every consumer uses; the
  block description becomes a structure theorem. See
  Hazard~\ref{haz:linear-def}.

\item[$\cover\sigma$ is data attached to irreducibility.]
  \texttt{Irreducible.exists\_isCover} produces it; there is no standalone
  function, because there is no cover without irreducibility. A pleasant
  by-product: with the definition taken literally, irreducibility \emph{implies}
  $\sigma\neq A^{2}$, since the empty family intersects to the universe. The
  source never says this.

\item[Instances are lists; scopes are tuples; the variable type is a parameter.]
  See Hazards~\ref{haz:instance-list} and~\ref{haz:renaming}. The last of these
  is the decision that has not yet been cashed out, and it should be made before
  anything about coverings is written.

\item[Nothing is stubbed with \texttt{sorry} at the level of a definition.]
  \texttt{IsConnectedInstance} and \texttt{IsExpandedCovering} are
  \emph{absent} rather than defined as \texttt{sorry}, because a
  \texttt{def~:~Prop~:=~sorry} is a junk constant that downstream proofs can
  silently lean on. Absence forces the design decision; a stub hides it.
\end{description}

\subsection{Order of attack}\label{sec:order}

\begin{enumerate}\tightlist
\item Lemma~\ref{lem:ba-center-implies}. Small, self-contained, used everywhere,
  and adjacent to what is already formalized.
\item The mixed-prime linear algebra of Hazard~\ref{haz:step2-linear-algebra}.
  Independent of everything else; can be done in parallel.
\item Lemma~\ref{lem:bridge-comp}. Twelve source lines of relational algebra;
  closes the two open items in \texttt{CSP/Bridge}.
\item The coproduct design for coverings, then
  Proposition~\ref{prop:covering-basics}.
\item Theorem~\ref{thm:stable-intersection}. The one hard statement in the
  interface, and the sole source of bridges.
\item Lemma~\ref{lem:bridge-from-instance}. Eight hypotheses, ten pages.
\item Theorem~\ref{thm:main-inductive}. Last.
\end{enumerate}

\subsection{Scope}\label{sec:scope}

The calibration point is that the predecessor formalized about $3.5$ printed
pages in about $1670$ lines of Lean --- roughly $480$ lines per page. Excluding
\S\ref{sec:complexity-layer} and \cite{ZhukSimplified}'s \S4, which is an
independent second result and not needed for the dichotomy, the transitive
closure of what must be proved is on the order of $100$--$130$ printed pages,
so $35\,000$--$60\,000$ lines of Lean.

That is a multi-person-year project, and it should be described as one. What is
finished here is its design: the route, the statements, the vocabulary at
formalization granularity, the hazard list, the missing-layer inventory, and a
building development whose foundations are proved and whose summit is stated.


\clearpage
\begin{thebibliography}{99}

\bibitem{BJK} A.~Bulatov, P.~Jeavons, A.~Krokhin,
  \emph{Classifying the complexity of constraints using finite algebras},
  SIAM J. Comput.\ \textbf{34} (2005), 720--742.

\bibitem{BradyNotes} Z.~Brady, \emph{Notes on CSPs and polymorphisms},
  arXiv:2210.07383.

\bibitem{BartoKozik} L.~Barto, M.~Kozik, \emph{Absorbing subalgebras, cyclic
  terms, and the constraint satisfaction problem}, Log. Methods Comput. Sci.
  \textbf{8} (2012), 1:07.

\bibitem{DecidingAbsorption} L.~Barto, A.~Kazda, \emph{Deciding absorption},
  Internat. J. Algebra Comput. \textbf{26} (2016), 1033--1060.

\bibitem{Bulatov} A.~Bulatov, \emph{A dichotomy theorem for nonuniform CSPs},
  arXiv:1703.03021; FOCS 2017.

\bibitem{FederVardi} T.~Feder, M.~Y.~Vardi, \emph{The computational structure of
  monotone monadic SNP and constraint satisfaction}, SIAM J. Comput.\
  \textbf{28} (1999), 57--104.

\bibitem{HobbyMcKenzie} D.~Hobby, R.~McKenzie, \emph{The structure of finite
  algebras}, Contemp. Math.~76, Amer. Math. Soc., 1988.

\bibitem{IstingerKaiser} M.~Istinger, H.~K.~Kaiser, \emph{A characterization of
  polynomially complete algebras}, J. Algebra \textbf{56} (1979), 103--110.

\bibitem{MarotiMcKenzie} M.~Mar\'oti, R.~McKenzie, \emph{Existence theorems for
  weakly symmetric operations}, Algebra Universalis \textbf{59} (2008),
  463--489.

\bibitem{MinimalTaylor} L.~Barto, Z.~Brady, A.~Bulatov, M.~Kozik, D.~Zhuk,
  \emph{Minimal Taylor algebras as a common framework for the three algebraic
  approaches to the CSP}, LICS 2021.

\bibitem{RossSlides} R.~Willard, \emph{Bridges, similarity and the centralizer
  relation}, slides, AAA/NSAC workshop; see \cite{ZhukSimplified} for the use made
  of them here.

\bibitem{ZhukJACM} D.~Zhuk, \emph{A proof of the CSP dichotomy conjecture},
  J. ACM \textbf{67} (2020), no.~5, 1--78; arXiv:1704.01914.

\bibitem{ZhukSimplified} D.~Zhuk, \emph{A simplified proof of the CSP dichotomy
  conjecture and XY-symmetric operations}, arXiv:2404.01080v2.

\bibitem{ZhukStrong} D.~Zhuk, \emph{Strong subalgebras and the constraint
  satisfaction problem}, J. Mult.-Valued Logic Soft Comput.\ \textbf{36} (2021);
  arXiv:2005.00593.

\end{thebibliography}

\end{document}
